%% learnheuristic_pmsp.tex
%% Compilar: xelatex -> bibtex learnheuristic_pmsp -> xelatex -> xelatex
%% Migrado para a classe elsarticle (Elsevier) na Entrada 76 (08/09/2026)
%% -- exigencia do EJOR (ver ANALISE_E_PLANO_EJOR.md, Secao 4). elsarticle.cls
%% e elsarticle-harv.bst (citacao estilo Harvard/autor-ano, pedido do EJOR)
%% vem junto deste arquivo, na mesma pasta -- nao fazem parte de nenhuma
%% instalacao TeX Live padrao (pacote texlive-publishers, confirmado ausente
%% tanto no ambiente de Carlos quanto no de verificacao).
%% Opcoes da classe atualizadas na Entrada 85 (09/09/2026, Task #45) para
%% bater exatamente com o artigo irmao do LCPP (Amaro Junior et al., em
%% preparacao para o EJOR, mesmo grupo/coautoria de Carlos): [final,3p]
%% no lugar de [preprint,review] -- decisao explicita de Carlos, coluna
%% unica em largura "de producao" (mais estreita que o preprint largo),
%% sem numeracao de linha de revisor. [times] ativa txfonts (Type 1,
%% compativel com xelatex sem fontspec) em vez da Computer Modern padrao,
%% [11pt] no lugar de [12pt] -- ambos tambem para bater com o LCPP.
%% \onehalfspacing (pacote setspace) e' carregado explicitamente depois
%% da classe pra manter 1.5 espaco mesmo com [final] (que por si so'
%% forcaria espaco simples) -- mesmo padrao usado no LCPP, apropriado
%% enquanto o artigo ainda esta' em fase de rascunho/revisao interna.
%% [authoryear] ativa a citacao estilo Harvard via natbib (\citep/\citet),
%% carregado automaticamente pela propria classe -- nao precisa de
%% \usepackage{natbib} manual.
\documentclass[final,3p,times,11pt,authoryear]{elsarticle}

\usepackage[utf8]{inputenc}
\usepackage[T1]{fontenc}
\usepackage{amsmath,amssymb,amsfonts}
\usepackage{graphicx}
\usepackage{textcomp}
\usepackage{subcaption}
\usepackage{booktabs}
\usepackage{longtable}
\usepackage{multirow}
\usepackage{url}
\usepackage{float}
\usepackage{algorithm}
\usepackage{algorithmic}
\floatstyle{ruled}
\restylefloat{algorithm}
\usepackage{setspace}
\onehalfspacing
\usepackage{enumitem}
\setlist{nosep,topsep=2pt,itemsep=1pt,parsep=0pt,leftmargin=*}
\setlength{\textfloatsep}{8pt plus 2pt minus 2pt}
\setlength{\floatsep}{6pt plus 2pt minus 2pt}
\setlength{\intextsep}{6pt plus 2pt minus 2pt}
\setlength{\abovecaptionskip}{4pt}
\setlength{\belowcaptionskip}{2pt}
\usepackage{xcolor}
\usepackage{bm}
\usepackage{xr}
\externaldocument[S-]{supplementary_material}

\sloppy

\journal{European Journal of Operational Research}

\begin{document}

\begin{frontmatter}

\title{Mechanism or policy? Attributing the gains of Q-Learning-guided
       Iterated Greedy for parallel machine scheduling}
%% Titulo trocado na Entrada 95 (04/10/2026), decisao de Carlos (Entrada
%% 94) -- RASCUNHO, texto final a revisar por Carlos. O titulo anterior
%% ("Learnheuristics with Adaptive Operator and Strategy Selection for the
%% Parallel Machine Scheduling Problem: Hybridizing Iterated Local Search
%% and Iterated Greedy with Online Q-Learning") superalegava o hibrido
%% (preliminar) e a enfase em aprendizado.
%% TODO(Task #44, Fase A -- rascunho para revisao de Carlos): titulo
%% ampliado em relacao ao original ("A Learnheuristic ... Hybridizing
%% Iterated Local Search with Online Q-Learning") para refletir que
%% ig_ql (Iterated Greedy + Q-Learning) e' agora contribuicao central,
%% nao so' uma variante secundaria -- decisao de Carlos, ver
%% diario_pesquisa.md, Entrada 79.

\author{Carlos Vinicio Dantas Araújo\corref{cor1}}
\ead{carlos.vinicio@aluno.uece.br}
\address{Graduate Program in Computer Science, State University of Ceará
         (UECE), Fortaleza, CE, Brazil}
\cortext[cor1]{Corresponding author.}
%% TODO(Carlos): coautoria (ex.: orientador) e' decisao sua, conforme
%% combinado -- adicione aqui com \author{}/\address{} extra se aplicavel,
%% antes da submissao.

%% Highlights adicionadas na Entrada 85 (09/09/2026, Task #45), seguindo
%% o padrao do artigo irmao do LCPP (formato de bullets curtos, sem
%% citacoes/formulas, <=85 caracteres cada, exigencia do EJOR) -- baseadas
%% nos numeros ja verificados da Secao 5/7, nao copiadas do LCPP.
\begin{highlights}
\item Search structure, not Q-Learning, drives the gain; the policy adds 0.01-0.05 pp.
\item Headroom is necessary, not sufficient for learning; offline tuning beats online learning.
\item Learning-free recipe: 75\% provably optimal; a local search is worth only its CPU.
\item Same decomposition on unrelated machines; the ILS policy is harmful in both problems.
\item Granularity decides IG vs ILS: six distributions, equal CPU, strong starts.
\end{highlights}

\begin{abstract}
Where do the gains of Q-Learning-guided metaheuristics for parallel machine
scheduling come from: the search structure or the learned policy? Two
learnheuristics embed tabular Q-Learning as adaptive operator selection for the
identical-machine problem (P$||$C$_{\max}$): \texttt{ig\_ql} chooses among
destruction strategies of an Iterated Greedy (IG), \texttt{ils\_ql} among
neighbourhood operators of an Iterated Local Search (ILS). On twelve synthetic
classes and 6\,780 published instances, \texttt{ig\_ql} is never significantly
beaten by LPT or seven ILS variants and beats simulated-annealing and Bee Colony
baselines, but MULTIFIT and the differencing method (LDM) beat it for large $n$.
Random-selection controls attribute most of the gain to the search structure.
The learned policy adds 0.01--0.05 percentage points of GAP; this replicates on
2\,690 new instances, does not depend on the state, and is limited by the
headroom of the actions: for the destruction size, where there is headroom,
neither more budget nor pre-training lets Q-Learning match choosing the
parameter offline by validation. That learning-free recipe, started from the best
of three constructive heuristics, is provably optimal on $75\%$ of the synthetic
instances (MULTIFIT: $49\%$) according to a certified arc-flow bound; a local
search added to it, or to the plain IG, helps at an equal number of iterations
but not at equal CPU time. On
unrelated machines the decomposition recurs (search structure $4.2\%$, policy
about $0.1\%$, ILS policy harmful) and the policy gain follows the headroom of the destruction strategies
across cells, also on the standard benchmark of 700 test instances (decomposition replicated: search structure $2.1\%$, policy $0.03$--$0.06\%$,
not confirmed), but the law fails prospectively for the destruction size, where an agent trained online or offline is worse than choosing the size by validation: headroom is necessary, not sufficient; on a third problem, the permutation flowshop, the active ingredient is the local search (the IG equals the ILS at equal CPU time) and the agent over the destruction size is again worse than random; exact methods remain ahead, but an Iterated Greedy started from the linear relaxation cuts the
GAP to the best known solutions from $2.2\%$ to $1.0\%$ ($0.3\%$ for $1000$ jobs), and $0.25\%$ when followed by a restricted MILP (CPLEX in two hours: $0.20\%$), and a general-purpose MILP with a 60-second limit beats the earlier heuristics up to $1000\times50$ and is matched by this one
(equal or better in $62\%$ of the instances; not significant overall). ILS is stronger where processing times are nearly all distinct, on
six distributions, at equal CPU time and with strong starts. Hypotheses were
specified before the tests, several on held-out data.
\end{abstract}

\begin{keyword}
Parallel machine scheduling \sep Iterated Greedy \sep Reinforcement
learning \sep Adaptive operator selection \sep Metaheuristics
%% TODO(Carlos): o EJOR exige 1 palavra-chave obrigatoria vinda da lista
%% oficial do sistema de submissao (ANALISE_E_PLANO_EJOR.md, Secao 4) --
%% as 5 acima sao um rascunho de trabalho; confirme/ajuste no portal
%% antes de submeter.
\end{keyword}

\end{frontmatter}

%% TODO(Task #44 / pre-submissao): secoes obrigatorias do EJOR ainda
%% faltando neste arquivo -- Highlights, CRediT authorship statement,
%% Declaration of Interest, Funding statement, Data Availability
%% Statement, Acknowledgments (ver ANALISE_E_PLANO_EJOR.md, Secoes 4-5).
%% Carlos optou por nao priorizar esta rodada (nao selecionou "Adiantar
%% textos administrativos curtos" na Entrada 74).

% =====================================================================
\section{Introduction}
\label{sec:intro}
%% TODO(Task #44, Fase A -- rascunho para revisao de Carlos): secao
%% inteira traduzida para ingles e reescrita nesta fase (Entrada 79).
%% Paragrafos de motivacao geral (PMSP/ILS/AOS/RL) sao traducao fiel do
%% original com uma introducao do framework IG logo no 2o paragrafo
%% (antes, IG so aparecia como afterthought bem mais adiante) e
%% citep{RuizStutzle2007} adicionada (referencia fundacional do IG,
%% existia no refs.bib desde a Entrada 46 mas nunca era citada). O
%% paragrafo de diferenciacao ig_ql-vs-QIG e' so' traducao do texto
%% corrigido na Entrada 78 (mesmo conteudo, sem novas alegacoes). A
%% lista de contribuicoes foi expandida de 4 para 5 itens cobrindo
%% AMBAS as learnheuristicas -- decisao de Carlos (Entrada 79) de
%% incluir tudo no corpo principal, mesmo que isso estique a paginacao
%% alem do teto do EJOR (ele ajusta manualmente depois).

Scheduling problems are among the most studied in combinatorial
optimisation, with applications in manufacturing, cloud computing,
task allocation in distributed systems, and logistics
\citep{Pinedo2016}.
The Parallel Machine Scheduling Problem (PMSP) consists of assigning
$n$ jobs to $m$ identical parallel machines so as to minimise the
makespan $C_{\max}$, defined as the completion time of the last job.
Even for $m=2$ the problem is NP-hard \citep{GareyJohnson1979}.
Exact methods nevertheless solve a large share of published benchmark
instances in well under a second \citep{OstojicEtAl2025}; they remain
expensive on the hardest instances, which is where heuristics are
used.

Metaheuristics are widely adopted for the PMSP.
Iterated Local Search (ILS) \citep{LourencaoEtal2019} is one of the
most popular approaches, alternating between a local-improvement phase
and a perturbation step to escape local optima.
Iterated Greedy (IG) \citep{RuizStutzle2007} is a closely related
framework that instead alternates cycles of destruction (removing a
subset of jobs from the current sequence) and reconstruction (greedily
reinserting them); the destruction-reconstruction cycle itself provides
both search and diversification, with no separate perturbation phase.
In both frameworks, when multiple neighbourhood operators or
destruction strategies are available, the question of which one to
apply at each iteration arises: Adaptive Operator Selection (AOS)
addresses this demand through mechanisms that estimate or learn the
relative effectiveness of each choice during the search
\citep{FialhoEtal2010,KarimiMamaghanEtal2022}.

Classical AOS strategies include uniform random selection
\citep{PisingerRopke2007}, bandit algorithms such as UCB (\textit{Upper
Confidence Bound}) \citep{AuerEtal2002}, and data-driven predictors
such as Random Forests (RF) \citep{MaturanaSaubion2008}.
A limitation of these approaches is that they generally do not exploit
the \emph{context} of the current search: solution quality, degree of
load imbalance, budget progress, and stagnation level.

Reinforcement Learning (RL), and Q-Learning in particular
\citep{WatkinsDayan1992,SuttonBarto2018}, offers in principle a
context-sensitive AOS mechanism: the agent maintains action-state value estimates and
updates its selection policy based on the return obtained at each
iteration.
Hybrid RL-metaheuristic approaches have shown promise in job-shop
scheduling \citep{ParkEtal2021}. In the scheduling literature, Q-Learning as an
AOS mechanism has been explored mainly for permutation flowshop,
integrated with iterated greedy \citep{KarimiMamaghanEtal2023} and with
dynamic operator-portfolio management schemes
\citep{KarimiMamaghanEtal2024}; for the PMSP specifically, RL has
already been combined with swarm-intelligence metaheuristics in
variants with setup and release times, but for dynamic parameter
tuning rather than operator selection \citep{Ozsoydan2025}. More
recently, Q-Learning has been integrated with
Iterated Greedy for the PMSP with sequence-dependent setup times,
targeting weighted earliness/tardiness -- not makespan -- via greedy
destruction-reconstruction rather than neighbourhood operators
\citep{MissaouiOSullivan2026}. The use of tabular Q-Learning as AOS
over neighbourhood operators within an ILS, specifically for the pure
PMSP (P||Cmax, identical machines, makespan minimisation), remains, to
date, unexplored.
The second \emph{learnheuristic} considered here, \texttt{ig\_ql}, employs
tabular Q-Learning analogously within an \emph{Iterated Greedy} scheme,
in the same general spirit as \citet{MissaouiOSullivan2026}, but it
differs in the problem variant and in every component of the learning
mechanism. They address sequence-dependent setup times, due-date windows
and weighted earliness/tardiness; \texttt{ig\_ql} addresses the pure
PMSP (P||Cmax, identical machines, no setups). There, Q-Learning selects
the destruction \emph{size} $d \in \{1,\ldots,5\}$ for a single fixed
destruction-reconstruction procedure, over a state space of the
solutions themselves, with a reward of
$(f_{\mathrm{before}} - f_{\mathrm{after}})/10$ on improvement and
simulated-annealing-style acceptance. Here, Q-Learning selects
\emph{among three qualitatively distinct destruction strategies}
(random removal, removal of the longest jobs, machine-concentrated
removal), with a fixed destruction size, over the same compact 81-state
representation as \texttt{ils\_ql} (Section~\ref{sec:state}), with a
raw-improvement reward truncated at zero and purely greedy acceptance.
This combination remains, to date, unexplored in both problem variant
and adaptation mechanism.
% RESOLVIDO (08/09/2026, Entrada 78): o TODO anterior (Entrada 72)
% ficava em aberto por falta de acesso ao texto completo do QIG.
% Carlos obteve o artigo diretamente com o Dr. Missaoui (email) e a
% comparacao acima ja incorpora o nivel de mecanismo, nao so
% objetivo/variante -- ver diario_pesquisa.md, Entrada 78. Traduzido
% para ingles na Fase A da Task #44 (Entrada 79) -- mesmo conteudo, so
% traducao; reposicionamento para a nova subsecao de Trabalhos
% Relacionados sobre IG/adaptacao de parametros fica pra Fase B.

This article makes five contributions.
(i) A controlled decomposition of the gains of a learnheuristic into search
structure and learned policy: a random-selection control at matched budget,
applied to two unrelated search mechanisms (ILS and IG) that share the state
representation and the learning agent, and repeated on a second problem
(unrelated machines, including the standard benchmark) and a third (permutation flowshop), where the active ingredient turns out to be a different one. The search structure accounts for most of the advantage;
the learned selection accounts for a small remainder (about $0.01$--$0.05$
percentage points for identical machines, about $0.1\%$ for unrelated ones)
(Sections~\ref{sec:f3f4-ablation} and~\ref{sec:unrelated}).
(ii) An explanation of why the remainder is so small, tested rather than
asserted: the destruction strategies leave almost no headroom; where the action
space has headroom (the destruction size), neither a larger budget, a dense
reward, a more informative state nor cross-instance pre-training lets
Q-Learning match choosing the parameter offline by validation; and across the
cells of the second problem the policy gain follows the headroom of the destruction strategies, although a prospective test with the
destruction size as the action, which has headroom, failed: the headroom is necessary, not sufficient
(Sections~\ref{sec:f3f4-ablation} and~\ref{sec:unrelated}).
(iii) The IG-based learnheuristic \texttt{ig\_ql} against stronger baselines
than usual (LPT, seven ILS variants, a simulated-annealing method following
\citet{LiPerkins2024}, a Bee Colony Optimisation algorithm
\citep{DavidovicEtal2012}, MULTIFIT and the differencing method LDM) on twelve
synthetic classes and 6\,780 published instances, with a certified lower bound
from an arc-flow relaxation. It is never significantly beaten by LPT or ILS, but
MULTIFIT and LDM beat it for large $n$ (Sections~\ref{sec:f2-sota} and~\ref{sec:exact-ref}); a learning-free
recipe, a destruction size chosen on separate instances and a start from the best
constructive solution, is the strongest method of the study ($75\%$ provably
optimal instances against $49\%$ for MULTIFIT; Section~\ref{sec:recipe}); a local
search added to it pays at an equal number of iterations but not at equal CPU
time, when the recipe with a larger iteration budget reaches $89\%$ ($97.5\%$ at $64{,}000$ iterations); for unrelated machines, an Iterated Greedy started from the
linear relaxation halves the GAP to the best known solutions of the standard benchmark and, followed by a restricted MILP, reaches $0.25\%$ (Section~\ref{sec:unrelated}).
(iv) An instance-level finding that the granularity of the processing times (how
many distinct values they take) decides in which direction IG or ILS is stronger,
while $n/m$ shifts the crossover (Section~\ref{sec:nm-structural}); a one-feature
rule that chooses between them was tested out of sample (Section~\ref{sec:dispatch}),
and both the finding and the rule hold across six processing-time distributions,
at equal CPU time and with strong starts in both mechanisms (Supplementary
Section~\ref{S-sec:robust}).
(v) Standards of evidence for studies of adaptive selection: hypotheses and
decision rules written down before every run (with hashes of the code), tests
specified on held-out and on fresh instances, equivalence margins, seed
averaging, exact references, and the reporting of every non-confirmation
(Supplementary Section~\ref{S-sec:prereg-record}). A companion study of the same
agent inside ILS (\texttt{ils\_ql}), whose null result and state-coverage artefact
motivated the decomposition, is summarised in Section~\ref{sec:ils-summary} and
reported in full in the supplement.
%% TODO(Task #44, Fase A): lista expandida de 4 para 5 contribuicoes
%% para cobrir ig_ql (novo item (i), headline) sem descartar os
%% achados do ils_ql (agora (iii)-(v)) -- decisao de Carlos (Entrada
%% 79) de incluir tudo no corpo principal. Nomes exatos dos 3
%% baselines SOTA em (i) a conferir contra a Secao~\ref{sec:experiments}
%% quando essa secao for reescrita (Fase D).

The remainder of this article is organised as follows:
Section~\ref{sec:related} reviews related work; Section~\ref{sec:problem}
defines the PMSP; Section~\ref{sec:method} describes the proposed
learnheuristics; Section~\ref{sec:experiments} presents the
computational study; Section~\ref{sec:discussion} discusses broader
implications; and Section~\ref{sec:conclusions} concludes the paper.
%% TODO(Task #44, Fase A): a versao original em portugues (linhas
%% 221-225 do .tex pre-Fase-A) tambem omitia a Secao de Discussao do
%% roteiro -- lacuna preexistente, nao introduzida por esta traducao.
%% Corrigida aqui (dentro do escopo da Introducao) ao inves de repetir
%% a omissao; ver Entrada 79 do diario.


% =====================================================================
\section{Related Work}
\label{sec:related}
%% TODO(Task #44, Fase B -- rascunho para revisao de Carlos, Entrada 80):
%% secao inteira traduzida para ingles nesta fase. As 3 subsecoes 2.1-2.3
%% sao traducao fiel do original (com correcao de um bug de citacao
%% descrito abaixo). A subsecao 2.4 e' NOVA, incorporando 4 referencias
%% que ja existiam em refs.bib desde a Entrada 74/78 mas nunca tinham
%% sido citadas no corpo do texto (RuizStutzle2007, MissaouiRuiz2022,
%% MissaouiOSullivanStrategicRecon2026, MissaouiOSullivan2026).
%% Bug corrigido nesta fase (achado ao traduzir, nao introduzido agora):
%% a migracao global \cite->\citep da Entrada 76 converteu tambem casos
%% em que o nome do autor ja estava escrito por extenso na prosa (ex.:
%% "Fialho et al. \citep{FialhoEtal2010}"), o que sob estilo Harvard
%% duplica o nome ("Fialho et al. (Fialho et al., 2010)"). Corrigido
%% para \citet{} nesses 8 casos (renderiza "Fialho et al. (2010)"), a
%% forma padrao quando o nome do autor e' o sujeito gramatical da frase.

This section reviews the literature underlying the two proposed
learnheuristics. It first covers adaptive operator and strategy
selection (AOS) mechanisms, reinforcement learning applied to
combinatorial optimisation more broadly, and the parallel machine
scheduling literature relevant to the problem studied here, before
turning to iterated greedy and parameter/strategy adaptation
(Section~\ref{sec:ig-relwork}), the framework underlying
\texttt{ig\_ql}.

\subsection{Adaptive Operator Selection}

AOS has emerged as a well-established research area within
evolutionary algorithms and hyper-heuristics \citep{BurkeEtal2013}.
\citet{FialhoEtal2010} conducted a comprehensive survey of bandit-based
AOS mechanisms applied to genetic algorithms, finding that simple
exploration-exploitation strategies frequently outperform more complex
methods in small operator spaces.
\citet{MaturanaSaubion2008} proposed a Random-Forest-based selector
that learns to map search states to operators via supervised
regression, obtaining gains on constraint satisfaction problems.
\citet{PisingerRopke2007} introduced ALNS (\textit{Adaptive Large
Neighbourhood Search}), in which operator weights are updated based on
recent performance, obtaining benchmark results on routing problems.

Despite these advances, most AOS methods do not explicitly model the
state of the search: the selected operator does not depend on features
such as proximity to a local optimum, stagnation, or budget progress.
This limitation motivated the investigation of RL-based mechanisms, in
which the agent observes a state and learns a state-dependent policy.

\subsection{Reinforcement Learning in Combinatorial Optimisation}

The integration of RL with metaheuristics has grown significantly over
the last decade \citep{BengioBengio2021}.
\citet{CalvetEtal2017} coined the term \emph{learnheuristic} for
metaheuristics coupled to a learning component, originally to handle
inputs that change as a solution is built; the term is used here in
the broader sense of a search whose decisions are guided by an online
learner.
\citet{ParkEtal2021} applied graph neural networks with RL to job-shop
scheduling, obtaining policies that generalise to unseen instances.
\citet{CappartEtal2021} reviewed the intersection between constraint
programming and machine learning, pointing to RL as one of the most
promising approaches for guiding search.

In the specific context of AOS, the use of tabular Q-Learning as an
operator selector within ILS remains little explored for classical,
discrete production scheduling: \citet{JuanEtal2014}
applied ILS to the flow shop scheduling problem with a deliberately
fixed, non-adaptive parametrisation, and Q-Learning-based operator
selection for scheduling has instead been explored within the closely
related Iterated Greedy framework \citep{KarimiMamaghanEtal2023}, not
ILS. \citet{VazquezRodriguezEtal2011} pursued a related but distinct
direction, using genetic programming to automatically discover novel
constructive-heuristic variants of the NEH procedure for flow shop
scheduling -- an automated-design approach rather than an adaptive
operator-selection mechanism within an iterative search. Outside this
setting, \citet{AlicastroEtal2021} embedded a Q-Learning-based Variable
Neighbourhood Search within ILS for makespan minimisation in
additive-manufacturing machine scheduling, a domain with different
constraints and decision variables from the discrete-parts PMSP
considered here.
This gap is addressed here by proposing and empirically evaluating a
tabular Q-Learning agent as an AOS mechanism within ILS for the PMSP.

\subsection{Parallel Machine Scheduling}

The PMSP is among the most extensively studied scheduling problems
\citep{Pinedo2016}.
The LPT (\textit{Longest Processing Time}) rule of \citet{Graham1969}
is the best-known constructive heuristic, with an approximation
guarantee of $(4/3 - 1/(3m))$.
\textsc{MultiFit} \citep{CoffmanGareyJohnson1978} binary-searches the
capacity and packs with First-Fit Decreasing, and the Largest
Differencing Method (LDM), the multiway form of the Karmarkar--Karp
differencing heuristic \citep{KarmarkarKarp1983,MichielsEtal2007},
repeatedly merges partial solutions so as to cancel the largest
differences; both are fast and deterministic. Local search with
multi-exchange neighbourhoods is a further classical route
\citep{FrangioniEtal2004}.
For unrelated machines, Iterated Greedy with local search is the established
heuristic \citep{FanjulPeyroRuiz2010}.
The problem is NP-hard \citep{GareyJohnson1979}, yet exact methods
have become strong: bin-packing-based and arc-flow formulations
\citep{DellAmicoEtal2008,MradSouayah2018,GharbiBamatraf2022}, engineered
branch-and-bound \citep{AkramEtal2024} and an iterative arc-flow method
solve most published benchmark instances in well under a second,
although the hardest instances remain costly \citep{OstojicEtAl2025}.
The linear relaxation of the arc-flow model is also a cheap, certified
lower bound, which this article uses (Section~\ref{sec:exact-ref}) to
tell near-optimal solutions apart from provably optimal ones.
Metaheuristics remain
relevant for those hard instances and for variants where exact methods
do not reach; the specific combination of AOS
mechanisms -- whether classical or RL-based -- with the two search
frameworks used in this article (ILS and IG) is examined next.

\subsection{Iterated Greedy and Parameter/Strategy Adaptation}
\label{sec:ig-relwork}

\citet{RuizStutzle2007} introduced Iterated Greedy for the permutation flowshop, showing that a simple
destruction--reconstruction cycle is competitive with far more complex metaheuristics; this motivates IG as the second
search framework here (Section~\ref{sec:method}). Its destruction size is, however, a sensitive parameter that usually
needs instance-specific tuning. \citet{MissaouiRuiz2022} adapted it automatically for a hybrid flowshop with setup times and
due-date windows, and introduced a ``strategic'' reconstruction (each removed job goes to its first feasible position) later
reused, with Q-Learning, by \citet{MissaouiOSullivan2026}. For identical parallel machines,
\citet{MissaouiOSullivanStrategicRecon2026} paired that reconstruction with a non-learning adaptive ``leap size'' for
weighted tardiness with setup times, and \citet{MissaouiOSullivan2026} replaced the deterministic rule by tabular
Q-Learning that selects the destruction size, for setups and weighted earliness/tardiness; this is the closest work to
\texttt{ig\_ql}, and Section~\ref{sec:intro} details the differences at the level of the learning mechanism, not
only of the problem variant. None of this line addresses the pure PMSP studied here (no setups, no due dates, makespan
objective): each adds at least one further feature or, for \citet{RuizStutzle2007}, is a different problem class.


% =====================================================================
\section{Problem Definition}
\label{sec:problem}

Let $\mathcal{J}=\{1,\ldots,n\}$ be the set of jobs and
$\mathcal{M}=\{1,\ldots,m\}$ the set of identical parallel machines.
Each job $j \in \mathcal{J}$ has a processing time $p_j > 0$.
A feasible solution assigns each job to exactly one machine; the
completion time of machine $i$ is $C_i = \sum_{j \in S_i} p_j$, where
$S_i \subseteq \mathcal{J}$ is the set of jobs assigned to machine $i$.
The objective is to minimise $C_{\max} = \max_{i \in \mathcal{M}} C_i$.

A classical lower bound \citep{Graham1969} is given by:
\begin{equation}
  \mathrm{LB} = \max\!\left(\max_{j \in \mathcal{J}} p_j,\;
  \left\lceil \frac{\sum_{j} p_j}{m} \right\rceil\right).
  \label{eq:lb}
\end{equation}

The optimality \emph{gap} is defined as
$\mathrm{GAP}(\%) = 100 \cdot (C_{\max} - \mathrm{LB})/\mathrm{LB}$;
$\mathrm{GAP} = 0$ indicates that the lower bound was attained.

\paragraph{Instance difficulty}
The LPT rule, which assigns jobs in non-increasing order of $p_j$,
tends to leave a significant GAP when the jobs-per-machine ratio $n/m$
is low, since there are few opportunities for load balancing.
The six original classes of this work span $n/m \in [3.3, 5.0]$, a
regime in which LPT produces mean gaps of 1.9--4.3\%, ensuring
meaningful differentiation between methods without trivial instances;
six further classes extend the range to $n/m \in [3.0, 15.0]$, where
the mean LPT gap falls as low as $0.16\%$ (Section~\ref{sec:setup}).

\paragraph{Relationship to PMSP structure}
The makespan is determined directly by machine loads.
Operators that transfer jobs between distinct machines reduce
$C_{\max}$ by redistributing load; operators that only reorder jobs
within the same machine leave $C_i$ unchanged and, therefore, cannot
improve the makespan.
This structural property is confirmed empirically by the analysis of
the learned policy in Section~\ref{S-sec:policy}.

\paragraph{The unrelated-machine variant}
Section~\ref{sec:unrelated} repeats the study for $\mathrm{R}||C_{\max}$, in which
the processing time of job $j$ on machine $i$ is $p_{ij}$ and the load of machine $i$
is $C_i=\sum_{j\in S_i}p_{ij}$. The classical lower bound becomes
$\max\{\max_j\min_i p_{ij},\lceil\sum_j\min_i p_{ij}/m\rceil\}$, and the linear
relaxation of the assignment model gives a stronger certified bound
\citep{LenstraEtal1990}.


% =====================================================================
\section{Proposed Learnheuristics}
\label{sec:method}
%% TODO(Task #44, Fase C -- rascunho para revisao de Carlos, Entrada 81):
%% Secao 3 e as subsecoes 4.1-4.4 traduzidas fielmente (com 1 correcao de
%% bug, ver abaixo). Subsecoes 4.5-4.6 sao NOVAS, descrevendo o mecanismo
%% do ig_ql (nunca antes descrito no corpo do artigo -- confirmado por
%% leitura completa da Secao 4 anterior a esta fase) com base em leitura
%% direta de pmsp/ig.py, pmsp/qlearning.py, pmsp/multifit.py, pmsp/sa.py,
%% pmsp/bco.py e experiment/runner.py/config.py (nao de memoria).
%% BUG DE CONTEUDO encontrado e corrigido nesta fase (nao introduzido
%% agora -- ja estava assim no .tex em portugues, antes desta traducao):
%% a Secao 4.3 (Agente Q-Learning) descrevia o cronograma de epsilon
%% como "epsilon=0,2, reduzido em 10% a cada 100 iteracoes", mas o
%% codigo real (pmsp/qlearning.py, QLearningAgent, confirmado via
%% experiment/config.py: QL_EPSILON_START=1.0, QL_EPSILON_MIN=0.05,
%% QL_EPSILON_DECAY=0.995) usa decaimento multiplicativo POR PASSO
%% (nao por bloco de 100), partindo de 1.0 (exploracao pura), piso 0.05,
%% atingido em ~591 passos -- o proprio docstring da classe confirma
%% esse calculo. Descricao factual do metodo estava desatualizada em
%% relacao ao codigo; corrigida para refletir o codigo real, verificado
%% nesta fase (nao assumido).

This section describes the ILS framework and its Q-Learning-based AOS
mechanism (Section~\ref{sec:ils}), then the
Iterated Greedy framework and its own Q-Learning-based
strategy-selection mechanism (Sections~\ref{sec:ig}--\ref{sec:igvariants}),
which reuses the same state representation and Q-Learning update rule.

\subsection{ILS Framework}
\label{sec:ils}

ILS starts from an LPT solution and iterates for $T$ steps. At each step the AOS module selects one of four neighbourhood
operators (inter-swap, reinsert, intra-swap, block-move), applied under best-improvement local search; after 50 consecutive
iterations without improvement a random perturbation (3 random moves) restarts the search from the best incumbent. The
intra-swap operator can never improve the makespan and serves as a structural control. The algorithm, the operator table
and the ILS variants used in the experiments (a pre-trained table, extended-budget variants with a random-selection
isolation control, UCB1 and online Random Forest selectors) are described in Supplementary Section~\ref{S-sec:ils-framework}.

\subsection{State Representation}
\label{sec:state}

The search context is encoded into four scalar features, each mapped
to three ordered bins ($b \in \{0,1,2\}$), producing $3^4 = 81$
discrete states:

\begin{itemize}
  \item $f_1$ (gap): current GAP relative to the LB; thresholds
        $\{0.10, 0.35\}$.
  \item $f_2$ (load CV): coefficient of variation of machine loads;
        thresholds $\{0.15, 0.40\}$.
  \item $f_3$ (progress): fraction of budget consumed; thresholds
        $\{0.33, 0.67\}$.
  \item $f_4$ (stagnation): fraction of recent iterations without
        improvement; thresholds $\{0.25, 0.60\}$.
\end{itemize}

The state index is $s = 27b_1 + 9b_2 + 3b_3 + b_4 \in \{0,\ldots,80\}$.
With these thresholds, however, the two features that describe the
incumbent carry almost no information: in the IG runs examined (four
classes, eight instances each) $f_1$ and $f_2$ stay in bin 0 at
\emph{every} step, because the gap of a near-optimal schedule is far
below $0.10$ and the load CV far below $0.15$. The agent then
distinguishes only $5$--$9$ states, which encode the progress and the
stagnation of the run rather than the quality or balance of the
incumbent. Section~\ref{sec:f3f4-ablation} tests whether recalibrating
the state changes any conclusion.

The choice of these four features is motivated by evidence from the
AOS literature \citep{FialhoEtal2010}: features related to solution
quality ($f_1$), structural diversity ($f_2$), search phase ($f_3$),
and stagnation ($f_4$) are meant to capture the aspects of the search
state most relevant to operator selection (Section~\ref{sec:f3f4-ablation}
shows that $f_1$ and $f_2$ are uninformative with these thresholds).
This representation is reused unchanged by the \texttt{ig\_ql} agent
introduced in Section~\ref{sec:ig}.

\subsection{Q-Learning Agent}
\label{sec:qlm}

The agent maintains a table $Q \in \mathbb{R}^{81 \times 4}$,
initialised to zero, updated at each step via off-policy Q-Learning
\citep{WatkinsDayan1992}:
\begin{equation}
  Q(s,a) \leftarrow Q(s,a) + \alpha \bigl[
    r + \gamma \max_{a'} Q(s', a') - Q(s,a)
  \bigr],
  \label{eq:ql}
\end{equation}
where $\alpha = 0.1$ is the learning rate, $\gamma = 0.9$ is the
discount factor, and $r = C_{\max}^{\mathrm{before}} -
C_{\max}^{\mathrm{after}} \geq 0$ is the reward (absolute makespan
reduction).
Operator selection follows an $\varepsilon$-greedy policy, with
$\varepsilon$ initialised at $1.0$ (pure exploration) and decayed
multiplicatively by a factor of $0.995$ after every step, down to a
floor of $0.05$ (reached after approximately 591 steps), so that
exploration dominates early in the run and exploitation dominates
later.

The Q-table is shared throughout a single instance's run, so the agent
simultaneously learns to explore different regions of the search space
over the $T$ iterations.
At the end of the run, the greedy policy $\pi(s) = \arg\max_a Q(s,a)$
can be inspected for interpretability.
The same update rule (Equation~\ref{eq:ql}), with the same $\alpha$,
$\gamma$, and $\varepsilon$-greedy schedule, is reused unchanged by
\texttt{ig\_ql} (Section~\ref{sec:ig}); the two methods differ only in
what the action $a$ selects, in the reward truncation, and in the
acceptance criterion applied to the resulting move.

\subsection{Iterated Greedy Framework}
\label{sec:ig}

The second learnheuristic, \texttt{ig\_ql}, embeds the same tabular
Q-Learning mechanism within an Iterated Greedy (IG) scheme instead of
ILS. IG (Algorithm~\ref{alg:ig}) alternates cycles of destruction --
removing a subset of jobs from the current solution -- and
reconstruction -- greedily reinserting them -- with no separate
perturbation phase: the destruction-reconstruction cycle itself
provides both search and diversification.

\begin{table}[ht]
\centering
\caption{The three destruction strategies used in IG.}
\label{tab:destroy}
\begin{tabular}{@{}clp{7.5cm}@{}}
\toprule
\textbf{Strategy} & \textbf{Name} & \textbf{Description} \\
\midrule
1 & random               & Removes $d$ uniformly random jobs. \\
2 & most-expensive       & Removes the $d$ jobs with the largest processing times. \\
3 & machine-concentrated & Removes $d$ jobs, all from a single, randomly chosen machine. \\
\bottomrule
\end{tabular}
\end{table}

\begin{algorithm}[tb]
\caption{IG with AOS for PMSP}
\label{alg:ig}
\begin{algorithmic}[1]
\REQUIRE Jobs $\mathcal{J}$, machines $\mathcal{M}$, iteration budget $T$, AOS module, destruction size $d$
\STATE $x \leftarrow \mathrm{LPT}(\mathcal{J},\mathcal{M})$; $x^{*} \leftarrow x$
\FOR{$t = 1, \ldots, T$}
\STATE $s \leftarrow \mathrm{encode}(x, t, T)$
\STATE $a \leftarrow \mathrm{AOS}(s)$
\STATE $x' \leftarrow \mathrm{Reconstruct}(\mathrm{Destroy}(x,\, \mathrm{Strategy}_a,\, d))$
\STATE $r \leftarrow \max\bigl(C_{\max}(x) - C_{\max}(x'),\, 0\bigr)$
\STATE Update AOS with $(s, a, r)$
\IF{$C_{\max}(x') \leq C_{\max}(x)$}
\STATE $x \leftarrow x'$
\ENDIF
\IF{$C_{\max}(x) < C_{\max}(x^{*})$}
\STATE $x^{*} \leftarrow x$
\ENDIF
\ENDFOR
\RETURN~$x^{*}$
\end{algorithmic}
\end{algorithm}

\paragraph{Destruction}
At each iteration, the AOS module selects one of the three destruction
strategies in Table~\ref{tab:destroy}.
The number of removed jobs, $d$, is fixed and proportional to instance
size, $d = \max(2, \mathrm{round}(0.3n))$, rather than learned. The
$30\%$ default was fixed in advance and never tuned on the evaluation
instances; the sensitivity of the results to $d$, to the learning rate,
the discount factor and the exploration schedule is reported in
Section~\ref{sec:f3f4-ablation} and Supplementary
Section~\ref{S-sec:hyper}.
%% TODO(Task #44, Fase D): a analise de sensibilidade de tamanho de
%% destruicao (Entradas 20-22) pode ser adicionada a Secao 5
%% (Experimentos) com mais detalhe, se Carlos quiser -- por ora,
%% mencionada aqui so' como justificativa da escolha de design.

\paragraph{Reconstruction}
Removed jobs are reinserted deterministically in non-increasing order
of processing time, each onto the currently least-loaded machine (an
NEH/LPT-style greedy rule), using a running per-machine load tracker so
that a removed job's original load contribution does not bias the
choice of destination machine for jobs reinserted after it.

\paragraph{Acceptance}
The reconstructed solution is accepted only if it does not worsen the
incumbent's makespan ($C_{\max}(x') \leq C_{\max}(x)$) -- the same
purely greedy acceptance rule used by ILS, chosen deliberately so that
any difference between the two frameworks reflects the search
mechanism itself (destruction-reconstruction vs.\ neighbourhood-operator
local search), not a difference in acceptance criteria.

\paragraph{Shared components and reward}
\texttt{ig\_ql} reuses, unchanged, the same 81-state representation
(Section~\ref{sec:state}) and the same Q-Learning update rule, learning
rate, discount factor, and $\varepsilon$-greedy schedule
(Section~\ref{sec:qlm}) as \texttt{ils\_ql}.
The reward is $r = \max\bigl(C_{\max}(x) - C_{\max}(x'),\, 0\bigr)$:
unlike best-improvement local search, which by construction never
worsens the current solution, the destruction-reconstruction step can
produce a worse candidate $x'$ that is then rejected by the acceptance
rule; truncating the reward at zero avoids assigning a negative
learning signal for a candidate that acceptance already discards.
The only structural differences from \texttt{ils\_ql} are therefore the
action space (three destruction strategies rather than four
neighbourhood operators), this reward truncation, and the
destruction-reconstruction search step itself.
Section~\ref{sec:ig-relwork} and Section~\ref{sec:intro} discuss how
this specific combination differs from the closest related use of
Q-Learning within Iterated Greedy for parallel-machine
scheduling.

\subsection{IG Method Variants and Baselines}
\label{sec:igvariants}

\paragraph{Ablation ladder}
To isolate which components of \texttt{ig\_ql} drive its performance,
four reduced variants are evaluated alongside the full method.
\texttt{ig\_random} replaces the Q-Learning agent with uniform random
selection among the three destruction strategies, isolating the
contribution of the IG search mechanism itself from any adaptive
selection policy.
The remaining three variants form a $2{\times}2$ factorial design that
crosses state-conditioning with temporal-difference (TD) learning --
the two components combined by Q-Learning: \texttt{ig\_ucb} uses a
global UCB1 bandit \citep{AuerEtal2002} over the three strategies (no
state, no TD); \texttt{ig\_ucb\_ctx} uses the same UCB1 formula
computed separately within each of the 81 states, updated by a simple
running average with no discounting or value propagation between
states (state, no TD); \texttt{ig\_ql\_flat} reuses the same tabular
Q-Learning update as \texttt{ig\_ql} -- bootstrapping via
$\gamma\max_{a'}Q(a')$, $\varepsilon$-greedy exploration -- but with a
single aggregated state (TD, no state); and \texttt{ig\_ql} itself
combines both (state and TD).
Because \texttt{ig\_ql\_flat} reuses the same underlying Q-Learning
implementation internally, with the state space collapsed to one state
rather than reimplemented, the only difference between
\texttt{ig\_ql\_flat} and \texttt{ig\_ql} is the size of the state
space.

\paragraph{Published baselines}
Three additional methods from the literature are compared against,
reimplemented from their original publications, with any necessary
deliberate deviation documented in the corresponding source-code
module.
MULTIFIT \citep{CoffmanGareyJohnson1978} is a deterministic constructive
heuristic with no local search: it binary-searches over a candidate
machine-capacity bound and repeatedly applies First-Fit-Decreasing
bin-packing until the tightest feasible capacity is found.
\texttt{sa\_li\_perkins} reimplements the simulated annealing method of
\citet{LiPerkins2024}: a neighbourhood of random inter-machine swaps and
moves, with worsening moves accepted with probability $\exp(-a/T)$ (a
fixed-parameter form that, as in the original article, does not scale
with the magnitude of the worsening) and $T$ decayed geometrically;
under the original parameterisation ($T_0=10$, decay $0.9$, $a=1$), the
temperature saturates within roughly the first 100 iterations
regardless of the total budget, so the method spends most of a longer
run in an essentially greedy regime -- a property of the published
design, not of the reimplementation used here.
BCO \citep{DavidovicEtal2012} is a population-based constructive
metaheuristic in which a colony of bees builds solutions job by job via
weighted-random task and machine selection, then recruits within the
colony via a roulette-wheel loyalty mechanism across iterations; the
``Modified BCO'' variant (their Section~5.3) is implemented, reported in
the original article as its best-performing configuration.
The baselines differ in how they are budgeted.
\texttt{sa\_li\_perkins} starts from the same LPT solution as
\texttt{ig\_ql} and runs under the same iteration budget (up to its own
early-stopping rule). \texttt{bco} is a constructive method: it starts
from an empty solution and, instead of the iteration budget, receives
the same total number of constructive moves, i.e.\
$I = \max(1,\lfloor T/(5n)\rfloor)$ outer iterations against $I=100$ in
the original publication (Section~\ref{sec:f2-sota} relaxes this rule).
MULTIFIT is a one-shot constructive method with no budget.


% =====================================================================
\section{Computational Experiments}
\label{sec:experiments}

This section reports the computational study; Table~\ref{tab:evidence} maps the
claims of the paper to the evidence. After the setup (Section~\ref{sec:setup}), it
presents the performance of \texttt{ig\_ql} against the ILS family and LPT
(Section~\ref{sec:f1-headline}), against three published baselines, with a
sensitivity analysis of their budget rules (Section~\ref{sec:f2-sota}), and against
stronger baselines (the differencing method, a start from the best constructive
solution) with a certified lower bound and a quality--time comparison
(Section~\ref{sec:exact-ref}); a decomposition of its advantage into a search
component and a learning component, and tests of why the latter is small, including a
comparison of a local-search hybrid at equal iterations and at equal CPU time
(Section~\ref{sec:f3f4-ablation}); a learning-free recipe that follows from them, and a test of
whether a local search added to it pays at equal CPU time (Section~\ref{sec:recipe}); an external validation on published benchmark instances,
with two pre-specified confirmatory tests and exact reference solutions
(Section~\ref{sec:external}); the structural limits of the ILS mechanism and the role
of the granularity of the processing times (Section~\ref{sec:nm-structural}); an
out-of-sample test of a one-feature rule that chooses between IG and ILS per instance
(Section~\ref{sec:dispatch}); the same protocol on unrelated machines, at sizes up to
$1000\times50$ against a MILP and on the standard benchmark (Section~\ref{sec:unrelated}); and a summary of the companion \texttt{ils\_ql} study
(Section~\ref{sec:ils-summary}). Convergence, further robustness checks (seeds, other
distributions, equal CPU time, hyperparameters), the certified bound and the record of
every pre-specified analysis are in the supplement.

\begin{table}[t]
\centering
\caption{Map of the claims of the study, the evidence for each and where it is reported. All hypotheses
were specified before the runs; non-confirmations (such as the budget law) are reported in the sections indicated.}
\label{tab:evidence}
{\scriptsize\setlength{\tabcolsep}{3pt}
\begin{tabular}{@{}p{0.27\textwidth}p{0.45\textwidth}p{0.18\textwidth}@{}}
\toprule
Claim & Evidence (size of the effect) & Sections \\
\midrule
1. The search structure, not the learned policy, drives the gain &
$\mathrm{P}||C_{\max}$: IG beats ILS by $0.3$--$0.7\%$ against a policy effect of $0.01$--$0.05\%$; $\mathrm{R}||C_{\max}$: IG with local search beats ILS by $4.2\%$ (also at equal CPU) against $0.15\%$ &
\ref{sec:f3f4-ablation}, \ref{sec:unrelated} \\[2pt]
2. The policy gain is small, real and replicates &
$0.01$--$0.05$ pp on $4{,}160$ held-out and $2{,}690$ new published instances, independent of the state; $-0.011\%$ on seed averages of 12 classes &
\ref{sec:external}, \ref{sec:f3f4-ablation} \\[2pt]
3. It is limited by headroom and by what an online learner can capture &
Strategy headroom $0.04\%$ ($\mathrm{P}||C_{\max}$) vs.\ $0.38\%$ ($\mathrm{R}||C_{\max}$); with real headroom (destruction size) neither more budget nor pre-training makes Q-Learning beat offline validation; at $1000\times50$ it captures a $1.7\%$ headroom in one distribution and is worse than random selection in two; for the destruction size at $700$ benchmark instances an online or pre-trained agent is worse than random choice ($+0.15\%$) and than validation ($+0.25$ to $+0.36\%$) &
\ref{sec:f3f4-ablation}, \ref{sec:unrelated} \\[2pt]
4. A learning-free recipe is the best method &
Tuned $d$ plus a start from the best constructive heuristic: $74.7\%$ provably optimal ($48.5\%$ for MULTIFIT), improves $26$--$43\%$ of the instances; $89.0\%$ with the CPU time of its local-search hybrid, which is then worse by $0.034\%$ &
\ref{sec:exact-ref}, \ref{sec:recipe} \\[2pt]
5. Granularity of the times decides IG vs.\ ILS &
Holds on six distributions, at equal CPU and with a strong start in both mechanisms; a one-feature rule exploits it &
\ref{sec:nm-structural}, \ref{sec:dispatch} \\[2pt]
6. The attribution is problem-dependent in magnitude &
Policy effect $13\times$ larger for unrelated machines; ILS policy harmful in both; MILP better than the IG there, up to $1000\times50$, but on par with the LP-started IG; replicated on the standard benchmark (700 instances): mechanism $2.1\%$ at equal CPU, policy $0.03$--$0.06\%$ not confirmed, validation beats the policy by $0.07\%$; an LP-started IG is $1.18\%$ better at equal CPU and $0.97\%$ above the best known (baseline $2.23\%$), $0.25\%$ with a restricted MILP &
\ref{sec:unrelated} \\
\bottomrule
\end{tabular}}
\end{table}

\subsection{Experimental Setup}
\label{sec:setup}

All experiments use the synthetic instance generator of
Section~\ref{sec:problem} ($p_j \sim \mathrm{Uniform}[1,100]$ unless
noted otherwise) across twelve instance classes, identified by
$(n,m)$: six were used throughout the project --
$(20,4)$, $(20,6)$, $(30,8)$, $(50,10)$, $(75,15)$ and $(100,20)$, with
$n/m$ ratios between $3.3$ and $5.0$ -- and six were added later to
probe larger scales and a wider range of $n/m$: $(40,5)$, $(60,20)$,
$(150,15)$, $(150,30)$, $(250,25)$ and $(300,20)$, with $n/m$ ratios
between $3.0$ and $15.0$. Each replication uses an independently
sampled random instance (distinct seed), not a repeated run on a
fixed instance, so that reported statistics reflect generalisation
across instances of a class, not just variance of the search
procedure.

The method roster grew from 8 to 23 named methods over the course of
the project. In its final form it comprises: LPT (constructive
baseline, Section~\ref{sec:problem}); three published
baselines reimplemented for this study, MULTIFIT
\citep{CoffmanGareyJohnson1978}, a simulated-annealing method after
\citet{LiPerkins2024} (\texttt{sa\_li\_perkins}), and a Bee Colony
Optimisation method after \citet{DavidovicEtal2012} (\texttt{bco});
eight ILS variants at the standard budget (\texttt{ils\_random},
\texttt{ils\_ucb}, \texttt{ils\_rf}, \texttt{ils\_ql},
\texttt{ils\_ql\_ws}, plus three mechanism-ablation controls,
\texttt{ils\_pm}, \texttt{ils\_ap} and \texttt{ils\_sarsa}, that
replace the Q-Learning agent with, respectively, Probability Matching,
Adaptive Pursuit and SARSA); four single-operator ILS controls
(\texttt{ils\_op1\_only} through \texttt{ils\_op4\_only}) that restrict
the local search to one operator; two extended-budget ILS variants
(\texttt{ils\_rand\_hi}, \texttt{ils\_ql\_hi}, Supplementary Section~\ref{S-sec:ils-framework});
and the five-method \texttt{ig\_*} family (Section~\ref{sec:igvariants}):
\texttt{ig\_random}, \texttt{ig\_ucb}, \texttt{ig\_ucb\_ctx},
\texttt{ig\_ql\_flat} and \texttt{ig\_ql}. All 23 methods were run on
all 12 classes, single-threaded, one CPU core per run, giving
$22{,}070$ individual runs in total (\texttt{results.csv}). All
methods, instance generation and results aggregation are implemented
in Python, using NumPy \citep{HarrisEtal2020} and pandas
\citep{McKinney2010} for numerical and tabular data handling. Later
batteries (LDM and seeded variants, state, reward and hyperparameter
variants, other distributions, seeds, equal CPU time) reuse the same
instances or generators and are described where they are used; each was
specified in the research diary, with hashes of the code, before it was
run (Supplementary Section~\ref{S-sec:prereg-record}).

Two replication regimes coexist, reflecting this chronological growth
rather than a deliberate factorial design. Table~\ref{tab:replication}
summarises them: the twelve original-ILS-family methods (the eight
standard-budget ILS variants plus the four single-operator controls)
were evaluated with $N=30$ independent instances on the six original
classes, matching the project's earliest phase
(Supplementary Section~\ref{S-sec:ils-chapter}); every other method-class combination
-- LPT, the three published baselines, the two extended-budget ILS-HI
variants, the entire \texttt{ig\_*} family, and \emph{all} methods on
the six later classes -- used $N=100$. The smallest class, $(20,4)$,
is the single exception that was never revisited: the original-ILS
methods there remain at $N=30$, and even \texttt{ig\_ql} itself was
only ever run there at $N=30$ (its four ablation siblings, introduced
later, were run there directly at $N=100$). This is not an oversight:
in this class essentially every method already reaches the lower bound
(the mean GAP of \texttt{ig\_ql} is $0.00\%$), leaving negligible room
for any method to differentiate itself, so re-running it at higher $N$
was not prioritised.

\begin{table}[ht]
\centering
\caption{Replications ($N$) by method group and instance-class group.
         $^\dagger$\texttt{ig\_ql} itself follows the top row at
         $(20,4)$; its four ablation siblings (introduced later)
         follow the bottom row throughout.}
\label{tab:replication}
{\small
\begin{tabular}{@{}lccc@{}}
\toprule
Method group & $(20,4)$ & 5 other original classes & 6 later classes \\
\midrule
Original ILS family (12 methods)$^\dagger$ & 30 & 30 & 100 \\
LPT, 3 published, 2 ILS-HI variants & 30 & 100 & 100 \\
\texttt{ig\_*} family (5 methods) & 100 & 100 & 100 \\
\bottomrule
\end{tabular}}
\end{table}

The standard iteration budget is $T=1{,}000$ for all methods except
MULTIFIT, which is a one-shot constructive procedure with no iteration
budget (Section~\ref{sec:igvariants}); the extended-budget \texttt{-HI}
variants use between $T=1{,}500$ and $T=8{,}000$ depending on the class,
re-derived from the fraction of the 81-state space visited at least
three times (Section~\ref{sec:qlm}) rather than a single fixed
multiplier. Unless stated otherwise, percentage gaps are
relative to the classical lower bound
$\mathrm{LB} = \max(\max_j p_j, \lceil \sum_j p_j / m \rceil)$
(Section~\ref{sec:problem}); Sections~\ref{sec:exact-ref} and
\ref{sec:external} also use a certified bound and exact optima.

\subsection{Headline Result: \texttt{ig\_ql} versus ILS and LPT}
\label{sec:f1-headline}

Table~\ref{tab:f1_headline} reports paired, two-sided Wilcoxon
signed-rank tests \citep{Wilcoxon1945}, computed with SciPy
\citep{VirtanenEtal2020} and Holm-Bonferroni corrected
\citep{Holm1979} within this family of $8 \times 12 = 96$ comparisons,
following standard practice for comparing metaheuristics
\citep{DerracEtal2011}, for \texttt{ig\_ql} against LPT and
the seven strongest ILS variants (\texttt{ils\_random},
\texttt{ils\_ucb}, \texttt{ils\_rf}, \texttt{ils\_ql},
\texttt{ils\_rand\_hi}, \texttt{ils\_ql\_ws}, \texttt{ils\_ql\_hi}),
across all 12 classes, together with the paired rank-biserial effect
size \citep{Kerby2014}. Of the 96 comparisons, 72 remain significant
after the Holm correction (79 before correction; 7 lose significance under the more
conservative family-wise procedure), and every one of the 72
significant comparisons favours \texttt{ig\_ql}: it is never
significantly beaten by any ILS variant or by LPT, in any class.

% Auto-generated by analysis/tables.py::latex_wilcoxon_family_table (Entrada 75)
\begin{table}[t]
\centering
\caption{Wilcoxon signed-rank test ($p$-values adjusted by Holm-Bonferroni) and effect size (paired rank-biserial correlation, in parentheses) for \textsc{IG-QL} against each competing method. Bold: $p_{\mathrm{holm}} < 0.05$. $\downarrow$ = \textsc{IG-QL} better; $\uparrow$ = \textsc{IG-QL} worse; \textit{tie} = all pairs identical; N/A = fewer than 3 valid pairs.}
\label{tab:f1_headline}
{\scriptsize\setlength{\tabcolsep}{3pt}
\resizebox{\textwidth}{!}{%
\begin{tabular}{rrcccccccc}
\toprule
$n$ & $m$ & \multicolumn{1}{c}{\textsc{IG-QL} vs \textsc{LPT}} & \multicolumn{1}{c}{\textsc{IG-QL} vs \textsc{ILS-Rand}} & \multicolumn{1}{c}{\textsc{IG-QL} vs \textsc{ILS-UCB}} & \multicolumn{1}{c}{\textsc{IG-QL} vs \textsc{ILS-RF}} & \multicolumn{1}{c}{\textsc{IG-QL} vs \textsc{ILS-QL}} & \multicolumn{1}{c}{\textsc{IG-QL} vs \textsc{ILS-Rand-HI}} & \multicolumn{1}{c}{\textsc{IG-QL} vs \textsc{ILS-QL-WS}} & \multicolumn{1}{c}{\textsc{IG-QL} vs \textsc{ILS-QL-HI}} \\
\midrule
20 & 4 & \textbf{0.0001}$\downarrow$ {\scriptsize(-1.00)} & 1.0000$\downarrow$ {\scriptsize(-1.00)} & \textbf{0.0383}$\downarrow$ {\scriptsize(-1.00)} & 1.0000$\downarrow$ {\scriptsize(-1.00)} & 1.0000$\downarrow$ {\scriptsize(-1.00)} & 1.0000$\downarrow$ {\scriptsize(-1.00)} & 1.0000$\downarrow$ {\scriptsize(-1.00)} & 1.0000$\downarrow$ {\scriptsize(-1.00)} \\
20 & 6 & \textbf{0.0000}$\downarrow$ {\scriptsize(-1.00)} & \textbf{0.0265}$\downarrow$ {\scriptsize(-1.00)} & \textbf{0.0014}$\downarrow$ {\scriptsize(-0.97)} & \textbf{0.0127}$\downarrow$ {\scriptsize(-1.00)} & \textbf{0.0081}$\downarrow$ {\scriptsize(-0.93)} & \textbf{0.0000}$\downarrow$ {\scriptsize(-0.85)} & \textbf{0.0127}$\downarrow$ {\scriptsize(-1.00)} & \textbf{0.0000}$\downarrow$ {\scriptsize(-0.93)} \\
\midrule
30 & 8 & \textbf{0.0000}$\downarrow$ {\scriptsize(-1.00)} & \textbf{0.0390}$\downarrow$ {\scriptsize(-0.75)} & \textbf{0.0004}$\downarrow$ {\scriptsize(-1.00)} & \textbf{0.0036}$\downarrow$ {\scriptsize(-1.00)} & \textbf{0.0055}$\downarrow$ {\scriptsize(-0.88)} & \textbf{0.0002}$\downarrow$ {\scriptsize(-0.78)} & \textbf{0.0036}$\downarrow$ {\scriptsize(-1.00)} & \textbf{0.0000}$\downarrow$ {\scriptsize(-0.92)} \\
\midrule
40 & 5 & \textbf{0.0000}$\downarrow$ {\scriptsize(-1.00)} & 1.0000$\uparrow$ {\scriptsize(-0.24)} & \textbf{0.0008}$\downarrow$ {\scriptsize(-0.89)} & 1.0000$\downarrow$ {\scriptsize(-0.42)} & 1.0000$\uparrow$ {\scriptsize(-0.06)} & 1.0000$\uparrow$ {\scriptsize(0.48)} & 1.0000$\uparrow$ {\scriptsize(0.11)} & 0.4968$\uparrow$ {\scriptsize(0.80)} \\
\midrule
50 & 10 & \textbf{0.0000}$\downarrow$ {\scriptsize(-1.00)} & \textbf{0.0013}$\downarrow$ {\scriptsize(-0.98)} & \textbf{0.0003}$\downarrow$ {\scriptsize(-1.00)} & \textbf{0.0023}$\downarrow$ {\scriptsize(-0.98)} & \textbf{0.0055}$\downarrow$ {\scriptsize(-0.90)} & \textbf{0.0004}$\downarrow$ {\scriptsize(-0.71)} & \textbf{0.0062}$\downarrow$ {\scriptsize(-0.95)} & \textbf{0.0000}$\downarrow$ {\scriptsize(-0.79)} \\
\midrule
60 & 20 & \textbf{0.0000}$\downarrow$ {\scriptsize(-1.00)} & \textbf{0.0000}$\downarrow$ {\scriptsize(-1.00)} & \textbf{0.0000}$\downarrow$ {\scriptsize(-1.00)} & \textbf{0.0000}$\downarrow$ {\scriptsize(-1.00)} & \textbf{0.0000}$\downarrow$ {\scriptsize(-1.00)} & \textbf{0.0000}$\downarrow$ {\scriptsize(-1.00)} & \textbf{0.0000}$\downarrow$ {\scriptsize(-1.00)} & \textbf{0.0000}$\downarrow$ {\scriptsize(-1.00)} \\
\midrule
75 & 15 & \textbf{0.0000}$\downarrow$ {\scriptsize(-1.00)} & \textbf{0.0024}$\downarrow$ {\scriptsize(-0.94)} & \textbf{0.0003}$\downarrow$ {\scriptsize(-0.98)} & \textbf{0.0055}$\downarrow$ {\scriptsize(-0.92)} & \textbf{0.0104}$\downarrow$ {\scriptsize(-0.87)} & \textbf{0.0000}$\downarrow$ {\scriptsize(-0.75)} & \textbf{0.0109}$\downarrow$ {\scriptsize(-0.82)} & \textbf{0.0000}$\downarrow$ {\scriptsize(-0.76)} \\
\midrule
100 & 20 & \textbf{0.0000}$\downarrow$ {\scriptsize(-1.00)} & \textbf{0.0030}$\downarrow$ {\scriptsize(-0.91)} & \textbf{0.0010}$\downarrow$ {\scriptsize(-0.98)} & \textbf{0.0047}$\downarrow$ {\scriptsize(-0.88)} & \textbf{0.0044}$\downarrow$ {\scriptsize(-0.88)} & \textbf{0.0000}$\downarrow$ {\scriptsize(-0.77)} & \textbf{0.0044}$\downarrow$ {\scriptsize(-0.90)} & \textbf{0.0000}$\downarrow$ {\scriptsize(-0.82)} \\
\midrule
150 & 15 & \textbf{0.0000}$\downarrow$ {\scriptsize(-1.00)} & 0.1347$\downarrow$ {\scriptsize(-0.43)} & \textbf{0.0000}$\downarrow$ {\scriptsize(-0.85)} & 0.0631$\downarrow$ {\scriptsize(-0.45)} & 0.2607$\downarrow$ {\scriptsize(-0.40)} & 1.0000$\downarrow$ {\scriptsize(-0.09)} & 0.1805$\downarrow$ {\scriptsize(-0.41)} & 1.0000$\downarrow$ {\scriptsize(-0.27)} \\
150 & 30 & \textbf{0.0000}$\downarrow$ {\scriptsize(-1.00)} & \textbf{0.0000}$\downarrow$ {\scriptsize(-0.96)} & \textbf{0.0000}$\downarrow$ {\scriptsize(-0.99)} & \textbf{0.0000}$\downarrow$ {\scriptsize(-0.98)} & \textbf{0.0000}$\downarrow$ {\scriptsize(-0.97)} & \textbf{0.0000}$\downarrow$ {\scriptsize(-0.91)} & \textbf{0.0000}$\downarrow$ {\scriptsize(-0.96)} & \textbf{0.0000}$\downarrow$ {\scriptsize(-0.92)} \\
\midrule
250 & 25 & \textbf{0.0000}$\downarrow$ {\scriptsize(-1.00)} & \textbf{0.0000}$\downarrow$ {\scriptsize(-0.78)} & \textbf{0.0000}$\downarrow$ {\scriptsize(-0.95)} & \textbf{0.0000}$\downarrow$ {\scriptsize(-0.78)} & \textbf{0.0001}$\downarrow$ {\scriptsize(-0.76)} & \textbf{0.0001}$\downarrow$ {\scriptsize(-0.74)} & \textbf{0.0001}$\downarrow$ {\scriptsize(-0.76)} & \textbf{0.0004}$\downarrow$ {\scriptsize(-0.71)} \\
\midrule
300 & 20 & \textbf{0.0000}$\downarrow$ {\scriptsize(-1.00)} & 1.0000$\downarrow$ {\scriptsize(-0.34)} & \textbf{0.0015}$\downarrow$ {\scriptsize(-0.73)} & 0.6465$\downarrow$ {\scriptsize(-0.40)} & 0.6856$\downarrow$ {\scriptsize(-0.40)} & 1.0000$\downarrow$ {\scriptsize(-0.32)} & 0.8624$\downarrow$ {\scriptsize(-0.37)} & 0.9166$\downarrow$ {\scriptsize(-0.36)} \\
\bottomrule
\end{tabular}}}
\end{table}


Against LPT specifically, the advantage is significant and large in
all 12 classes (effect sizes close to $-1$, meaning \texttt{ig\_ql}
wins essentially every paired instance). Restricting attention to the
six strongest ILS competitors -- \texttt{ils\_ql}, \texttt{ils\_ql\_hi},
\texttt{ils\_ql\_ws}, \texttt{ils\_rand\_hi}, \texttt{ils\_rf} and
\texttt{ils\_ucb}, the methods a reviewer is most likely to ask about
-- $52$ of the $72$ comparisons ($6$ competitors $\times$ $12$
classes) are Holm-significant within the family of $96$, and
\emph{all 52} favour \texttt{ig\_ql}: not one of these stronger
comparisons ever reverses. The $24$ non-significant cells of the
family of $96$ fall in four classes, six cells each -- $(20,4)$,
$(40,5)$, $(150,15)$ and $(300,20)$ -- where the gap available to any
method is already tiny (mean GAP below $0.3\%$ for both sides) and the
paired differences, while occasionally nominally in favour of an ILS
variant, are too small relative to their variance to survive
correction. This is the most robust and least qualified result of the
present study.

\subsection{\texttt{ig\_ql} versus Published Baselines}
\label{sec:f2-sota}

Table~\ref{tab:f2_sota} repeats the same protocol for \texttt{ig\_ql}
against the three reimplemented published baselines
(Section~\ref{sec:igvariants}): MULTIFIT \citep{CoffmanGareyJohnson1978},
\texttt{sa\_li\_perkins} \citep{LiPerkins2024} and \texttt{bco}
\citep{DavidovicEtal2012}. Of the $3 \times 12 = 36$ comparisons, 35
are Holm-significant -- the sole exception, discussed below, is a
genuine statistical tie -- but, unlike Section~\ref{sec:f1-headline},
the direction is not uniform even among the significant cases.

% Auto-generated by analysis/tables.py::latex_wilcoxon_family_table (Entrada 75)
\begin{table}[t]
\centering
\caption{Wilcoxon signed-rank test ($p$-values adjusted by Holm-Bonferroni) and effect size (paired rank-biserial correlation, in parentheses) for \textsc{IG-QL} against each competing method. Bold: $p_{\mathrm{holm}} < 0.05$. $\downarrow$ = \textsc{IG-QL} better; $\uparrow$ = \textsc{IG-QL} worse; \textit{tie} = all pairs identical; N/A = fewer than 3 valid pairs.}
\label{tab:f2_sota}
{\scriptsize\setlength{\tabcolsep}{3pt}
\begin{tabular}{rrccc}
\toprule
$n$ & $m$ & \multicolumn{1}{c}{\textsc{IG-QL} vs \textsc{MultiFit}} & \multicolumn{1}{c}{\textsc{IG-QL} vs \textsc{SA (Li\&Perkins)}} & \multicolumn{1}{c}{\textsc{IG-QL} vs \textsc{BCO}} \\
\midrule
20 & 4 & \textbf{0.0000}$\downarrow$ {\scriptsize(-1.00)} & \textbf{0.0000}$\downarrow$ {\scriptsize(-1.00)} & \textbf{0.0000}$\downarrow$ {\scriptsize(-1.00)} \\
20 & 6 & \textbf{0.0000}$\downarrow$ {\scriptsize(-1.00)} & \textbf{0.0000}$\downarrow$ {\scriptsize(-1.00)} & \textbf{0.0000}$\downarrow$ {\scriptsize(-1.00)} \\
\midrule
30 & 8 & \textbf{0.0000}$\downarrow$ {\scriptsize(-1.00)} & \textbf{0.0000}$\downarrow$ {\scriptsize(-1.00)} & \textbf{0.0000}$\downarrow$ {\scriptsize(-1.00)} \\
\midrule
40 & 5 & \textbf{0.0000}$\downarrow$ {\scriptsize(-0.99)} & \textbf{0.0000}$\downarrow$ {\scriptsize(-1.00)} & \textbf{0.0000}$\downarrow$ {\scriptsize(-1.00)} \\
\midrule
50 & 10 & \textbf{0.0005}$\downarrow$ {\scriptsize(-0.63)} & \textbf{0.0000}$\downarrow$ {\scriptsize(-1.00)} & \textbf{0.0000}$\downarrow$ {\scriptsize(-1.00)} \\
\midrule
60 & 20 & 0.6116$\downarrow$ {\scriptsize(-0.09)} & \textbf{0.0000}$\downarrow$ {\scriptsize(-1.00)} & \textbf{0.0000}$\downarrow$ {\scriptsize(-1.00)} \\
\midrule
75 & 15 & \textbf{0.0000}$\uparrow$ {\scriptsize(0.96)} & \textbf{0.0000}$\downarrow$ {\scriptsize(-1.00)} & \textbf{0.0000}$\downarrow$ {\scriptsize(-1.00)} \\
\midrule
100 & 20 & \textbf{0.0000}$\uparrow$ {\scriptsize(1.00)} & \textbf{0.0000}$\downarrow$ {\scriptsize(-1.00)} & \textbf{0.0000}$\downarrow$ {\scriptsize(-1.00)} \\
\midrule
150 & 15 & \textbf{0.0000}$\uparrow$ {\scriptsize(1.00)} & \textbf{0.0000}$\downarrow$ {\scriptsize(-1.00)} & \textbf{0.0000}$\downarrow$ {\scriptsize(-1.00)} \\
150 & 30 & \textbf{0.0000}$\uparrow$ {\scriptsize(1.00)} & \textbf{0.0000}$\downarrow$ {\scriptsize(-1.00)} & \textbf{0.0000}$\downarrow$ {\scriptsize(-1.00)} \\
\midrule
250 & 25 & \textbf{0.0000}$\uparrow$ {\scriptsize(1.00)} & \textbf{0.0000}$\downarrow$ {\scriptsize(-1.00)} & \textbf{0.0000}$\downarrow$ {\scriptsize(-1.00)} \\
\midrule
300 & 20 & \textbf{0.0000}$\uparrow$ {\scriptsize(1.00)} & \textbf{0.0000}$\downarrow$ {\scriptsize(-1.00)} & \textbf{0.0000}$\downarrow$ {\scriptsize(-1.00)} \\
\bottomrule
\end{tabular}}
\end{table}


\texttt{ig\_ql} beats \texttt{sa\_li\_perkins} and \texttt{bco} in all
12 classes each, with effect sizes at or near $-1$ throughout and mean
gap differences that widen, not narrow, with instance size (e.g.\ mean
GAP against \texttt{bco} grows from $0.72\%$ at $(20,4)$ to $19.8\%$ at
$(60,20)$ and $25.4\%$ at $(150,30)$). Against MULTIFIT the picture
reverses at scale: \texttt{ig\_ql} wins five of the six classes with
$n \le 60$ -- the exception, $(60,20)$, is a genuine statistical tie
($p_{\mathrm{holm}}=0.612$, effect size $-0.09$), not a loss -- but
loses all six classes with $n \ge 75$, all six losses Holm-significant
with effect size $+1$. This crossover is structurally explained, not
anomalous. MULTIFIT is a one-shot bin-packing-style
construction with no iteration budget at all -- it runs once to
completion regardless of $n$ -- and its absolute GAP improves with
scale under classical bin-packing approximation theory (from
$0.76\%$ at the smallest classes to $0.17\%$ at the largest); it is
also, on the evidence surveyed independently in a 2025 systematic
review of P$\|$C$_{\max}$ heuristics that itself co-authors the BCO
baseline used here \citep[Section~5]{OstojicEtAl2025}, one of the
best-regarded classical constructive heuristics for this problem, not
a deliberately weak comparator. \texttt{ig\_ql}, by contrast, uses a
\emph{fixed} iteration budget ($T=1{,}000$) independent of $n$, so its
absolute GAP grows slowly but steadily with scale (from near $0\%$ to
$0.67\%$ across the same range) -- enough, at the two-to-four largest
classes, to fall behind MULTIFIT's scale-improving curve. \texttt{bco}
degrades with scale for the opposite structural reason: its per-outer-%
iteration cost is $B \times n$ constructive moves for a fixed
population parameter $B=5$ \citep{DavidovicEtal2012}, so under the
same total-moves budget used to equalise it with the other methods
(Section~\ref{sec:igvariants}) its number of outer iterations collapses
to $I = \max(1, T/(Bn))$, i.e.\ one or two outer iterations once
$n \gtrsim 80$ -- a consequence of the equal-total-moves budget rule
adopted in this study, not of the published algorithm, which uses
$I=100$ outer iterations (the sensitivity analysis below relaxes this
rule). The crossover with MULTIFIT, by contrast, was tested directly:
running \texttt{ig\_ql} at up to $10\times$ its standard budget
($T=10{,}000$) on the most extreme class, $(300,20)$, shows that the
crossover is not an artefact of an unfairly small budget for
\texttt{ig\_ql}: mean GAP improves only
from $0.094\%$ ($T=1{,}000$) to $0.077\%$ ($T=5{,}000$), then plateaus
exactly at $0.077\%$ at $T=10{,}000$ -- an order of magnitude more
search than the standard protocol still leaves \texttt{ig\_ql} roughly
$30\times$ worse than MULTIFIT's $0.003\%$ on this class. The net
result -- \texttt{ig\_ql} wins 16 of 18 comparisons with published baselines at the
original six classes and loses specifically to a 1978 heuristic with
no search budget at all, only at large $n$ -- is reported without
euphemism: where MULTIFIT wins it does so at roughly one-fiftieth of the
CPU time of \texttt{ig\_ql} (at $(300,20)$), so for $n \ge 75$ on this
instance family MULTIFIT dominates \texttt{ig\_ql} in both quality and
time. Section~\ref{sec:exact-ref} adds LDM and an IG that starts from the
best constructive solution.

\paragraph{Sensitivity of the SA and BCO comparisons to budget rules}
Two design features of these baselines could flatter \texttt{ig\_ql}
and are therefore examined directly. First, \texttt{sa\_li\_perkins}
halts once its best solution is within $1\%$ of the lower bound; in
the classes $(150,15)$, $(250,25)$ and $(300,20)$ the initial LPT
solution already satisfies this, so no search takes place, and in
every class with $n \ge 60$ its mean GAP equals that of LPT to three
decimals. Second, \texttt{bco} receives only
$I = \max(1,\lfloor T/(5n)\rfloor)$ outer iterations, against $I=100$
in the original publication. Table~\ref{tab:f7_sota_sensitivity}
reruns both under relaxed rules on the same instances: SA without its
stopping rule (all $T=1{,}000$ iterations); BCO with $I$ calibrated so
that its CPU time equals that of \texttt{ig\_ql} ($I$ from $31$ at
$n=20$ to $6$ at $n=300$); and BCO with the published $I=100$, which
costs $3.3$--$14.8\times$ the CPU time of \texttt{ig\_ql} (measured in
the same session). \texttt{ig\_ql} remains Holm-significantly better
than all three variants in all 12 classes (36 of 36 tests), so the
\emph{direction} of the comparison does not depend on the budget rule.
The \emph{magnitudes} do: BCO's mean GAP at $(100,20)$ falls from
$14.9\%$ to $1.7\%$ at $I=100$, so the large gaps reported for BCO
above (up to $25\%$) largely reflect the equal-total-moves rule rather
than the capability of the algorithm; and SA without its stopping rule
still leaves the LPT solution unchanged in $(60,20)$, $(75,15)$ and
$(100,20)$, where a single random move or swap almost never improves a
balanced LPT schedule. MULTIFIT has no budget and is unaffected.
Finally, the comparison does not depend on the learning component:
\texttt{ig\_random}, tested against the original SA, BCO and MULTIFIT,
beats SA and BCO in all 12 classes and, against MULTIFIT, wins five
classes, ties one and loses six -- the same pattern as \texttt{ig\_ql}.
Re-executing \texttt{ig\_ql} on all 850 instance replications of this
analysis reproduced every stored makespan exactly, confirming that the
iteration-budget protocol is deterministic.

% Auto-generated by analysis/make_table_sota_sensitivity.py (Entrada 92)
\begin{table}[t]
\centering
\caption{Mean GAP (\%) of \textsc{IG-QL} and of the published baselines under alternative budget settings, on the same replications ($N=100$ for $n \le 100$; $N=30$ for $n \ge 150$ and for $(20,4)$). SA and BCO are the versions used in Table~\ref{tab:f2_sota} (BCO with $I=\max(1,\lfloor T/(5n)\rfloor)$ outer iterations). $^{\dagger}$SA without its early-stopping rule (all $T=1{,}000$ iterations). $^{\ddagger}$BCO with $I$ calibrated so that its CPU time equals that of \textsc{IG-QL} ($I$ from 31 at $n=20$ to 6 at $n=300$). $^{\S}$BCO with the published $I=100$ (3.3--14.8$\times$ the CPU time of \textsc{IG-QL}). Bold: best in row. \textsc{IG-QL} is Holm-significantly better than each of the three variants $^{\dagger}$, $^{\ddagger}$ and $^{\S}$ in all 12 classes (36 of 36 tests).}
\label{tab:f7_sota_sensitivity}
{\scriptsize\setlength{\tabcolsep}{3pt}
\begin{tabular}{@{}rrrrrrrrr@{}}
\toprule
$n$ & $m$ & \textsc{IG-QL} & \textsc{MultiFit} & SA & SA$^{\dagger}$ & BCO & BCO$^{\ddagger}$ & BCO$^{\S}$ \\
\midrule
20 & 4 & \textbf{0.00} & 0.76 & 0.73 & 0.38 & 0.72 & 0.38 & 0.15 \\
20 & 6 & \textbf{0.43} & 1.71 & 1.58 & 1.57 & 3.01 & 1.81 & 1.36 \\
30 & 8 & \textbf{0.21} & 1.12 & 2.10 & 2.10 & 3.98 & 2.37 & 1.39 \\
40 & 5 & \textbf{0.02} & 0.20 & 0.52 & 0.25 & 0.74 & 0.32 & 0.14 \\
50 & 10 & \textbf{0.29} & 0.39 & 1.72 & 1.72 & 4.06 & 1.77 & 1.02 \\
60 & 20 & \textbf{0.85} & 0.89 & 6.49 & 6.49 & 19.80 & 6.73 & 3.98 \\
75 & 15 & 0.56 & \textbf{0.24} & 2.02 & 2.02 & 11.76 & 2.67 & 1.31 \\
100 & 20 & 0.67 & \textbf{0.17} & 2.01 & 2.01 & 14.93 & 4.00 & 1.66 \\
150 & 15 & 0.18 & \textbf{0.03} & 0.42 & 0.42 & 7.12 & 1.12 & 0.37 \\
150 & 30 & 0.87 & \textbf{0.08} & 2.21 & 2.21 & 24.60 & 6.81 & 2.18 \\
250 & 25 & 0.16 & \textbf{0.00} & 0.36 & 0.36 & 10.54 & 3.02 & 0.53 \\
300 & 20 & 0.09 & \textbf{0.00} & 0.16 & 0.16 & 4.73 & 1.26 & 0.23 \\
\bottomrule
\end{tabular}}
\end{table}


\subsection{Stronger Baselines and an Exact Reference}
\label{sec:exact-ref}

Two objections to Section~\ref{sec:f2-sota} call for a direct answer: the
baseline set omits the Largest Differencing Method (LDM), and
\texttt{ig\_ql} starts from LPT rather than from the best constructive
solution at hand. Moreover, a GAP to the classical bound mixes
algorithmic error with bound slack. Both points were examined with
analyses specified before the runs (hypotheses, decision rules and
hashes of the frozen code are in the research diary).

\paragraph{A certified bound}
The linear relaxation of the arc-flow model
\citep{MradSouayah2018,GharbiBamatraf2022} yields a valid lower bound,
LB$^{\mathrm{AF}}$: the smallest capacity whose relaxation needs at most
$m$ bins. It was solved with the interior-point method of HiGHS (median
$0.2$\,s). Against the proven optima of the iAF method
\citep{OstojicEtAl2025} on 150 published instances it never exceeded the
optimum and equalled it in 148 ($98.7\%$); brute-force checks confirm
its validity on small instances. On the $1{,}200$ synthetic instances it
equals the best makespan found by the 23 original methods in $903$ cases
($75.3\%$), so for those the optimum is \emph{proven}.

\paragraph{Stronger baselines}
LDM is the multiway Karmarkar--Karp differencing heuristic
\citep{KarmarkarKarp1983,MichielsEtal2007}. \textsc{IG-QL}$^{+}$ is
\texttt{ig\_ql} started from the best of LPT, MULTIFIT and LDM (never
worse than its seed by construction; same budget and parameters).
Table~\ref{tab:f12_stronger} gives the mean GAP against
LB$^{\mathrm{AF}}$ and the share of provably optimal solutions on the
$1{,}130$ synthetic instances for which all methods were run, and
Table~\ref{tab:f13_tests} the eight pre-specified tests, on the
synthetic instances and on the 6\,780 published ones. All eight
hypotheses are confirmed after Holm correction, but their magnitudes
qualify the picture of Section~\ref{sec:f2-sota}.

% Auto-generated by analysis/make_table_stronger_baselines.py (Entrada 102)
\begin{table}[t]
\centering
\caption{Stronger baselines against the certified arc-flow lower bound (LB$^{\mathrm{AF}}$): mean GAP (\%) and share of instances solved \emph{provably} to optimality (makespan $=$ LB$^{\mathrm{AF}}$). \textsc{IG-QL}$^{+}$ is \textsc{IG-QL} started from the best of LPT, \textsc{MultiFit} and \textsc{LDM}. Classes sorted by $n$; $N=100$ per class ($30$ for $(20,4)$). Bold: lowest GAP in the row (ties all bold).}
\label{tab:f12_stronger}
{\scriptsize\setlength{\tabcolsep}{3.5pt}
\begin{tabular}{@{}rr rrrr rrrr@{}}
\toprule
 & & \multicolumn{4}{c}{GAP vs LB$^{\mathrm{AF}}$ (\%)} & \multicolumn{4}{c}{Provably optimal (\%)} \\
\cmidrule(lr){3-6}\cmidrule(lr){7-10}
$n$ & $m$ & \textsc{MultiFit} & \textsc{LDM} & \textsc{IG-QL} & \textsc{IG-QL}$^{+}$ & \textsc{MultiFit} & \textsc{LDM} & \textsc{IG-QL} & \textsc{IG-QL}$^{+}$ \\
\midrule
20 & 4 & 0.759 & 0.493 & \textbf{0.000} & 0.026 & 10 & 27 & 100 & 93 \\
20 & 6 & 1.454 & 2.231 & \textbf{0.169} & 0.173 & 10 & 1 & 74 & 73 \\
30 & 8 & 1.123 & 1.789 & 0.211 & \textbf{0.165} & 6 & 3 & 59 & 68 \\
40 & 5 & 0.199 & 0.044 & 0.021 & \textbf{0.009} & 45 & 82 & 91 & 96 \\
50 & 10 & 0.390 & 0.692 & 0.291 & \textbf{0.243} & 24 & 12 & 31 & 42 \\
60 & 20 & 0.851 & 5.550 & 0.820 & \textbf{0.775} & 9 & 2 & 12 & 9 \\
75 & 15 & 0.239 & 0.898 & 0.559 & \textbf{0.223} & 44 & 4 & 4 & 46 \\
100 & 20 & \textbf{0.175} & 1.047 & 0.666 & \textbf{0.175} & 56 & 1 & 1 & 56 \\
150 & 15 & 0.039 & 0.026 & 0.166 & \textbf{0.016} & 80 & 87 & 21 & 92 \\
150 & 30 & \textbf{0.101} & 1.480 & 0.816 & \textbf{0.101} & 74 & 0 & 0 & 74 \\
250 & 25 & \textbf{0.002} & 0.026 & 0.175 & \textbf{0.002} & 99 & 87 & 21 & 99 \\
300 & 20 & 0.003 & 0.004 & 0.074 & \textbf{0.000} & 98 & 97 & 44 & 100 \\
\midrule
\multicolumn{2}{c}{All (1{,}130)} & 0.425 & 1.233 & 0.351 & \textbf{0.167} & 48.5 & 34.0 & 34.3 & 69.3 \\
\bottomrule
\end{tabular}}
\end{table}

% Auto-generated by analysis/make_table_stronger_baselines.py (Entrada 102)
\begin{table}[t]
\centering
\caption{Pre-specified paired tests with stronger baselines (Wilcoxon signed-rank, two-sided; Holm over all eight tests). $\bar d$: mean relative makespan difference $100(C_A-C_B)/C_B$ in \% (negative: first method better); W/L/T: instances where the first method is better/worse/tied; $r$: rank-biserial correlation (negative: first better). All eight hypotheses are confirmed, but P1 is practically negligible ($|\bar d|<0.05\%$).}
\label{tab:f13_tests}
{\scriptsize\setlength{\tabcolsep}{3.5pt}
\begin{tabular}{@{}llrrrrrr@{}}
\toprule
Set & Comparison & $N$ & W/L/T & $\bar d$ (\%) & $r$ & $p_{\mathrm{Holm}}$ & Neglig. \\
\midrule
Synthetic & P1: \textsc{IG-QL}$^{+}$ vs.\ \textsc{IG-Random}$^{+}$ & 1{,}130 & 46/22/1062 & -0.010 & -0.35 & 0.036 & yes \\
 & P2: \textsc{IG-QL}$^{+}$ vs.\ \textsc{IG-QL} & 1{,}130 & 548/48/534 & -0.183 & -0.88 & $3\times10^{-62}$ & no \\
 & P3: \textsc{IG-QL}$^{+}$ vs.\ best constructive & 1{,}130 & 218/0/912 & -0.172 & -1.00 & $5\times10^{-37}$ & no \\
 & P4: \textsc{IG-QL} vs.\ \textsc{LDM} & 1{,}130 & 580/238/312 & -0.845 & -0.69 & $8\times10^{-88}$ & no \\
\midrule
External & P1: \textsc{IG-QL}$^{+}$ vs.\ \textsc{IG-Random}$^{+}$ & 6{,}780 & 959/754/5067 & -0.013 & -0.08 & $2\times10^{-9}$ & yes \\
 & P2: \textsc{IG-QL}$^{+}$ vs.\ \textsc{IG-QL} & 6{,}780 & 3717/825/2238 & -0.180 & -0.74 & $<10^{-300}$ & no \\
 & P3: \textsc{IG-QL}$^{+}$ vs.\ best constructive & 6{,}780 & 2275/0/4505 & -0.339 & -1.00 & $<10^{-300}$ & no \\
 & P4: \textsc{IG-QL} vs.\ \textsc{LDM} & 6{,}780 & 3045/2420/1315 & -0.975 & 0.01 & $2\times10^{-189}$ & no \\
\bottomrule
\end{tabular}}
\end{table}


\emph{(a)} LDM improves on LPT (GAP$^{\mathrm{AF}}$ $1.23\%$ against
$2.18\%$) but is worse than MULTIFIT ($0.43\%$; MULTIFIT wins $585$ of
the $725$ decisive synthetic instances), so it does not alter the ranking
of the constructive heuristics. Plain \texttt{ig\_ql} is better than LDM
on average (P4) but loses to it in $21\%$ of the synthetic and $36\%$ of
the published instances, and in $67$, $70$ and $53$ of the $100$ instances
of $(150,15)$, $(250,25)$ and $(300,20)$.
\emph{(b)} Starting IG from the best constructive solution recovers the
large-$n$ classes: \textsc{IG-QL}$^{+}$ ties with the best constructive
heuristic in every instance with $n \geq 100$ (hence matches MULTIFIT
there) and improves on it in $19\%$ (synthetic) and $34\%$ (published)
of the instances, never worsening it (P3). Its mean GAP$^{\mathrm{AF}}$
is $0.17\%$, against $0.35\%$ for \texttt{ig\_ql} and $0.43\%$ for
MULTIFIT, and $69.3\%$ of its solutions are provably optimal, against
$34.3\%$ and $48.5\%$ (P2). Seeding is not uniformly beneficial: at
$(20,4)$ the LPT start is marginally better ($0.000\%$ against
$0.026\%$).
\emph{(c)} The Q-Learning policy adds almost nothing to the strong
start: $\bar d = -0.010\%$ (synthetic) and $-0.013\%$ (published) against
random destruction from the same seed (P1), inside the $\pm0.05\%$
practical-equivalence margin although statistically detectable with
$6{,}780$ instances.

\paragraph{A symmetric strong start}
Starting only the IG from the best constructive solution could favour it, so ILS
was started the same way (\textsc{ILS-QL}$^{+}$ and \textsc{ILS-Random}$^{+}$;
Table~\ref{tab:f17_ils_seeded}). The Q-Learning policy does not help the seeded ILS
either ($+0.015\%$ synthetic, $+0.024\%$ published: slightly worse than random
selection). With the same strong start the mechanism contrast keeps its dependence
on granularity: \textsc{IG-QL}$^{+}$ is better on the synthetic instances
($-0.093\%$; $164$ wins against $10$) and on average on the published ones
($-0.065\%$), but \textsc{ILS-QL}$^{+}$ wins more published instances ($554$ against
$232$), in particular in family B (322 of 400) while the IG wins by magnitude in C
and E. The ILS improves on the best constructive heuristic in $10.1\%$ of the
synthetic and $52.5\%$ of the published instances, against $19.3\%$ and $32.9\%$
for the IG. For $n \geq 150$ the two seeded methods are identical in every synthetic
instance: nothing is left to improve on the constructive solution there.

% Auto-generated by analysis/make_tables_cycle2.py (Entrada 122)
\begin{table}[t]
\centering
\caption{IG and ILS both started from the best of LPT, \textsc{MultiFit} and \textsc{LDM} (``$^{+}$''); Wilcoxon, Holm over six. Synthetic: 12 classes $\times$ 100; published: 1{,}500 instances (B 400, C 400, E 400, F 300). $\bar d$ in \% of the makespan (negative: first better).}
\label{tab:f17_ils_seeded}
{\scriptsize\setlength{\tabcolsep}{3.5pt}
\begin{tabular}{@{}llrrrr@{}}
\toprule
Set & Comparison & $N$ & W/L/T & $\bar d$ (\%) & $p_{\mathrm{Holm}}$ \\
\midrule
Synthetic & I1: \textsc{ILS-QL}$^{+}$ vs.\ \textsc{ILS-Random}$^{+}$ & 1{,}200 & 15/41/1144 & +0.015 & 0.001 \\
 & I2: \textsc{IG-QL}$^{+}$ vs.\ \textsc{ILS-QL}$^{+}$ & 1{,}200 & 164/10/1026 & -0.093 & $1\times10^{-28}$ \\
 & I3: \textsc{ILS-QL}$^{+}$ vs.\ best constructive & 1{,}130 & 114/0/1016 & -0.077 & $8\times10^{-20}$ \\
\midrule
Published & I1: \textsc{ILS-QL}$^{+}$ vs.\ \textsc{ILS-Random}$^{+}$ & 1{,}500 & 224/305/971 & +0.024 & $4\times10^{-5}$ \\
 & I2: \textsc{IG-QL}$^{+}$ vs.\ \textsc{ILS-QL}$^{+}$ & 1{,}500 & 232/554/714 & -0.065 & 0.005 \\
 & I3: \textsc{ILS-QL}$^{+}$ vs.\ best constructive & 1{,}500 & 788/0/712 & -0.257 & $8\times10^{-130}$ \\
\bottomrule
\end{tabular}}
\end{table}


\paragraph{Quality versus time}
Figure~\ref{fig:pareto} places the methods on the quality--time plane;
CPU times come from one sequential session on an otherwise idle machine
(those stored with the earlier batteries were measured on different
machines and are not used for ratios). \textsc{IG-QL}$^{+}$ lies on the
Pareto frontier in all three size ranges, at $0.2$--$0.7$\,s per run.
For the same $T$, \texttt{ils\_ql} costs between $1.3\times$ ($n=20$) and
$24.8\times$ ($n=300$) the CPU time of \texttt{ig\_ql}, and BCO with its
published $I=100$ costs $3.2$--$16\times$; both are dominated in the
panel means. MULTIFIT costs $50$--$500$ times less than \texttt{ig\_ql}.
Because an ILS iteration is costlier, comparing ILS and IG at equal $T$
favours ILS in CPU terms; Supplementary Section~\ref{S-sec:eqcpu} checks that this
does not drive the conclusions.

\begin{figure}[t]
\centering
\includegraphics[width=\textwidth]{fig_pareto.pdf}
\caption{Quality--time trade-off on the synthetic classes, in three size
ranges (each marker is the mean over the four classes of the panel; GAP
against the certified bound; CPU time is the geometric mean of the
class medians, log scale). The dashed line is the Pareto frontier.
\textsc{IG-QL}$^{+}$: \texttt{ig\_ql} started from the best of LPT,
MULTIFIT and LDM.}
\label{fig:pareto}
\end{figure}

\subsection{Mechanism versus Policy: Ablation and Isolation}
\label{sec:f3f4-ablation}

Sections~\ref{sec:f1-headline}--\ref{sec:f2-sota} establish that
\texttt{ig\_ql} wins; this section asks \emph{why}, using the ablation
ladder and isolation control introduced in Section~\ref{sec:igvariants}.

\paragraph{The Q-Learning-specific contribution is small, and invisible
to per-class tests} Table~\ref{tab:f3_ablation} repeats the Wilcoxon/Holm
protocol for \texttt{ig\_ql} against its four ablation siblings --
\texttt{ig\_random}, \texttt{ig\_ucb}, \texttt{ig\_ucb\_ctx} and
\texttt{ig\_ql\_flat} -- across all 12 classes ($4 \times 12 = 48$
comparisons).

% Auto-generated by analysis/tables.py::latex_wilcoxon_family_table (Entrada 75)
\begin{table}[t]
\centering
\caption{Wilcoxon signed-rank test ($p$-values adjusted by Holm-Bonferroni) and effect size (paired rank-biserial correlation, in parentheses) for \textsc{IG-QL} against each competing method. Bold: $p_{\mathrm{holm}} < 0.05$. $\downarrow$ = \textsc{IG-QL} better; $\uparrow$ = \textsc{IG-QL} worse; \textit{tie} = all pairs identical; N/A = fewer than 3 valid pairs.}
\label{tab:f3_ablation}
{\scriptsize\setlength{\tabcolsep}{3pt}
\begin{tabular}{rrcccc}
\toprule
$n$ & $m$ & \multicolumn{1}{c}{\textsc{IG-QL} vs \textsc{IG-Random}} & \multicolumn{1}{c}{\textsc{IG-QL} vs \textsc{IG-UCB}} & \multicolumn{1}{c}{\textsc{IG-QL} vs \textsc{IG-UCB-Ctx}} & \multicolumn{1}{c}{\textsc{IG-QL} vs \textsc{IG-QL-Flat}} \\
\midrule
20 & 4 & 1.0000$\downarrow$ {\scriptsize(-1.00)} & 1.0000$\downarrow$ {\scriptsize(-1.00)} & 1.0000$\downarrow$ {\scriptsize(-1.00)} & 1.0000$\downarrow$ {\scriptsize(-1.00)} \\
20 & 6 & 1.0000$\uparrow$ {\scriptsize(0.22)} & 1.0000$\uparrow$ {\scriptsize(0.21)} & 1.0000$\uparrow$ {\scriptsize(0.35)} & 1.0000$\downarrow$ {\scriptsize(-0.23)} \\
\midrule
30 & 8 & 0.6329$\downarrow$ {\scriptsize(-0.60)} & 1.0000$\downarrow$ {\scriptsize(-0.15)} & 1.0000$\downarrow$ {\scriptsize(-0.09)} & 1.0000$\downarrow$ {\scriptsize(-0.08)} \\
\midrule
40 & 5 & 1.0000$\uparrow$ {\scriptsize(0.24)} & 1.0000$\uparrow$ {\scriptsize(0.20)} & 1.0000$\downarrow$ {\scriptsize(-0.06)} & 1.0000$\downarrow$ {\scriptsize(-0.43)} \\
\midrule
50 & 10 & 1.0000$\downarrow$ {\scriptsize(-0.41)} & 1.0000$\downarrow$ {\scriptsize(-0.38)} & 1.0000$\downarrow$ {\scriptsize(-0.36)} & 1.0000$\downarrow$ {\scriptsize(-0.26)} \\
\midrule
60 & 20 & 1.0000$\uparrow$ {\scriptsize(0.06)} & 1.0000$\uparrow$ {\scriptsize(-0.02)} & 1.0000$\downarrow$ {\scriptsize(-0.20)} & 1.0000$\downarrow$ {\scriptsize(-0.05)} \\
\midrule
75 & 15 & 1.0000$\downarrow$ {\scriptsize(-0.07)} & 1.0000$\downarrow$ {\scriptsize(-0.13)} & 1.0000$\uparrow$ {\scriptsize(0.34)} & 1.0000$\uparrow$ {\scriptsize(0.03)} \\
\midrule
100 & 20 & 1.0000$\downarrow$ {\scriptsize(-0.36)} & 1.0000$\downarrow$ {\scriptsize(-0.21)} & 1.0000$\downarrow$ {\scriptsize(-0.31)} & 0.2236$\downarrow$ {\scriptsize(-0.56)} \\
\midrule
150 & 15 & 1.0000$\uparrow$ {\scriptsize(-0.02)} & 1.0000$\downarrow$ {\scriptsize(-0.11)} & 1.0000$\downarrow$ {\scriptsize(-0.41)} & 1.0000$\downarrow$ {\scriptsize(-0.73)} \\
150 & 30 & 1.0000$\downarrow$ {\scriptsize(-0.20)} & 0.6329$\downarrow$ {\scriptsize(-0.49)} & 1.0000$\downarrow$ {\scriptsize(-0.27)} & 1.0000$\downarrow$ {\scriptsize(-0.33)} \\
\midrule
250 & 25 & 1.0000$\downarrow$ {\scriptsize(-0.59)} & 1.0000$\downarrow$ {\scriptsize(-0.36)} & 1.0000$\downarrow$ {\scriptsize(-0.37)} & 0.9274$\downarrow$ {\scriptsize(-0.76)} \\
\midrule
300 & 20 & 1.0000$\uparrow$ {\scriptsize(0.21)} & 1.0000$\uparrow$ {\scriptsize(0.03)} & 1.0000$\downarrow$ {\scriptsize(-0.47)} & 1.0000$\uparrow$ {\scriptsize(0.52)} \\
\bottomrule
\end{tabular}}
\end{table}


Only 7 of the 48 raw $p$-values fall below $0.05$, and \emph{none}
survive the Holm correction. The closest case, \texttt{ig\_ql} versus
\texttt{ig\_ql\_flat} at $(100,20)$, has $p_{\mathrm{holm}} = 0.224$.
Class by class, then, no cell of the $2\times2$ design that crosses
state-conditioning with TD-learning (Section~\ref{sec:igvariants})
shows a detectable advantage: neither UCB1 with state-conditioning
alone (\texttt{ig\_ucb\_ctx}) nor tabular Q-Learning collapsed to a
single aggregated state (\texttt{ig\_ql\_flat}) is separable from
uniform random selection (\texttt{ig\_random}), and the full
\texttt{ig\_ql} agent is not separable from any of them. This is a
statement about \emph{detectability}, not about exact equality. With
$N \le 100$ instances per class and a Holm correction over 48 tests,
this design can only expose effects far larger than the ones that
separate \texttt{ig\_ql} from the ILS family
(Sections~\ref{sec:f1-headline}--\ref{sec:f2-sota}); an effect of a few
percent of the (already small) GAP is below its resolution. Section~\ref{sec:external}
shows that such an effect does exist: pooled within each published
benchmark family, \texttt{ig\_ql} is slightly but significantly better
than \texttt{ig\_random}, and a test pre-specified before running
4\,160 held-out instances confirms it in the three families that
admit a confirmatory test. The accurate summary of the ablation is
therefore that the learning-specific contribution is \emph{small} -- a
few hundredths of a percentage point of GAP -- while the search
mechanism accounts for the overwhelming majority of the advantage over
the ILS family, as the paragraph on the search mechanism below shows.

\paragraph{The original state was nearly degenerate; a more informative one
does not change the conclusion}
In the original runs the load features are constant: in the IG runs
examined, $f_1$ and $f_2$ fall in bin 0 at every step and the agent
visits only $5$--$7$ distinct states on average, which encode progress
and stagnation (Section~\ref{sec:state}). The ablation above therefore
tests state-conditioning on that reduced information. To check whether
this drives the null, two recalibrated states were tested with the same
agent and hyperparameters (Table~\ref{tab:f10_state}): thresholds at the
terciles of each feature in calibration runs on separate instances
(``cal''), and features expressed relative to their value at the start
of the run (``rel''), which removes the instance-specific component that
cannot help an agent that learns from scratch in each run. Both arms
make the features vary (about $18$ distinct states per IG run and $27$
per ILS run; the most occupied bin of $f_1$ or $f_2$ holds $54$--$72\%$
of the steps, against $100\%$). With $N=100$ instances in each of the 12
classes (hypotheses, decision rules and calibration thresholds were
written down in the research diary before the runs),
\texttt{ig\_ql} with either state is statistically indistinguishable
from random destruction ($-3.1\%$ and $-3.6\%$ pooled GAP,
$p_{\mathrm{holm}}=0.28$), and \texttt{ils\_ql} with either state is
significantly \emph{worse} than random operator selection ($+4.8\%$ and
$+4.7\%$, $p_{\mathrm{holm}}=0.009$). Neither differs from the original
state (IG: $p \geq 0.79$; ILS, six original classes:
$p_{\mathrm{holm}} \geq 0.56$). The only gain is against the stateless
TD learner (\texttt{ig\_ql\_flat}), which is itself nominally worse than
random (exploratory: $+1.9\%$, $p=0.07$): conditioning on an informative
state removes that penalty ($-4.9\%$ and $-5.4\%$,
$p_{\mathrm{holm}}=0.012$ and $0.004$) but does not add to random
selection. A more informative state therefore does not rescue the
learned policy and, in ILS, hurts it; the paragraphs below test other
explanations (headroom, reward signal, hyperparameters).

% Auto-generated by analysis/make_table_state_recalibration.py (Entrada 98)
\begin{table}[t]
\centering
\caption{Q-Learning with a recalibrated state ($N=100$ instances in each of the 12 classes). ``States'' is the mean number of distinct states visited per run; ``Top bin'' is the mean share of the most occupied bin of $f_1$ / $f_2$ (1.00 in the original state, whose load features stay in one bin at every step). $\Delta$ is the pooled relative change in GAP against random selection with the same search mechanism (negative = Q-Learning better); $p_{\mathrm{holm}}$ is Holm-adjusted over the two arms of each mechanism (paired Wilcoxon on class-standardised differences). ``cal'': thresholds at the terciles of each feature in calibration runs; ``rel'': features relative to their value at the start of the run, with tercile thresholds. The last column compares the IG arms with the stateless TD learner (\texttt{ig\_ql\_flat}).}
\label{tab:f10_state}
{\scriptsize\setlength{\tabcolsep}{4pt}
\begin{tabular}{@{}llrcccc@{}}
\toprule
Mechanism & State & States & Top bin ($f_1$ / $f_2$) & $\Delta$ vs random (\%) & $p_{\mathrm{holm}}$ & $\Delta$ vs stateless (\%) \\
\midrule
IG & original & 5--7 & 1.00 / 1.00 & see Table~\ref{tab:f3_ablation} & & \\
IG & cal & 18.1 & 0.68 / 0.62 & $-$3.1 & 0.278 & $-$4.9 (0.012) \\
IG & rel & 18.0 & 0.72 / 0.66 & $-$3.6 & 0.278 & $-$5.4 (0.004) \\
\midrule
ILS & cal & 26.8 & 0.61 / 0.58 & $+$4.8 & 0.009 & -- \\
ILS & rel & 26.8 & 0.64 / 0.54 & $+$4.7 & 0.009 & -- \\
\bottomrule
\end{tabular}}
\end{table}


\paragraph{The search mechanism itself is what wins}
Table~\ref{tab:f4_isolation} isolates the contribution of the destroy-%
and-reconstruct search mechanism from any adaptive selection at all,
by comparing \texttt{ig\_random} (uniform random destruction strategy,
zero learning) against \texttt{ils\_random} (uniform random operator
selection, zero learning) -- the same comparison, at the same fixed
budget, with learning removed from both sides.

% Auto-generated by analysis/tables.py::latex_wilcoxon_family_table (Entrada 75)
\begin{table}[t]
\centering
\caption{Wilcoxon signed-rank test ($p$-values adjusted by Holm-Bonferroni) and effect size (paired rank-biserial correlation, in parentheses) for \textsc{IG-Random} against each competing method. Bold: $p_{\mathrm{holm}} < 0.05$. $\downarrow$ = \textsc{IG-Random} better; $\uparrow$ = \textsc{IG-Random} worse; \textit{tie} = all pairs identical; N/A = fewer than 3 valid pairs.}
\label{tab:f4_isolation}
{\scriptsize\setlength{\tabcolsep}{3pt}
\begin{tabular}{rrc}
\toprule
$n$ & $m$ & \multicolumn{1}{c}{\textsc{IG-Random} vs \textsc{ILS-Rand}} \\
\midrule
20 & 4 & 0.4427$\downarrow$ {\scriptsize(-1.00)} \\
20 & 6 & \textbf{0.0030}$\downarrow$ {\scriptsize(-0.96)} \\
\midrule
30 & 8 & \textbf{0.0207}$\downarrow$ {\scriptsize(-0.71)} \\
\midrule
40 & 5 & 0.4427$\downarrow$ {\scriptsize(-0.38)} \\
\midrule
50 & 10 & \textbf{0.0003}$\downarrow$ {\scriptsize(-0.98)} \\
\midrule
60 & 20 & \textbf{0.0000}$\downarrow$ {\scriptsize(-1.00)} \\
\midrule
75 & 15 & \textbf{0.0018}$\downarrow$ {\scriptsize(-0.83)} \\
\midrule
100 & 20 & \textbf{0.0018}$\downarrow$ {\scriptsize(-0.88)} \\
\midrule
150 & 15 & \textbf{0.0258}$\downarrow$ {\scriptsize(-0.42)} \\
150 & 30 & \textbf{0.0000}$\downarrow$ {\scriptsize(-0.95)} \\
\midrule
250 & 25 & \textbf{0.0003}$\downarrow$ {\scriptsize(-0.67)} \\
\midrule
300 & 20 & \textbf{0.0207}$\downarrow$ {\scriptsize(-0.55)} \\
\bottomrule
\end{tabular}}
\end{table}


Ten of the 12 comparisons are Holm-significant, and all ten favour
\texttt{ig\_random}; the two exceptions are again the near-ceiling
classes $(20,4)$ and $(40,5)$. Combined with the ablation of
Table~\ref{tab:f3_ablation}, this is the central mechanistic finding
of the study: \texttt{ig\_ql}'s advantage over the ILS family comes
overwhelmingly from the destroy-and-reconstruct search mechanism
itself, not from the Q-Learning policy layered on top of it. This
decomposition mirrors, for \texttt{ig\_ql}, exactly the question that
Supplementary Section~\ref{S-sec:ils-chapter} asks of \texttt{ils\_ql} -- there, the
same style of analysis found that Q-Learning does not beat random
operator selection at all; here, it finds that Q-Learning adds only a
small improvement to a search mechanism that is already winning on its
own.

\paragraph{Which part of the search structure?}
The advantage of the IG over the ILS mixes two things: the destruction--reconstruction
cycle and the absence of an exhaustive local search. To separate them, a hybrid
(IG with a short best-improvement descent after each reconstruction, \texttt{igls})
was run on the same instances (Table~\ref{tab:f25_igls_p}). The pure IG already beats
the ILS ($-0.59\%$, $381$ wins against $61$ on the $780$ paired instances); the local
search adds a further $0.15\%$ to the IG ($436$ against $16$ wins; the descent is not
a hindrance), and the hybrid beats the ILS by $0.72\%$; Q-Learning adds nothing to the
hybrid ($-0.004\%$). The destruction--reconstruction cycle thus supplies most of the
advantage ($0.59$ of $0.72$ points), the local search a smaller complement, and the
policy none; for unrelated machines the pure IG does not beat the ILS and the combination
is what wins (Section~\ref{sec:unrelated}).

% Auto-generated by analysis/make_tables_cycle3.py (Entrada 125)
\begin{table}[t]
\centering
\caption{$\mathrm{P}||C_{\max}$: is the active ingredient the destruction--reconstruction cycle or the local search? IG+LS adds a best-improvement descent (inter-swap, reinsert; two passes) after each reconstruction ($T=1000$, 12 classes; the ILS has 30 instances in the six original classes). $\bar d$ in \% of the makespan (negative: first better); Wilcoxon, Holm over three.}
\label{tab:f25_igls_p}
{\scriptsize\setlength{\tabcolsep}{3.5pt}
\begin{tabular}{@{}lrrrrr@{}}
\toprule
Comparison & $N$ & W/L/T & $\bar d$ (\%) & 90\% interval & $p_{\mathrm{Holm}}$ \\
\midrule
P1: IG+LS vs.\ IG (random selection) & 1{,}200 & 436/16/748 & -0.150 & [-0.162, -0.139] & $7\times10^{-65}$ \\
P2: IG+LS vs.\ ILS (random selection) & 780 & 459/16/305 & -0.722 & [-0.796, -0.650] & $2\times10^{-75}$ \\
P3: IG+LS, \textsc{QL} vs.\ random selection & 1{,}200 & 98/77/1025 & -0.004 & [-0.011, +0.004] & 0.646 \\
\bottomrule
\end{tabular}}
\end{table}


That $0.15\%$ is, however, a comparison at an equal number of iterations, and the
descent multiplies the CPU time per iteration by a factor between $2.1$ and $52.4$
depending on the class. Giving the IG without local search the same CPU time as the hybrid
($T_{\mathrm{eq}}$ between $2{,}064$ and $52{,}425$ iterations, calibrated sequentially;
Table~\ref{tab:f28_igls_eq}) reduces the advantage of the local search to $0.010\%$
(90\% interval $0.002$--$0.018$; the pre-specified test is formally significant,
$p=0.015$, but the interval lies inside $\pm0.05\%$): $93\%$ of the gain is an effect of
budget. In four classes ($20\times6$, $30\times8$, $50\times10$, $75\times15$) the hybrid
is still ahead by $0.04$--$0.07\%$, whereas in the other eight the plain IG is level or ahead
(by up to $0.04\%$, in $250\times25$). Both reach about $55\%$ of
provably optimal solutions, against $37\%$ for the IG at $T=1{,}000$. On
identical machines, then, the active ingredient is destruction--reconstruction and
the number of iterations it is given; the descent mostly substitutes for iterations.

% Auto-generated by analysis/make_tables_cycle5.py (Entrada 129)
\begin{table}[t]
\centering
\caption{Is the gain of the local search in the hybrid of Table~\ref{tab:f25_igls_p} an effect of search structure or of budget? \textsc{IG} (random selection, $d=0.3n$, no learning) with $T_{\mathrm{eq}}$ iterations has the same CPU time as the hybrid with $T=1000$ ($T_{\mathrm{eq}}=1000K$, $K$ the sequential CPU ratio per class, from $2.1$ to $52.4$; checked again on an idle machine: ratios $0.96$--$1.02$). 12 classes $\times$ 100 instances; $\bar d$ in \% of the makespan (negative: first better); Wilcoxon. Q1 was the pre-specified test; Q2 and the first row are descriptive.}
\label{tab:f28_igls_eq}
{\scriptsize\setlength{\tabcolsep}{3.5pt}
\begin{tabular}{@{}lrrrrr@{}}
\toprule
Comparison & $N$ & W/L/T & $\bar d$ (\%) & 90\% interval & $p$ \\
\midrule
IG+LS vs.\ IG at equal iterations ($T=1000$) & 1{,}200 & 436/16/748 & -0.150 & -- & -- \\
Q1: IG+LS ($T=1000$) vs.\ IG at \emph{equal CPU} ($T_{\mathrm{eq}}$) & 1{,}200 & 112/104/984 & -0.010 & [-0.018, -0.002] & 0.015 \\
Q2: IG at equal CPU vs.\ IG ($T=1000$) & 1{,}200 & 421/0/779 & -0.140 & [-0.150, -0.129] & $1\times10^{-70}$ \\
\midrule
\multicolumn{6}{l}{Certified bound (1{,}200 instances): provably optimal / mean GAP (\%)} \\
\multicolumn{6}{l}{\quad IG ($T=1000$) 36.9\% / 0.347; IG+LS ($T=1000$) 55.2\% / 0.196; IG at equal CPU 54.5\% / 0.206} \\
\bottomrule
\end{tabular}}
\end{table}


\paragraph{Why the policy adds so little: headroom}
Rival explanations for the small policy effect were tested with analyses
specified beforehand (Table~\ref{tab:f14_policy}). If the three destruction
strategies were interchangeable, no selector would have anything to learn.
Running IG with each strategy fixed shows that the best fixed strategy per
class (chosen on the same instances, hence optimistic) is no better than
random selection ($\bar d = 0.000\%$), and that even a selector that knew the
best strategy of every instance would gain only $0.037\%$. One strategy
(removing the longest jobs) is by itself $0.06$--$3.7\%$ worse than random,
since it is deterministic and stalls, yet avoiding it is worth less than
$0.04\%$: the search saturates within $T=1{,}000$ even when a third of the
iterations are spent on it.

\paragraph{Signal and pre-training}
The sparse reward (about three non-zero values per run) could hide a real
effect. A dense reward, the relative reduction of the sum of squared machine
loads, makes Q-Learning significantly \emph{worse} than random ($+0.084\%$,
$p_{\mathrm{Holm}}<10^{-20}$), because improving the balance is not improving
the makespan; pre-training the table of each class on 50 independent
instances leaves it equivalent to random ($+0.002\%$). The null is thus not
an artefact of reward sparsity or of learning from scratch, at least for these
designs; other dense rewards were not tried.

\paragraph{Hyperparameters}
Changing $\alpha$, $\gamma$ or the $\varepsilon$ decay changes nothing
($|\bar d| \leq 0.023\%$). The destruction size matters more, and the policy
matters more when it is small: with $d=0.1n$, $0.2n$ and $0.4n$ the advantage
over the random control with the same $d$ is $0.11\%$, $0.06\%$ and $0.04\%$.
The default $d=0.3n$ was never tuned on the evaluation instances: at $d=0.1n$
the mean GAP of \texttt{ig\_ql} falls from $0.62\%$ to $0.19\%$ at $(100,20)$
and from $0.84\%$ to $0.34\%$ at $(150,30)$, so the crossover with MULTIFIT of
Section~\ref{sec:f2-sota} moves to larger $n$ rather than disappearing.

\paragraph{Seeds}
The main batteries run each method once per instance. Ten algorithm seeds on
20 instances in each of the twelve classes (240 instances; Supplementary
Section~\ref{S-sec:seedvar}) show that the instance explains a median of
$44$--$45\%$ (IG) and $72$--$74\%$ (ILS) of the variance of the GAP, and that averaging over
seeds sharpens the resolution: the policy effect is $\bar d=-0.011\%$ (90\%
interval $-0.019$ to $-0.003$, $p=0.009$), real but strictly negligible
(the interval lies inside $\pm0.05\%$), and only one class, $(30,8)$, detects it
individually after correction ($-0.076\%$). The IG-versus-ILS contrast, by
contrast, is $-0.49\%$ (random) and $-0.53\%$ (Q-Learning) on seed averages,
with the same sign as the single-run contrast in $23$ of $24$ class--variant
pairs. The policy effect is thus real but of the order of $0.01$--$0.05\%$ at
default settings (up to $0.11\%$ with small destruction sizes), against
$0.3$--$0.7\%$ for the mechanism.

% Auto-generated by analysis/make_table_policy_robustness.py (Entrada 112)
\begin{table}[t]
\centering
\caption{Effect of the Q-Learning policy over random destruction, under variations of the state, reward, training, hyperparameters and seed averaging. $\bar d$: mean relative makespan difference (\%; negative: Q-Learning better) against the paired random control; 90\% percentile-bootstrap interval. Positive values mean that Q-Learning is \emph{worse} than random. The two rows of the middle block are not learners: they bound the headroom, i.e.\ what any selector among the three destruction strategies could gain (best fixed strategy per class, chosen in-sample; omniscient per-instance choice).}
\label{tab:f14_policy}
{\scriptsize\setlength{\tabcolsep}{4pt}
\begin{tabular}{@{}lrrr@{}}
\toprule
Variant & $N$ & $\bar d$ (\%) & 90\% interval \\
\midrule
Original \texttt{ig\_ql} (sparse reward, default state) & 1{,}130 & -0.016 & [-0.027, -0.006] \\
Informative state (\texttt{cal}) & 1{,}200 & -0.011 & [-0.022, -0.001] \\
Dense reward (load balance) & 1{,}200 & +0.084 & [+0.070, +0.098] \\
Pre-trained per class & 1{,}200 & +0.002 & [-0.009, +0.012] \\
Dense reward + pre-training & 1{,}200 & +0.099 & [+0.084, +0.115] \\
\midrule
Best fixed strategy per class (in-sample) & 1{,}200 & +0.000 & [-0.011, +0.012] \\
Per-instance oracle of the 3 strategies & 1{,}200 & -0.037 & [-0.048, -0.027] \\
\midrule
Destruction size $d=0.1n$ & 300 & -0.112 & [-0.148, -0.075] \\
Destruction size $d=0.2n$ & 300 & -0.057 & [-0.084, -0.032] \\
Destruction size $d=0.4n$ & 300 & -0.037 & [-0.054, -0.020] \\
$\alpha$, $\gamma$, $\varepsilon$-decay (6 settings; range) & 300 & --- & [-0.023, +0.007] \\
Seed-averaged (10 seeds; 3 classes) & 60 & -0.045 & [-0.063, -0.027] \\
\bottomrule
\end{tabular}}
\end{table}


\paragraph{An action space with real headroom}
The headroom argument can be turned around: if the action is the destruction
\emph{size}, where Section~\ref{sec:f3f4-ablation} found up to $0.45\%$ at stake,
a learner should gain. Table~\ref{tab:f15_dsel} reports the pre-specified test
with five fractions $d/n \in \{0.05,\ldots,0.40\}$ and random destruction
($1{,}200$ instances, $T=1{,}000$). The best fraction chosen \emph{by validation}
on 30 separate instances per class beats random choice of $d$ out of sample
($\bar d=-0.089\%$): the headroom is real. Q-Learning over the same five
actions does not capture it ($+0.015\%$; $2\%$ of the headroom), neither does
UCB1 ($+0.007\%$), and Q-Learning is worse than the validated fraction by
$0.103\%$ in $318$ of the $375$ decisive instances. The hypothesis that the
policy matters once there is headroom is therefore refuted for this agent.

% Auto-generated by analysis/make_tables_cycle2.py (Entrada 122)
\begin{table}[t]
\centering
\caption{The action is the destruction size ($d/n\in\{0.05,0.10,0.20,0.30,0.40\}$, random destruction), a space with real headroom. Pre-specified tests on 1{,}200 instances ($T=1000$; Wilcoxon, Holm over four). $\bar d$: mean relative makespan difference (\%; negative: first better); W/L/T: first better/worse/tied. ``Best $d$ by validation'' is chosen per class on 30 separate instances.}
\label{tab:f15_dsel}
{\scriptsize\setlength{\tabcolsep}{3.5pt}
\begin{tabular}{@{}lrrrr@{}}
\toprule
Comparison & W/L/T & $\bar d$ (\%) & 90\% interval & $p_{\mathrm{Holm}}$ \\
\midrule
A1: best $d$ by validation vs.\ random $d$ & 324/47/829 & -0.089 & [-0.100, -0.077] & $6\times10^{-30}$ \\
A2: \textsc{IG-QL} (over $d$) vs.\ random $d$ & 229/213/758 & +0.015 & [+0.001, +0.030] & 0.123 \\
A3: UCB1 (over $d$) vs.\ random $d$ & 165/178/857 & +0.007 & [-0.005, +0.018] & 0.366 \\
A4: \textsc{IG-QL} (over $d$) vs.\ best $d$ by validation & 57/318/825 & +0.103 & [+0.092, +0.116] & $3\times10^{-35}$ \\
\bottomrule
\end{tabular}}
\end{table}


\paragraph{Less budget, more budget, pre-training}
Two further explanations were tested (Tables~\ref{tab:f20_budget}
and~\ref{tab:f21_dsel_budget}). \emph{Budget.} If the policy only matters
where the budget binds, tightening $T$ should reveal it. The headroom does
grow ($-0.103\%$ at $T=100$, $-0.067\%$ at $250$, $-0.037\%$ at $1000$; test
B4 confirmed), but the policy effect does not ($-0.011\%$, $-0.015\%$,
$-0.011\%$; B1--B3 not confirmed): the fraction of the headroom that Q-Learning
captures rises with the horizon ($0.11$, $0.22$, $0.30$), i.e.\ the agent needs
time to learn. \emph{Horizon and cross-instance learning.} With five times the
budget ($T=5000$) Q-Learning over destruction sizes still does not capture it
significantly ($-0.010\%$, interval $-0.028$ to $+0.008$) although headroom remains
($-0.049\%$); and a table pre-trained per
class on 50 independent instances gains $0.032\%$ against random choice
(not confirmed; $0.10$--$0.15\%$ in the two largest classes) and loses $0.063\%$
against the fraction chosen by validation. In this problem the gain available
from the action space is realised by choosing the parameter \emph{offline},
on separate instances, not by the online or the cross-instance tabular
learning tested here (Section~\ref{sec:recipe}).

% Auto-generated by analysis/make_tables_cycle3.py (Entrada 125)
\begin{table}[t]
\centering
\caption{Budget law in $\mathrm{P}||C_{\max}$: policy effect (informative state, \texttt{ig\_ql\_cal} vs.\ \texttt{ig\_random}) and headroom of the three destruction strategies (omniscient per-instance choice vs.\ random) as the budget $T$ shrinks (12 classes $\times$ 100 instances; 90\% bootstrap intervals). Last rows: pre-specified tests (Holm over four): B1, B2 policy effect at $T=100$ and $250$; B3 policy effect, $T=100$ vs.\ $1000$ (12 class pairs); B4 headroom, $T=100$ vs.\ $1000$.}
\label{tab:f20_budget}
{\scriptsize\setlength{\tabcolsep}{3.5pt}
\begin{tabular}{@{}rrrrr@{}}
\toprule
$T$ & Policy $\bar d$ (\%) & 90\% interval & Headroom $\bar d$ (\%) & Capture \\
\midrule
100 & -0.011 & [-0.031, +0.010] & -0.103 & 0.11 \\
250 & -0.015 & [-0.029, +0.001] & -0.067 & 0.22 \\
1000 & -0.011 & [-0.022, -0.001] & -0.037 & 0.30 \\
\midrule
\multicolumn{2}{l}{B1: policy at $T=100$} & -0.011 & & $p_{\mathrm{Holm}}=1.000$ \\
\multicolumn{2}{l}{B2: policy at $T=250$} & -0.015 & & $p_{\mathrm{Holm}}=0.533$ \\
\multicolumn{2}{l}{B3: policy, $T=100$ vs.\ $1000$} & +0.000 & & $p_{\mathrm{Holm}}=1.000$ \\
\multicolumn{2}{l}{B4: headroom, $T=100$ vs.\ $1000$} & -0.066 & & $p_{\mathrm{Holm}}=0.027$ \\
\bottomrule
\end{tabular}}
\end{table}

% Auto-generated by analysis/make_tables_cycle3.py (Entrada 125)
\begin{table}[t]
\centering
\caption{Can the policy capture the headroom of the destruction-size action space with five times the budget or with cross-instance pre-training? 6 classes $\times$ 50 instances; pre-training: one Q-table per class trained on 50 independent instances. $\bar d$ in \% of the makespan (negative: first better); Wilcoxon, Holm over four.}
\label{tab:f21_dsel_budget}
{\scriptsize\setlength{\tabcolsep}{3.5pt}
\begin{tabular}{@{}lrrrr@{}}
\toprule
Comparison & W/L/T & $\bar d$ (\%) & 90\% interval & $p_{\mathrm{Holm}}$ \\
\midrule
D1: \textsc{IG-QL} (over $d$) vs.\ random $d$, $T=5000$ & 37/29/234 & -0.010 & [-0.028, +0.008] & 0.788 \\
D2: best $d$ by validation vs.\ random $d$, $T=5000$ & 41/1/258 & -0.049 & [-0.062, -0.037] & $5\times10^{-7}$ \\
D3: pre-trained \textsc{IG-QL} vs.\ random $d$, $T=1000$ & 63/37/200 & -0.032 & [-0.054, -0.010] & 0.371 \\
D4: pre-trained \textsc{IG-QL} vs.\ best $d$ by validation, $T=1000$ & 19/67/214 & +0.063 & [+0.044, +0.083] & $2\times10^{-6}$ \\
\bottomrule
\end{tabular}}
\end{table}


\paragraph{No variance reduction}
Median-centred Levene \citep{BrownForsythe1974} and
Fligner--Killeen \citep{FlignerKilleen1976} tests, Holm-corrected over the
ten classes then available, reject equality of variance in none of them
(the single nominal case, $(300,20)$, has a variance ratio of $0.996$),
so Q-Learning does not detectably reduce the variance of the GAP either.

\subsection{A Practical Recipe: Tuned Destruction Size and a Strong Start}
\label{sec:recipe}

The previous sections leave a constructive question: if learning does not pay,
what does? The evidence points to two cheap, learning-free ingredients: a
strong start (Section~\ref{sec:exact-ref}) and a destruction size chosen
\emph{offline}, by validation on separate instances. \textsc{IG-Tuned}$^{+}$
combines them: Iterated Greedy with random destruction, $d/n$ chosen per
class on 30 separate instances (per family on 553 published instances that
were never run, for the external data), started from the best of LPT,
MULTIFIT and LDM, with the same budget $T=1{,}000$. Table~\ref{tab:f22_igcv}
reports the pre-specified comparison, with all eight tests confirmed.

% Auto-generated by analysis/make_tables_cycle3.py (Entrada 125)
\begin{table}[t]
\centering
\caption{Tuned-$d$ Iterated Greedy started from the best constructive solution (\textsc{IG-Tuned}$^{+}$: random destruction, $d/n$ chosen on separate instances, no learning). Synthetic: 1{,}130 instances; published: 1{,}500 (B 400, C 400, E 400, F 300; $d/n$ per family chosen on 553 instances never run). $\bar d$ in \% (negative: \textsc{IG-Tuned}$^{+}$ better); Wilcoxon, Holm over eight.}
\label{tab:f22_igcv}
{\scriptsize\setlength{\tabcolsep}{3.5pt}
\begin{tabular}{@{}llrrr@{}}
\toprule
Set & Comparison & W/L/T & $\bar d$ (\%) & $p_{\mathrm{Holm}}$ \\
\midrule
Synthetic & E1: vs.\ \textsc{IG-QL}$^{+}$ (default $d$, learned strategies) & 111/30/989 & -0.034 & $1\times10^{-6}$ \\
 & E2: vs.\ \textsc{MultiFit} & 378/0/752 & -0.288 & $4\times10^{-63}$ \\
 & E3: vs.\ \textsc{LDM} & 707/0/423 & -1.061 & $1\times10^{-116}$ \\
 & E4: vs.\ best constructive & 292/0/838 & -0.207 & $4\times10^{-49}$ \\
\midrule
Published & E1: vs.\ \textsc{IG-QL}$^{+}$ (default $d$, learned strategies) & 408/127/965 & -0.100 & $3\times10^{-30}$ \\
 & E2: vs.\ \textsc{MultiFit} & 1098/0/402 & -0.849 & $3\times10^{-180}$ \\
 & E3: vs.\ \textsc{LDM} & 845/0/655 & -1.145 & $4\times10^{-139}$ \\
 & E4: vs.\ best constructive & 645/0/855 & -0.420 & $1\times10^{-106}$ \\
\bottomrule
\end{tabular}}
\end{table}


Tuning $d$ improves on the default-$d$ IG-QL$^{+}$ ($-0.034\%$ synthetic,
$-0.100\%$ published; the gain is concentrated in $50 \leq n \leq 75$ and in
family C, and nil for $n \geq 150$, where the constructive heuristics are
already hard to improve). \textsc{IG-Tuned}$^{+}$ beats MULTIFIT by $0.29\%$
(synthetic) and $0.85\%$ (published) and LDM by $1.06\%$ and $1.15\%$, never
returns a worse solution than the best constructive heuristic, and improves on
it in $25.8\%$ of the synthetic and $43.0\%$ of the published instances. Against
the certified bound, $74.7\%$ of its solutions are provably optimal (IG-QL$^{+}$:
$69.3\%$; MULTIFIT: $48.5\%$; LDM: $34.0\%$; best constructive: $55.3\%$) and its mean GAP is
$0.133\%$, against $0.167\%$, $0.425\%$ and $1.233\%$. The recipe is a
statement about this problem and these budgets, not about what a better
learner could do; but at this budget it is the strongest method of the study and
contains no learned component.

\paragraph{Does a local search help the recipe? Only at an equal number of iterations}
The hybrid analysis of Section~\ref{sec:f3f4-ablation} suggests adding the descent of
the hybrid to the recipe. Table~\ref{tab:f26_igcvls} tests it with a pre-specified
design, on the same synthetic and published instances. At an equal number of
iterations the local search helps: $-0.036\%$ on the synthetic and $-0.072\%$ on
the published instances (F1, F3), and the recipe with local search reaches $81.2\%$
of provably optimal solutions against $74.7\%$. But the descent is expensive. One
iteration costs between $2.3$ and $82.9$ times more CPU time, depending on the
class (on average $0.35$ against $6.7$ seconds per synthetic instance, and $0.37$
against $37.5$ seconds per published instance). When the recipe \emph{without}
local search is given the same CPU time ($T_{\mathrm{eq}}$ between $2{,}298$ and
$82{,}855$ iterations, calibrated sequentially per class on an idle machine; the timings
of the parallel run, which depend on the load, gave ratios between $0.74$ and $1.17$),
it is better, not worse: the local search loses by $0.034\%$ (F2, $24$ wins against
$115$ losses, $p_{\mathrm{Holm}}=3\times10^{-8}$), and the equal-CPU recipe reaches
$89.0\%$ of provably optimal solutions and a mean GAP of $0.058\%$, against $0.094\%$.
The advantage of the local search is therefore an effect of budget, not of search
structure: the same CPU time spent on more destruction--reconstruction iterations is
better spent. Two consequences follow: the strongest variant of the study is the
learning-free recipe \emph{with a larger iteration budget} (on six classes it reaches $86.8\%$, $92.7\%$ and $97.5\%$ of provably optimal solutions at $T=4{,}000$, $16{,}000$ and $64{,}000$, against $75.5\%$ at $T=1{,}000$; Supplementary Table~\ref{S-tab:s_recipe_budget}); and the results at
$T=1{,}000$ throughout the paper are comparisons at a fixed, small budget, not claims
about what each method can reach.

% Auto-generated by analysis/make_tables_cycle5.py (Entrada 129)
\begin{table}[t]
\centering
\caption{Does adding the local search of the hybrid to the practical recipe help? Both variants use random destruction with $d/n$ chosen on separate instances and start from the best of LPT, MULTIFIT and LDM. IG+LS adds the descent (inter-swap, reinsert; two passes) after each reconstruction ($T=1000$). \textsc{IG-Tuned}$^{+}$ at equal CPU runs the recipe without local search for $T_{\mathrm{eq}}$ iterations, $T_{\mathrm{eq}}=1000K$ with $K$ the CPU ratio measured per class ($K$ from $2.3$ to $82.9$; $T_{\mathrm{eq}}$ from $2{,}298$ to $82{,}855$). Synthetic: 1{,}200 instances (1{,}130 with a constructive reference); published: 1{,}500. $\bar d$ in \% of the makespan (negative: IG+LS better); Wilcoxon, Holm over five.}
\label{tab:f26_igcvls}
{\scriptsize\setlength{\tabcolsep}{3.5pt}
\begin{tabular}{@{}lrrrrr@{}}
\toprule
Comparison & $N$ & W/L/T & $\bar d$ (\%) & 90\% interval & $p_{\mathrm{Holm}}$ \\
\midrule
F1: IG+LS vs.\ IG ($T=1000$), synthetic & 1{,}200 & 103/21/1076 & -0.036 & [-0.045, -0.028] & $5\times10^{-10}$ \\
F2: IG+LS vs.\ IG at \emph{equal CPU} ($T_{\mathrm{eq}}$ iterations), synthetic & 1{,}200 & 24/115/1061 & +0.034 & [+0.026, +0.042] & $3\times10^{-8}$ \\
F3: IG+LS vs.\ IG ($T=1000$), published & 1{,}500 & 576/36/888 & -0.072 & [-0.081, -0.063] & $1\times10^{-76}$ \\
F4: IG+LS vs.\ best constructive, synthetic & 1{,}130 & 340/0/790 & -0.245 & [-0.269, -0.220] & $5\times10^{-57}$ \\
F5: IG+LS vs.\ best constructive, published & 1{,}500 & 818/0/682 & -0.492 & [-0.534, -0.452] & $8\times10^{-135}$ \\
\midrule
\multicolumn{6}{l}{Synthetic, certified bound (1{,}130 instances): provably optimal / mean GAP (\%)} \\
\multicolumn{6}{l}{\quad best constructive 55.3\% / 0.342; recipe ($T=1000$) 74.7\% / 0.133; recipe + local search 81.2\% / 0.094; recipe at equal CPU 89.0\% / 0.058} \\
\bottomrule
\end{tabular}}
\end{table}


\subsection{External Validation on Published Benchmark Instances}
\label{sec:external}

To test whether the \texttt{ig\_ql} advantage generalises to instances
this project did not design, the 13-method core roster (LPT, the three
published baselines, \texttt{ils\_random}, \texttt{ils\_ucb},
\texttt{ils\_ql}, \texttt{ils\_rf} and the five \texttt{ig\_*} variants)
was run once on each of $6{,}780$ instances from the benchmark released
with the review of \citet{OstojicEtAl2025}: a stratified sample of
$2{,}000$ instances from each of families B, C and E (50 generator
classes per family) and all $780$ of F (26 classes), with the same
protocol as Section~\ref{sec:setup}, analysed per generator class with the
Wilcoxon/Holm procedure used throughout (Holm over all classes and
baselines of a family).\footnote{Wide-range instances of family B
($U[1,10^{18}]$ at the extreme) overflow 64-bit load accumulators; the
code switches to arbitrary-precision integers above a safe threshold, with
no change of behaviour on any other instance (verified by an independent
recomputation and regression tests).}

\paragraph{The claim against the published baselines replicates}
\texttt{ig\_ql} against SA, BCO and MULTIFIT is Holm-significant and
favourable in a large majority of tests in all four families: B ($133$ of
$150$ tests significant, $110$ favourable), C ($121/150$, $116$), E
($128/150$, all $128$) and F ($67/78$, $61$). Two caveats of
Section~\ref{sec:f2-sota} carry over: most comparisons are against SA and
BCO, and the informative one is MULTIFIT, which \texttt{ig\_ql} beats on
$1{,}388$ of the $2{,}703$ instances with an exact reference (below) and
loses to on $759$.

\paragraph{The mechanism-specific claims do not replicate uniformly}
Within IG, \texttt{ig\_ql} against its four ablation siblings gives only
$10$ of $704$ Holm-significant comparisons, inconsistent in direction:
favourable in B ($4$ of $5$) and E ($2$ of $2$), none in C, and
unfavourable in all three significant cases in F. These per-subfolder
tests are low-powered; the family-level test below refines the picture.
IG against ILS is strongly family-dependent: ILS wins essentially all
significant comparisons in family B ($157$ of $163$ unfavourable to IG;
its processing times are wide-range and almost duplicate-free, and
Supplementary Section~\ref{S-sec:hybrid} traces the loss to a concrete
cause), IG wins nearly all of family C ($100$ of $102$), and E and F are
mixed ($33$ of $56$ and $13$ of $31$ unfavourable to IG). Neither
mechanism dominates universally: the processing-time profile of an
instance matters more than either fixed choice.

\paragraph{A small Q-Learning effect, detected at family level and
confirmed on held-out instances}
Per-class and per-subfolder tests split the data into many small cells and
correct for all of them, which leaves little power against an effect of a
few percent of the GAP. A family-level analysis of the $6{,}780$ instances,
one paired Wilcoxon test per family, suggested that \texttt{ig\_ql} has a
lower GAP than \texttt{ig\_random} in B, C and E and a marginally higher
one in F. Since the hypothesis arose after looking at the data, it was
tested on $4{,}160$ instances never run before (disjoint from all earlier
ones: $1{,}500$ each from B and C, $1{,}160$ from E; F had none left), with
the hypothesis, the test (paired two-sided Wilcoxon, Holm over three
families) and the decision rule fixed in writing, with hashes of the
analysis files, before any of them was run (a dated diary record, not a
public registry). Table~\ref{tab:f6_holdout_ql} shows that the rule was
met in all three families. The effect is small: $-3.0\%$ and $-3.5\%$ of
the GAP in B and E, matching the exploratory magnitude, but only
$-0.85\%$ in C, where the interval of the mean difference includes zero;
in absolute terms $0.01$--$0.03$ percentage points of GAP. Family F, where
Q-Learning looked slightly \emph{worse} ($+0.5\%$, $p_{\mathrm{holm}}=0.078$),
cannot be tested confirmatorily. Against the other three ablation siblings
the evidence is partial (\texttt{ig\_ql\_flat}: B and E; \texttt{ig\_ucb}
and \texttt{ig\_ucb\_ctx}: E only), so the gain cannot be attributed to
state-conditioning or to temporal-difference learning alone.

% Auto-generated by analysis/make_table_holdout_ql.py (Entrada 91)
\begin{table}[t]
\centering
\caption{Family-level paired comparison of \textsc{IG-QL} against \textsc{IG-Random} on the external families. Differences are \textsc{IG-QL} minus \textsc{IG-Random} in GAP percentage points (negative = \textsc{IG-QL} better), with a bootstrap 95\% interval for the mean; ``Better/worse/tie'' counts instances. The rank-biserial correlation is negative when \textsc{IG-QL} is better. $p_{\mathrm{holm}}$ is Holm-adjusted over the three families of each block. The exploratory block generated the hypothesis; the confirmatory block applies a rule frozen before the held-out instances were run. Family F has no unseen instances and is therefore not testable confirmatorily.}
\label{tab:f6_holdout_ql}
{\scriptsize\setlength{\tabcolsep}{3pt}
\begin{tabular}{@{}lrcccrr@{}}
\toprule
Family & $N$ & Better/worse/tie & $\Delta$GAP [95\% CI] (pp) & $\Delta$GAP (\%) & $r_{rb}$ & $p_{\mathrm{holm}}$ \\
\midrule
\multicolumn{7}{@{}l}{\emph{Exploratory (6{,}780-instance sample)}} \\
B & 2000 & 974/741/285 & $-$0.010 [$-$0.014, $-$0.006] & $-$2.98 & $-$0.14 & $<10^{-3}$ \\
C & 2000 & 545/378/1077 & $-$0.044 [$-$0.071, $-$0.018] & $-$2.25 & $-$0.17 & $<10^{-3}$ \\
E & 2000 & 577/374/1049 & $-$0.022 [$-$0.031, $-$0.013] & $-$3.17 & $-$0.23 & $<10^{-3}$ \\
\midrule
\multicolumn{7}{@{}l}{\emph{Confirmatory (4{,}160 held-out instances)}} \\
B & 1500 & 729/553/218 & $-$0.009 [$-$0.014, $-$0.005] & $-$3.02 & $-$0.12 & $<10^{-3}$ \\
C & 1500 & 379/297/824 & $-$0.017 [$-$0.049, $+$0.018] & $-$0.85 & $-$0.11 & 0.017 \\
E & 1160 & 271/196/693 & $-$0.026 [$-$0.048, $-$0.013] & $-$3.50 & $-$0.21 & $<10^{-3}$ \\
\bottomrule
\end{tabular}}
\end{table}

\paragraph{A second confirmatory test, with an informative state}
The held-out test above used the original state. A second test, specified
before it was run, used $2{,}690$ published instances never run before (B 980;
C 1{,}050, from 75 classes of C that had not been sampled, since the earlier
classes were exhausted; E 660), with the original state (\texttt{ig\_ql}) and
with an informative one (\texttt{ig\_ql\_rel}, features relative to their value
at the start of the run, thresholds calibrated on 553 disjoint instances).
Table~\ref{tab:f16_confirm2} shows that the small effect replicates in all three
families ($-5.4\%$, $-1.8\%$, $-3.2\%$ of the GAP, i.e.\ $0.015$, $0.029$ and
$0.027$ percentage points; $p_{\mathrm{Holm}} \leq 0.003$) and that the
informative state reproduces it ($-2.9\%$, $-3.1\%$, $-3.3\%$) without adding to
the original one (pooled $-0.6\%$ of the GAP, $p=0.77$). The effect therefore does
not depend on the information in the state.

% Auto-generated by analysis/make_tables_cycle2.py (Entrada 122)
\begin{table}[t]
\centering
\caption{Second confirmatory test on 2{,}690 published instances never run before (B 980; C 1{,}050 from 75 classes not sampled earlier; E 660), with the original state (\texttt{ig\_ql}) and with an informative state (\texttt{ig\_ql\_rel}; thresholds calibrated on 553 disjoint instances). $\Delta$GAP: mean difference against \texttt{ig\_random} in percentage points (negative: Q-Learning better); Holm over six tests.}
\label{tab:f16_confirm2}
{\scriptsize\setlength{\tabcolsep}{4pt}
\begin{tabular}{@{}llrrrr@{}}
\toprule
Family & State & $N$ & W/L/T & $\Delta$GAP (pp) [\% of GAP] & $p_{\mathrm{Holm}}$ \\
\midrule
B & original & 980 & 507/348/125 & -0.0147 [-5.4\%] & $2\times10^{-7}$ \\
B & informative & 980 & 493/367/120 & -0.0079 [-2.9\%] & 0.002 \\
\midrule
C & original & 1050 & 268/182/600 & -0.0294 [-1.8\%] & 0.003 \\
C & informative & 1050 & 284/177/589 & -0.0497 [-3.1\%] & $5\times10^{-5}$ \\
\midrule
E & original & 660 & 133/88/439 & -0.0268 [-3.2\%] & 0.002 \\
E & informative & 660 & 130/78/452 & -0.0279 [-3.3\%] & $9\times10^{-4}$ \\
\bottomrule
\end{tabular}}
\end{table}




\paragraph{Gap to exact reference solutions}
The GAP elsewhere is measured against the classical bound, which can be
loose. For families C and F the benchmark provides the solutions of an exact
method (iAF; \citealp{OstojicEtAl2025}). The match of each instance file to
its reference row was verified rather than assumed: sizes agree, the
reference bound is never below ours, and no method of this study ever
obtains a makespan below the reference value (a deliberately mismatched
join fails the last check for $30$--$45\%$ of the instances).
Table~\ref{tab:f5_gap_iaf} reports the mean gap to the reference on the
$1{,}983$ C and $720$ F instances whose optimality the reference proves.
\texttt{ig\_ql} is $1.13\%$ above the optimum in C, solving $21.3\%$ to
optimality, against $1.68\%$ and $18.3\%$ for MULTIFIT; in F it is $0.17\%$
above ($33.6\%$ solved). The classical bound is loose by $0.72\%$ (C) and
$1.01\%$ (F) on average, so a sizeable part of the GAP reported elsewhere is
bound slack rather than sub-optimality (in C, about $38\%$ of the $1.90\%$
GAP of \texttt{ig\_ql}); conversely \texttt{ig\_ql}, the best method of the
study in C, is not near-optimal there: $79\%$ of its solutions exceed the
optimum. In F, \texttt{ils\_ql} reaches the optimum more often ($38.2\%$)
but errs more when it misses ($0.25\%$ against $0.17\%$).

% Auto-generated by analysis/make_table_gap_to_iaf.py (Entrada 90)
\begin{table}[t]
\centering
\caption{Mean gap to the exact reference solution (iAF; \citealp{OstojicEtAl2025}) on the external families C ($N=1983$) and F ($N=720$), restricted to instances whose optimality the reference proves. ``Opt.'' is the share of instances solved to the reference value; ``GAP$_{\mathrm{LB}}$'' is the gap to the classical lower bound used elsewhere in this paper, shown for comparison. Bold: best in column.}
\label{tab:f5_gap_iaf}
{\scriptsize\setlength{\tabcolsep}{4pt}
\begin{tabular}{@{}lcccccc@{}}
\toprule
 & \multicolumn{3}{c}{Family C} & \multicolumn{3}{c}{Family F} \\
\cmidrule(lr){2-4}\cmidrule(lr){5-7}
Method & GAP$_{\mathrm{opt}}$ (\%) & Opt.\ (\%) & GAP$_{\mathrm{LB}}$ (\%) & GAP$_{\mathrm{opt}}$ (\%) & Opt.\ (\%) & GAP$_{\mathrm{LB}}$ (\%) \\
\midrule
\textsc{LPT} & 4.57 & 8.3 & 5.39 & 0.48 & 23.8 & 1.63 \\
\textsc{MultiFit} & 1.68 & 18.3 & 2.45 & 0.65 & 33.2 & 1.80 \\
\textsc{SA (Li\&Perkins)} & 4.18 & 8.9 & 4.99 & 0.41 & 23.9 & 1.56 \\
\textsc{BCO} & 27.97 & 1.7 & 29.03 & 3.18 & 10.0 & 4.38 \\
\textsc{ILS-QL} & 3.20 & 17.9 & 4.01 & 0.25 & \textbf{38.2} & 1.40 \\
\textsc{IG-Random} & 1.18 & \textbf{21.4} & 1.95 & \textbf{0.16} & 33.1 & \textbf{1.31} \\
\textsc{IG-QL} & \textbf{1.13} & 21.3 & \textbf{1.91} & 0.17 & 33.6 & 1.32 \\
\bottomrule
\end{tabular}}
\end{table}


\subsection{Structural Limits at High $n/m$ Ratios}
\label{sec:nm-structural}

The six classes added later reach higher $n/m$ ratios (up to $15.0$,
against $3.3$--$5.0$ for the original six), which exposed a limitation of
the ILS mechanism itself, independent of the operator-selection policy.

\paragraph{At high $n/m$, LPT alone can already be near-optimal}
At $(150,15)$ ($n/m=10$) every ILS variant converged to essentially the
same solution, with zero or one accepted improving moves in $1{,}000$
iterations for most seeds, and not even $8\times$ the standard budget
changed this. With ten jobs per machine on average, LPT's guarantee is
tight and, for uniform times, the makespan concentrates near the mean
load: there is little room for any local search. Three of the six large
classes have $n/m \geq 10$ and share this property to varying degrees;
only $(150,30)$ ($n/m=5$) keeps a difficulty profile comparable to the
original six.

\paragraph{The ILS perturbation does not scale with instance size}
At $(150,30)$ the same stall recurred (three improvements of the best
makespan in $5{,}000$ iterations). The fixed perturbation of
$n\_perturb=3$ random moves is a vanishing fraction of a $150$-job
schedule, but scaling it with $n$ did not help (results were slightly
worse). A joint grid over $n\_perturb \times$ perturbation threshold, with
\texttt{ils\_random} to isolate the mechanism, showed that the threshold
was the culprit: the mean GAP improves as the threshold drops from $400$ to
about $15$ and worsens below about $10$, whereas $n\_perturb$ matters little
near that optimum. At $N=100$ on $(150,30)$ a threshold of $15$ gives a mean
GAP of $1.56\%$ against $1.69\%$ at the default $50$ (better in $23$ of
$100$ pairs, worse in $5$; $p=0.0003$), and under it \texttt{ils\_ql}
remains indistinguishable from \texttt{ils\_random} ($p=0.87$). The
threshold of $15$ is used for the ILS family on the four large classes
(Section~\ref{sec:setup}); it was not adopted as a global default, since it
was validated only at this scale.

\paragraph{Iterated Greedy does not share this ceiling}
With the corrected ILS, the four large classes were run with the full
\texttt{ig\_*} family: all five variants beat LPT in all four classes ($20$
tests, Holm-significant, $p_{\mathrm{holm}}$ from $4.9\times10^{-9}$ to
$1.2\times10^{-16}$), including \texttt{ig\_random}, which has no learning.
The ceiling is therefore a limitation of the ILS mechanism (a single
fixed-size perturbation followed by exhaustive reoptimisation around one
bottleneck machine) rather than of the instances: IG reconstructs up to
$30\%$ of the schedule at each iteration and does not stall in the same way.

\paragraph{What decides IG versus ILS: the granularity of the processing
times, with $n/m$ shifting the crossover}
An earlier version of this experiment contrasted $U[1,100]$ with
$U[1,10^8]$ and called the difference ``variance''. That label is
wrong. The coefficient of variation of a uniform distribution does not
depend on its range ($0.576$ to $0.583$ across the five ranges tested
below; Kruskal--Wallis $p=0.30$), and the main methods of this study are
exactly scale invariant (multiplying every $p_j$ by $1{,}000$
multiplies the makespan by $1{,}000$ in all six test instances, for
\texttt{ig\_ql}, \texttt{ils\_ql} and their random variants; the
raw-reward UCB baseline is the exception). What the range changes is the
number of distinct processing times, that is, the share of jobs whose
time duplicates another's: with $n=150$ it falls from $0.48$ at
$U[1,10^2]$ to $0.07$, $0.009$ and $0$ at $U[1,10^3]$, $U[1,10^4]$ and
wider ranges. Table~\ref{tab:f9_granularity} varies this range and $n/m$
orthogonally ($n=150$, $m \in \{50,15,5\}$, five ranges, $30$ instances
per cell, the same instance indices in every cell) -- something the
external families of Section~\ref{sec:external} cannot do, since range
and $n/m$ vary together there.

% Auto-generated by analysis/make_table_granularity.py (Entrada 98)
\begin{table}[t]
\centering
\caption{Granularity of the processing times versus $n/m$ ($n=150$, $p_j \sim U[1,H]$, 30 instances per cell, the same instance indices in every cell). Mean GAP (\%) of \textsc{IG-QL} / \textsc{ILS-QL}; bold marks the method that is Holm-significantly better in the cell (paired Wilcoxon, Holm over the 15 cells), no bold means a tie. ``Dup.'' is the mean share of jobs whose processing time duplicates another's; ``CV'' is the coefficient of variation of the processing times.}
\label{tab:f9_granularity}
{\scriptsize\setlength{\tabcolsep}{4pt}
\begin{tabular}{@{}lrrccc@{}}
\toprule
$H$ & Dup. & CV & $n/m=3$ & $n/m=10$ & $n/m=30$ \\
\midrule
$10^{2}$ & 0.482 & 0.576 & \textbf{1.25} / 6.68 & 0.13 / 0.16 & 0.000 / 0.000 \\
$10^{3}$ & 0.069 & 0.581 & \textbf{1.75} / 4.41 & 0.22 / \textbf{0.08} & 0.008 / \textbf{0.001} \\
$10^{4}$ & 0.009 & 0.582 & 1.80 / 1.83 & 0.22 / \textbf{0.02} & 0.010 / \textbf{0.000} \\
$10^{6}$ & 0.000 & 0.582 & 1.72 / \textbf{0.57} & 0.24 / \textbf{0.01} & 0.011 / \textbf{0.000} \\
$10^{8}$ & 0.000 & 0.583 & 1.86 / \textbf{0.54} & 0.23 / \textbf{0.01} & 0.011 / \textbf{0.000} \\
\bottomrule
\end{tabular}}
\end{table}


\emph{Granularity decides the direction.} At $n/m=3$, \texttt{ig\_ql}
wins at the two coarsest ranges ($1.25\%$ against $6.68\%$, and $1.75\%$
against $4.41\%$), the two methods tie at $U[1,10^4]$, and
\texttt{ils\_ql} wins at the two finest ranges ($0.57\%$ against
$1.72\%$, and $0.54\%$ against $1.86\%$). The pre-specified test
(Spearman correlation between $\log_{10} H$ and the GAP difference,
Holm over the three $n/m$ levels) is confirmed at every level
($\rho = 0.88$, $0.67$, $0.57$; $p_{\mathrm{holm}} < 10^{-13}$). The
same pattern appears with learning removed (\texttt{ig\_random} against
\texttt{ils\_random}), so it is a property of the two mechanisms.
\emph{The ratio $n/m$ shifts the crossover and scales the gaps.} At
$n/m=10$ and $n/m=30$ IG wins nowhere -- a tie at $U[1,10^2]$ and ILS
from $U[1,10^3]$ on -- and the absolute differences shrink by two orders
of magnitude because both methods approach the lower bound. Section
\ref{S-sec:hybrid} traces the ILS advantage on fine-grained instances
to a mechanistic cause.

\subsection{Choosing the Mechanism per Instance}
\label{sec:dispatch}

Section~\ref{sec:nm-structural} identifies the granularity of the
processing times as the instance property that decides the direction
between IG and ILS. A cheap statistic that measures it is the share of jobs whose
processing time duplicates that of another job,
$\mathrm{dup\_ratio} = 1 - |\{p_j\}|/n$: instances with many repeated
values are the coarse-grained regime in which the stochastic
reconstruction of IG suffices. The ratio alone cannot be exact, however:
at a duplication ratio of $0.07$, below the rule's threshold, IG is still
better at $n/m=3$ but ILS is better at $n/m \geq 10$
(Table~\ref{tab:f9_granularity}). The dispatch rule runs IG when
$\mathrm{dup\_ratio} > 0.1$ and ILS otherwise, with no search and no
learning in the decision. Its threshold was fixed before the
experiments reported here (a preliminary evaluation on $95$ instances
is in Supplementary Section~\ref{S-sec:hybrid}) and never retuned. The
rule was tested on the $4{,}160$ held-out instances of
Section~\ref{sec:external}, which played no part in choosing it, with
the hypotheses and decision rule written down before \texttt{ils\_ql}
was run on them: the rule must beat both fixed policies, always IG and
always ILS, in a pooled paired Wilcoxon test, Holm-corrected.

Both hypotheses are confirmed (both $p_{\mathrm{holm}} < 10^{-100}$),
but the practical picture is more modest than the $p$-values suggest
(Table~\ref{tab:f8_dispatch}). Against always using ILS the rule lowers
the mean GAP by $0.74$ percentage points; against always using IG, the
better fixed policy, by only $0.013$ points ($1.020\% \to 1.007\%$),
which closes $16\%$ of the distance to the per-instance oracle. The
gain is concentrated in family B, where ILS is the stronger mechanism
and the rule sends $92\%$ of the instances to it (accuracy $93\%$; mean
GAP $0.180\%$ against $0.299\%$ for IG); it is neutral in E and
\emph{negative} in C, where the rule sends $9\%$ of the instances to ILS
and loses $0.083$ points relative to always using IG
($p_{\mathrm{holm}} = 2.5 \times 10^{-7}$). IG is therefore a robust
default on these families, and a one-feature rule recovers the ILS
advantage where it exists at a small pooled gain and a family-specific
cost. The rule, with the same threshold, again beat both fixed policies
on $900$ new instances of six processing-time distributions (by $0.022\%$
against always-IG and $0.263\%$ against always-ILS; Supplementary
Section~\ref{S-sec:newdist}).

% Auto-generated by analysis/make_table_dispatch.py (Entrada 96)
\begin{table}[t]
\centering
\caption{Held-out test of the one-feature dispatch rule (IG if $\mathrm{dup\_ratio}>0.1$, ILS otherwise; threshold fixed before the test). ``ILS chosen'' is the share of instances sent to ILS; ``Acc.'' is the share of instances, among those where the two methods differ, on which the rule picked the better one. Mean GAP (\%) of always using \textsc{IG-QL}, always using \textsc{ILS-QL}, the rule, and the per-instance oracle (minimum of the two). Bold: best of the three non-oracle policies.}
\label{tab:f8_dispatch}
{\scriptsize\setlength{\tabcolsep}{4pt}
\begin{tabular}{@{}lrrrrrrr@{}}
\toprule
Family & $N$ & ILS chosen (\%) & Acc.\ (\%) & Always IG & Always ILS & Rule & Oracle \\
\midrule
B & 1500 & 92.1 & 93.1 & 0.299 & 0.187 & \textbf{0.180} & 0.163 \\
C & 1500 & 9.3 & 72.8 & \textbf{1.971} & 4.049 & 2.054 & 1.926 \\
E & 1160 & 58.6 & 71.0 & 0.724 & 0.793 & \textbf{0.722} & 0.659 \\
\midrule
All & 4160 & 52.9 & 81.7 & 1.020 & 1.748 & \textbf{1.007} & 0.937 \\
\bottomrule
\end{tabular}}
\end{table}


\subsection{Another Problem: Unrelated Machines}
\label{sec:unrelated}

Everything above concerns identical machines, where exact methods are strong
and the decisions of the search saturate quickly. To see whether the
decomposition is a property of the problem or of the method, the protocol was
repeated on the unrelated-machine problem $\mathrm{R}||C_{\max}$, in which the
processing time $p_{ij}$ depends on the machine, a classical hard variant for
which Iterated Greedy methods with local search are well established
\citep{FanjulPeyroRuiz2010} and whose assignment MILP does not scale gracefully
\citep[LP rounding gives a $2$-approximation;][]{LenstraEtal1990}.

\paragraph{Methods and design}
The same agent, state design and decision rules were reused. The constructive
heuristic assigns jobs, in decreasing order of $\min_i p_{ij}$, to the machine
with the earliest completion. The IG removes $d$ jobs with one of four
strategies (random; largest processing time on the current machine; random
machine; critical machine) and reinserts them greedily, with or without a
local search (a descent that moves or swaps jobs from the critical machine,
at most 20 improvements per iteration); the ILS applies a best-improvement move,
swap or double move from the critical machine and perturbs three random jobs
after 50 stagnant iterations. Instances were generated with four distributions
($U[1,100]$, $U[10,100]$, $U[100,200]$, $U[100,120]$) and five sizes
($50\times10$, $100\times10$, $100\times20$, $200\times20$, $500\times30$), 20
instances per cell ($400$ in all), $T=1{,}000$. Following the lesson of
Section~\ref{sec:f3f4-ablation}, the fraction $d/n$ was not fixed: it was
chosen per size on separate instances (0.30, 0.20, 0.10, 0.10, 0.05), and the
thresholds of the informative state were calibrated on further separate
instances. References: the LP relaxation of the assignment model and a MILP
(HiGHS) with 5 and 60 seconds, on 160 of the instances. All hypotheses were
specified beforehand (Supplementary Section~\ref{S-sec:prereg-record}).

\paragraph{The mechanism matters, as before}
Table~\ref{tab:f18_unrelated} shows the same ordering as in
Section~\ref{sec:f3f4-ablation}, with a different cause. The IG with local search
beats the ILS by $4.16\%$ (384 of 400 instances), also at equal CPU time ($2.16\%$,
the ILS receiving 8{,}000 iterations, $1.25\times$ the CPU; Table~\ref{tab:f23_unrelated_budget}).
But the \emph{pure} IG, without local search, does not beat the ILS ($+1.39\%$;
it is better only in $U[1,100]$), and local search improves it by $5.40\%$: for
unrelated machines the active ingredient is the combination of destruction--reconstruction
with a descent, whereas for identical machines destruction--reconstruction
alone was enough. Both heuristics improve greatly on the constructive start
($-13.1\%$ and $-9.3\%$).

\paragraph{The policy matters more, but still little, and the ILS policy hurts}
Q-Learning over the destruction strategies of the IG with local search gains
$0.145\%$ over random selection (90\% interval $0.054$--$0.234$; $p_{\mathrm{Holm}}=0.011$), about thirteen times the
identical-machine value, yet about thirty times smaller than the mechanism
effect; its interval does not lie inside the $\pm0.10\%$ margin used for
identical machines.
In the ILS, Q-Learning is \emph{worse} than random operator selection ($+0.96\%$, $p_{\mathrm{Holm}}<10^{-14}$),
as for identical machines with an informative state, and remains worse with
eight times the budget ($+0.92\%$). The headroom explains the difference in the IG: the
omniscient per-instance choice among the four strategies gains $0.38\%$ (up
to $1.12\%$ in the largest $U[1,100]$ cell) against $0.04\%$ for identical machines, and
Table~\ref{tab:f19_unrelated_cells} shows that the policy gain is large where
the headroom is large and null in $U[100,120]$, where the machines are almost
identical and the headroom vanishes. Pooled over cells, Q-Learning captures
$38\%$ of the headroom.

\paragraph{Exact references: the metaheuristic does not win here}
The MILP proves optimality in $66.9\%$ of the 160 instances within 60 seconds
($31.9\%$ within five). The IG with local search and Q-Learning, which takes a few
seconds, ends $2.04\%$ above the best lower bound (MILP: $0.43\%$) and is no
worse than the MILP incumbent in $28.1\%$ of the instances; test M4 is not
confirmed in its favour ($+1.59\%$). Five times the budget reduces the
difference to $+0.95\%$ and raises the share to $42.5\%$. The IG is
competitive only in the largest and narrowest cells ($U[100,120]$, $500\times30$:
$100\%$ of the instances). Unlike for identical machines, then, the heuristics
studied here are not a practical alternative to the exact method at these sizes;
the purpose of this section is the attribution, not the solver.

\paragraph{Larger instances do not reverse it}
A pilot on a single $1000\times50$ instance suggested that the MILP might fail at
that size, so the question was tested with a pre-specified design on $32$ new
instances ($n=1000$, $m=50$, the four distributions, eight instances each; $d/n$
chosen on $24$ separate instances; Table~\ref{tab:f27_unrelated_large}). It does
not fail: within 60 seconds the MILP returned an incumbent in all $32$ instances
and ended $2.04\%$ above the best lower bound, against $4.15\%$ for the IG with
local search and Q-Learning at $T=1{,}000$ (makespan $2.24\%$ higher, MILP better in
$28$ of $32$ instances, $p_{\mathrm{Holm}}=10^{-4}$) and $2.78\%$ at $T=5{,}000$
($+0.89\%$, interval $-0.46$ to $+1.95$, with $37$ seconds of CPU time against $51$).
The only distribution in which the heuristic is better on average is $U[10,100]$,
where the MILP is erratic (gaps from $1.3\%$ to $28.1\%$ among the eight instances)
and yet better than the $T=1{,}000$ heuristic in seven of the eight pairs. The
conclusion stands for these heuristics. Started from the linear relaxation (the LP-IG described below), the IG is no longer dominated: its mean GAP to the best lower bound is $1.09\%$ against $2.04\%$ for the MILP, in $6.5$ seconds against $51$, and it equals or beats the MILP in $62.5\%$ of the instances (against $28\%$ without the LP start). A second pre-specified test does not confirm the superiority, however ($-0.76\%$, interval $-2.10$ to $+0.13$, $p_{\mathrm{Holm}}=0.72$): the MILP is better in $U[1,100]$ ($1.23\%$ against $2.46\%$) and the LP-IG in the other three distributions, most of all in $U[10,100]$, where the MILP is erratic ($6.25\%$ against $1.59\%$); the search does not improve beyond $T=1000$ at this size.

% Auto-generated by analysis/make_tables_cycle5.py (Entrada 129)
\begin{table}[t]
\centering
\caption{$\mathrm{R}||C_{\max}$ at $n=1000$, $m=50$ (32 new instances; four distributions $\times$ eight replications; $d/n$ chosen on 24 separate instances): the iterated greedy with local search against a MILP with a 60\,s limit. $\bar d$ in \% of the makespan (negative: IG+LS better); Wilcoxon, Holm over two, on the instances in which the MILP returned an incumbent.}
\label{tab:f27_unrelated_large}
{\scriptsize\setlength{\tabcolsep}{3.5pt}
\begin{tabular}{@{}lrrrrr@{}}
\toprule
Comparison & $N$ & W/L/T & $\bar d$ (\%) & 90\% interval & $p_{\mathrm{Holm}}$ \\
\midrule
G1: IG+LS \textsc{QL} ($T=1000$) vs.\ MILP (60\,s) & 32 & 3/28/1 & +2.241 & [+0.664, +3.573] & $1\times10^{-4}$ \\
G2: IG+LS \textsc{QL} ($T=5000$) vs.\ MILP (60\,s) & 32 & 8/23/1 & +0.887 & [-0.455, +1.952] & 0.023 \\
\midrule
\multicolumn{6}{l}{Second pre-specified test: LP-IG (Section on the standard benchmark), $d/n=0.02$; Holm over two} \\
M1: LP-IG ($T=1000$) vs.\ MILP (60\,s) & 32 & 15/12/5 & -0.758 & [-2.102, +0.128] & 0.722 \\
M2: LP-IG ($T=5000$) vs.\ MILP (60\,s) & 32 & 15/12/5 & -0.758 & [-2.102, +0.128] & 0.722 \\
\bottomrule
\end{tabular}}
\end{table}


\paragraph{The policy at this size: large where one strategy dominates, harmful elsewhere}
The iteration budget is far from saturating the search at $1000\times50$ (the GAP falls
from $4.15\%$ to $2.78\%$ with five times the budget), and on the $32$ instances above the
policy gain over random selection was $0.52\%$, five times that at $n\leq500$. Whether it
survives, and at equal CPU time (the policy costs $1.39$ times the CPU time of random
selection here), was tested on $96$ further instances with a pre-specified design
(Table~\ref{tab:f29_policy_large}). Pooled over the four distributions the gain is
$0.35\%$ at an equal $T$ ($p_{\mathrm{Holm}}=0.19$) and $0.19\%$ at equal CPU time
($p_{\mathrm{Holm}}=0.67$): neither is confirmed, and the prediction recorded before the
run (that the first would be) was wrong. The reason is that the effect is not uniform. In a
post hoc breakdown (Holm over the four distributions), in $U[1,100]$ one strategy,
random destruction, is by itself $1.66\%$ better than the uniform mixture of the four, and
Q-Learning finds it and gains $1.73\%$ ($17$ wins, no losses; $p_{\mathrm{Holm}}=8\times10^{-4}$),
capturing all of the headroom of the fixed strategies. In $U[100,120]$ and $U[100,200]$,
by contrast, every fixed strategy is worse than the mixture ($0.9$--$4.5\%$), the policy
is worse than random selection ($+0.10\%$ and $+0.38\%$) and, since it costs up to $1.63$
times the CPU time (the state), at equal CPU time random selection wins in $23$ of $24$
instances in both ($+0.11\%$ and $+0.89\%$); $U[10,100]$ is in between. The policy gain thus
follows the headroom of the fixed strategies once more, with a case in which it is
captured completely, but ``the policy pays where the search is not saturated'' is not a
general rule.

% Auto-generated by analysis/make_tables_cycle5.py (Entrada 129)
\begin{table}[t]
\centering
\caption{Does the policy pay at $1000\times50$ ($\mathrm{R}||C_{\max}$, 96 new instances, 24 per distribution; IG with local search, $d/n=0.02$, $T=1000$)? Pre-specified pooled tests Q1 (\textsc{QL} vs.\ random selection, equal $T$) and Q2 (\textsc{QL} vs.\ random selection with the same CPU time, $T_{\mathrm{eq}}=1000$, $1164$, $1734$, $1436$ for the four distributions), Holm over two. Below the line, a post hoc breakdown by distribution (Holm over four within each column): $\bar d$ in \% of the makespan (negative: \textsc{QL} better); the last column is the best single destruction strategy (fixed) against the uniform mixture.}
\label{tab:f29_policy_large}
{\scriptsize\setlength{\tabcolsep}{3.5pt}
\begin{tabular}{@{}lrrrrr@{}}
\toprule
Comparison & W/L/T & $\bar d$ (\%) & 90\% interval & $p_{\mathrm{Holm}}$ & Best fixed (\%) \\
\midrule
Q1: \textsc{QL} vs.\ random, equal $T$ (pooled) & 34/43/19 & -0.353 & [-0.556, -0.151] & 0.193 & -- \\
Q2: \textsc{QL} vs.\ random, equal CPU (pooled) & 23/51/22 & -0.190 & [-0.402, +0.018] & 0.666 & -- \\
\midrule
\multicolumn{6}{l}{Post hoc, by distribution: \textsc{QL} vs.\ random at equal $T$ / at equal CPU} \\
$U[1,100]$, equal $T$ & 17/0/7 & -1.729 & -- & $8\times10^{-4}$ & -1.657 \\
\quad equal CPU & 17/0/7 & -1.729 & -- & $6\times10^{-4}$ & \\
$U[10,100]$, equal $T$ & 10/4/10 & -0.156 & -- & 0.217 & -0.034 \\
\quad equal CPU & 6/4/14 & -0.034 & -- & 1.000 & \\
$U[100,120]$, equal $T$ & 0/23/1 & +0.098 & -- & $1\times10^{-4}$ & +1.691 \\
\quad equal CPU & 0/24/0 & +0.110 & -- & $5\times10^{-5}$ & \\
$U[100,200]$, equal $T$ & 7/16/1 & +0.376 & -- & 0.217 & +0.923 \\
\quad equal CPU & 0/23/1 & +0.893 & -- & $1\times10^{-6}$ & \\
\bottomrule
\end{tabular}}
\end{table}


Whether offline validation does the same as the policy for this choice, as it does for the
destruction size on identical machines, was tested on $96$ further instances
(Table~\ref{tab:f30_unrelated_valid}). One of five candidates (each strategy fixed, or the
uniform mixture) was chosen per distribution on $24$ separate instances (random
destruction in $U[1,100]$, the mixture in the other three), and compared with Q-Learning
on new instances. The validated choice captures the headroom ($-0.46\%$ against random
selection, $17$ wins and no losses, $p_{\mathrm{Holm}}=8\times10^{-4}$) at a lower CPU
time ($3.3$ against $4.7$ seconds), and the comparison with Q-Learning is inconclusive
($-0.14\%$, interval $-0.34$ to $+0.07$, $p_{\mathrm{Holm}}=0.30$): Q-Learning beats the
best fixed strategy by $0.45\%$ in $U[1,100]$ and loses by $1.1\%$ to the mixture in
$U[100,200]$. The pooled gain of Q-Learning over random selection is again not
confirmed ($0.31\%$, $p_{\mathrm{Holm}}=0.30$), with the same sign pattern by distribution
as in the first sample. At this size, as at the others, the pooled data show no advantage
of learning online over choosing offline, although in one distribution the agent does
better than the best fixed strategy.

% Auto-generated by analysis/make_tables_cycle5.py (Entrada 129)
\begin{table}[t]
\centering
\caption{Offline validation against the policy at $1000\times50$ ($\mathrm{R}||C_{\max}$, 96 new instances, 24 per distribution, $T=1000$). The candidate (one of the four destruction strategies fixed, or the uniform mixture) was chosen per distribution on 24 separate instances (fixed strategy 0, random destruction, in $U[1,100]$; the mixture in the other three). $\bar d$ in \% of the makespan (negative: first better); Wilcoxon, Holm over three. Mean CPU time per run: validated choice 3.3\,s, \textsc{QL} 4.7\,s, random selection 3.3\,s.}
\label{tab:f30_unrelated_valid}
{\scriptsize\setlength{\tabcolsep}{3.5pt}
\begin{tabular}{@{}lrrrr@{}}
\toprule
Comparison & W/L/T & $\bar d$ (\%) & 90\% interval & $p_{\mathrm{Holm}}$ \\
\midrule
V1: validated choice vs.\ \textsc{QL} & 47/23/26 & -0.135 & [-0.337, +0.070] & 0.296 \\
V2: validated choice vs.\ random selection & 17/0/79 & -0.461 & [-0.648, -0.282] & $8\times10^{-4}$ \\
V3: \textsc{QL} vs.\ random selection & 38/45/13 & -0.312 & [-0.586, -0.044] & 0.296 \\
\bottomrule
\end{tabular}}
\end{table}


\paragraph{External validation on the standard benchmark}
All the instances above are synthetic, so the section was repeated on the standard
$\mathrm{R}||C_{\max}$ benchmark of \citet{FanjulPeyroRuiz2010}, in the copy distributed with
\citet{AbladEtAl2021}: $1{,}400$ instances (seven types, $m\in\{10,\ldots,50\}$,
$n\in\{100,\ldots,1000\}$, ten instances per cell) with, for each, the best solution known
from CPLEX (two hours), Gurobi and an exact method (300 seconds each), proved optimal in
$73\%$ of the test instances. Two instances per cell ($280$) were used for tuning ($d/n$ per $n$;
the strategy per type among the four fixed ones and the mixture; the CPU ratios) and five
($700$) for a pre-specified test (Table~\ref{tab:f31_rcmax_benchmark}); the rest were not used.
The decomposition replicates. The IG with local search beats the ILS at equal CPU time by
$2.14\%$ ($529$ wins against $99$; $3.23\%$ at equal $T$), the policy gain over random
selection is $0.055\%$ at equal $T$ and $0.033\%$ at equal CPU time (neither confirmed,
$p_{\mathrm{Holm}}=0.15$ and $0.48$), the validated strategy beats Q-Learning by $0.066\%$
($p_{\mathrm{Holm}}=0.007$), and the headroom law is confirmed on third-party data: the
Spearman correlation between the headroom of a cell and the policy gain over the $28$ cells
is $0.74$ ($p_{\mathrm{Holm}}=10^{-5}$). The gain is concentrated where a fixed strategy
dominates: $0.61\%$ in $U[1,100]$ ($46$ wins, $9$ losses) and $0.15\%$ in $U[10,100]$, against
$+0.20\%$ (worse) in machine-correlated instances. The heuristics are not competitive with the exact references: the IG with
local search and Q-Learning ends $2.08\%$ above the best known solution on average ($0.13\%$ in $U[1000,1100]$
to $5.64\%$ in $U[1,100]$), against $0.07\%$ for Gurobi and $0.09\%$ for the exact method
at 300 seconds, and equals or beats it in $10.9\%$ of the instances ($0.1\%$ strictly better).

% Auto-generated by analysis/make_tables_cycle5.py (Entrada 129)
\begin{table}[t]
\centering
\caption{External validation on the standard $\mathrm{R}||C_{\max}$ benchmark of \citet{FanjulPeyroRuiz2010} (700 test instances: 7 types $\times$ 5 machine counts $\times$ 4 job counts $\times$ 5 instances; tuning on 280 other instances; $T=1000$). $\bar d$ in \% of the makespan (negative: first better); Wilcoxon, Holm over five. B3 is the Spearman correlation between the headroom of the destruction strategies and the policy gain over the 28 cells (type $\times$ $n$), one-sided. Below the line, the mean GAP (\%) to the best known solution (the minimum of CPLEX in 2 h, Gurobi and the LR method in 300 s; \citealt{AbladEtAl2021}), by instance type.}
\label{tab:f31_rcmax_benchmark}
{\scriptsize\setlength{\tabcolsep}{3.5pt}
\begin{tabular}{@{}lrrrr@{}}
\toprule
Comparison & W/L/T & $\bar d$ (\%) or $\rho$ & 90\% interval & $p_{\mathrm{Holm}}$ \\
\midrule
B1: IG+LS vs.\ ILS at equal CPU (mechanism) & 529/99/72 & -2.135 & [-2.347, -1.926] & $6\times10^{-68}$ \\
B2a: \textsc{QL} vs.\ random selection, equal $T$ & 229/213/258 & -0.055 & [-0.106, -0.004] & 0.150 \\
B2b: \textsc{QL} vs.\ random selection, equal CPU & 218/222/260 & -0.033 & [-0.084, +0.016] & 0.483 \\
B3: Spearman(headroom, policy gain), 28 cells & -- & +0.74 & -- & $1\times10^{-5}$ \\
B4: validated strategy (per type) vs.\ \textsc{QL} & 249/182/269 & -0.066 & [-0.120, -0.015] & 0.007 \\
\midrule
\multicolumn{5}{l}{Mean GAP (\%) to the best known solution: IG+LS \textsc{QL} / ILS at equal CPU (type)} \\
\quad job-correlated & \multicolumn{4}{l}{1.79 / 1.84} \\
\quad machine-correlated & \multicolumn{4}{l}{1.87 / 6.60} \\
\quad $U[1000,1100]$ & \multicolumn{4}{l}{0.13 / 0.25} \\
\quad $U[100,120]$ & \multicolumn{4}{l}{0.33 / 0.93} \\
\quad $U[100,200]$ & \multicolumn{4}{l}{0.66 / 2.26} \\
\quad $U[10,100]$ & \multicolumn{4}{l}{4.17 / 7.62} \\
\quad $U[1,100]$ & \multicolumn{4}{l}{5.64 / 12.34} \\
\bottomrule
\end{tabular}}
\end{table}


\paragraph{A stronger heuristic: starting from the linear relaxation}
The distance to the exact references suggested improving the heuristic rather than the policy. The IG with local search above starts from a constructive solution. The linear relaxation of the assignment model (solved with HiGHS in $0.4$ seconds on average, descent included) has a vertex solution with fewer fractional jobs than machines (at most $49$ of $100$ jobs in the smallest cells, $29$ of $1{,}000$ in the largest); rounding it (the integral jobs stay, the fractional ones go to the machine of earliest completion within their support) and descending gives a start that is poor by itself ($5.2\%$ above the best known solution on the tuning instances) but from which the same IG is far better (LP-IG; Table~\ref{tab:f32_lpig}). On the $420$ instances of the benchmark that had never been used, with $d/n$ tuned on $280$ others and a pre-specified design, LP-IG beats the IG started from the constructive solution by $1.18\%$ at equal CPU time ($308$ wins against $21$), and the mean GAP to the best known solutions falls from $2.23\%$ to $0.97\%$ ($0.31\%$ for $n=1000$; $T=1000$, $3.3$ seconds per instance). With $T=5000$ it is $0.85\%$ ($15.7$ seconds), equals or beats the best known solution of CPLEX, Gurobi and the exact method in $27\%$ of the instances and improves on it in $3.8\%$. It meets the competitiveness criterion fixed beforehand ($\leq1.0\%$ overall and $\leq0.5\%$ for $n=1000$), but the exact references remain ahead ($0.09$--$0.10\%$ on average, with 300 seconds), and the heuristic is weak where the relaxation is loose or the instances are small: $1.2$--$1.7\%$ for $n\leq200$ and $2$--$3\%$ in $U[1,100]$ and $U[10,100]$. Q-Learning does not help it ($+0.055\%$, not confirmed).

% Auto-generated by analysis/make_tables_cycle5.py (Entrada 129)
\begin{table}[t]
\centering
\caption{A stronger heuristic for $\mathrm{R}||C_{\max}$: Iterated Greedy with local search started from the rounded linear relaxation of the assignment model (LP-IG), on the 420 instances of the standard benchmark that were never used before (the $d/n$ was tuned on 280 other instances). $\bar d$ in \% of the makespan (negative: first better); Wilcoxon, Holm over three. Below the line, the mean GAP (\%) to the best known solution (the minimum of CPLEX, Gurobi and LR; Table~\ref{tab:f31_rcmax_benchmark}): the baseline IG, LP-IG ($T=1000$ and $5000$) and the exact references (Gurobi, LR in 300\,s; CPLEX in 2\,h), by job count $n$.}
\label{tab:f32_lpig}
{\scriptsize\setlength{\tabcolsep}{3.5pt}
\begin{tabular}{@{}lrrrr@{}}
\toprule
Comparison & W/L/T & $\bar d$ (\%) & 90\% interval & $p_{\mathrm{Holm}}$ \\
\midrule
L1: LP-IG vs.\ IG from the constructive start, equal CPU & 308/21/91 & -1.176 & [-1.338, -1.018] & $4\times10^{-46}$ \\
L2: LP-IG $T=5000$ vs.\ $T=1000$ & 104/0/316 & -0.124 & [-0.163, -0.088] & $2\times10^{-18}$ \\
L3: LP-IG with \textsc{QL} vs.\ random selection & 107/67/246 & +0.055 & [-0.033, +0.149] & 0.122 \\
\midrule
\multicolumn{5}{l}{Mean GAP (\%) to the best known solution: baseline IG / LP-IG ($T=1000$) / LP-IG ($T=5000$) / Gurobi / LR} \\
\quad $n=100$ & \multicolumn{4}{l}{1.85 / 1.21 / 0.92 / 0.26 / 0.00} \\
\quad $n=200$ & \multicolumn{4}{l}{2.48 / 1.74 / 1.53 / 0.09 / 0.01} \\
\quad $n=500$ & \multicolumn{4}{l}{2.41 / 0.63 / 0.62 / 0.04 / 0.20} \\
\quad $n=1000$ & \multicolumn{4}{l}{2.18 / 0.31 / 0.31 / 0.02 / 0.15} \\
\quad all instances & \multicolumn{4}{l}{2.23 / 0.97 / 0.85 / 0.10 / 0.09} \\
\bottomrule
\end{tabular}}
\end{table}


\paragraph{A prospective test of the headroom law: the destruction size}
The headroom law was derived from, and tested on, destruction strategies. To test it prospectively in an action space with headroom, the LP-IG was run with the destruction size $d/n\in\{0.01,0.02,0.05,0.10,0.20\}$ as the action (the $700$ benchmark instances of the test above; design fixed beforehand; Table~\ref{tab:f33_dpolicy}), against random choice of $d$ at each iteration, the choice of $d$ validated on $10$ separate instances per type and job count, an online Q-Learning agent, a UCB1 bandit and a Q-Learning table pre-trained offline on those same $10$ instances. The headroom is real but modest: validation improves on random choice by $0.065\%$ per $n$ and $0.201\%$ per type and $n$ (a per-instance oracle: $0.325\%$). The policy does not capture it. Online Q-Learning is \emph{worse} than random choice ($+0.152\%$, $p_{\mathrm{Holm}}=10^{-5}$) and than the validated choice ($+0.360\%$), the pre-trained table also loses to validation ($+0.249\%$), and over the $28$ cells of type and $n$ the policy gain is unrelated to the headroom (Spearman $-0.37$; the criterion fixed beforehand was $\geq0.5$): the largest loss ($+0.44\%$) is in $U[1,100]$, where the headroom is largest ($1.05\%$). Most of the predictions recorded before the run were wrong. The law is therefore refuted as a general statement: headroom is necessary but not sufficient. In an exploratory diagnostic on $21$ instances, the immediate reward ranks the fixed sizes consistently with their final quality (rank correlation $0.6$--$0.7$ early in the run), but more than $90\%$ of it arrives in the first $100$ of the $1{,}000$ iterations (mean reward of the best size: $0.21$, then $0.011$ and $0.002$ in the next windows), leaving the agent a very short window, with high exploration, in which to learn. In a pilot on $126$ tuning instances, bandits designed for such a short window (Thompson sampling, with and without discounting, and a successive-elimination race) did not beat random choice of $d$ either ($+0.14\%$, $+0.07\%$ and $+0.06\%$; Q-Learning $+0.05\%$; the validated size $-0.18\%$, optimistic), so none of five learners captures the headroom. Choosing $d$ is, in these data, a configuration problem solved offline, not a control problem solved online.

% Auto-generated by analysis/make_tables_cycle5.py (Entrada 129)
\begin{table}[t]
\centering
\caption{A prospective test of the headroom law with the destruction size as the action (LP-IG, 700 instances of the standard $\mathrm{R}||C_{\max}$ benchmark, $T=1000$; actions $d/n\in\{0.01,0.02,0.05,0.10,0.20\}$). Validated choices and the pre-trained tables use 10 tuning instances per type and $n$. $\bar d$ in \% of the makespan (negative: first better); Wilcoxon, Holm over five. D5 is the Spearman correlation between the headroom of a cell (type $\times$ $n$; 28 cells) and the policy gain, one-sided, with the pre-specified criterion $\rho\geq0.5$. Last row: mean gain over random $d/n$ (\%) of each choice.}
\label{tab:f33_dpolicy}
{\scriptsize\setlength{\tabcolsep}{3.5pt}
\begin{tabular}{@{}lrrrr@{}}
\toprule
Comparison & W/L/T & $\bar d$ (\%) or $\rho$ & 90\% interval & $p_{\mathrm{Holm}}$ \\
\midrule
D1: validated $d/n$ (per $n$) vs.\ random $d/n$ & 189/181/330 & -0.065 & [-0.126, +0.004] & 0.006 \\
D2: \textsc{QL} over $d/n$ vs.\ random $d/n$ & 144/221/335 & +0.152 & [+0.080, +0.223] & $1\times10^{-5}$ \\
D3: \textsc{QL} vs.\ validated $d/n$ (per type, $n$) & 85/298/317 & +0.360 & [+0.285, +0.443] & $2\times10^{-30}$ \\
D4: pre-trained \textsc{QL} vs.\ validated $d/n$ (per type, $n$) & 96/234/370 & +0.249 & [+0.184, +0.317] & $2\times10^{-16}$ \\
D5: Spearman(headroom, \textsc{QL} gain), 28 cells & -- & -0.37 & -- & 0.975 \\
\midrule
\multicolumn{5}{p{0.96\textwidth}}{Gain over random $d/n$ (\%): validated per $n$ -0.065; validated per type and $n$ -0.201; per-instance oracle -0.325; \textsc{QL} +0.152; UCB1 -0.018; pre-trained \textsc{QL} +0.040} \\
\bottomrule
\end{tabular}}
\end{table}


\paragraph{Closing the distance to the exact references: a restricted MILP}
The LP-IG still ends $0.97\%$ above the best known solutions, against $0.1\%$ for Gurobi and the exact method, and its worst cells are the small instances and the wide-range processing times ($1.2$--$1.7\%$ for $n\leq200$; $2$--$3\%$ in $U[1,100]$ and $U[10,100]$). Following the size-reduction idea of the literature on unrelated machines, a MILP in which each job may only go to its six cheapest machines (and to the machine it occupies in the incumbent, so that the restricted problem is feasible and contains the LP-IG solution) was solved with HiGHS for at most $30$ seconds after the LP-IG. The design (tuning instances only: $k=3$ against $k=6$) is described in the diary; the test was run on the $420$ instances of the benchmark that had never been used (Table~\ref{tab:f34_lpigsr}). It never worsens the solution and improves $54.8\%$ of the LP-IG solutions; against more iterations of the LP-IG at a similar CPU time ($T=8000$, $27.3$ against $23.3$ seconds) it is better by $0.49\%$ ($196$ wins against $33$, $p_{\mathrm{Holm}}=10^{-24}$), and the mean GAP to the best known solutions falls from $0.97\%$ (LP-IG) and $0.78\%$ ($T=8000$) to $0.25\%$ ($0.11\%$ for $n=1000$; baseline IG: $2.23\%$). That is within the range of CPLEX with two hours ($0.20\%$) and $0.15$ points from Gurobi and the exact method with $300$ seconds ($0.10$ and $0.09\%$); it equals or beats the best known solution in $50\%$ of the instances and improves on it in $4.8\%$. The competitiveness criterion fixed beforehand (GAP $\leq0.3\%$ overall and $\leq0.5\%$ in every job count and every type) is met overall and for every $n$ (at most $0.36\%$) but not for the job-correlated instances ($0.67\%$). No learned component takes part in any of the three steps: the gains come from the linear relaxation ($-1.18\%$ at equal CPU) and the restricted MILP ($-0.49\%$ at equal CPU). A variant designed on the tuning instances, which chooses the six candidate machines of each job by the reduced cost of the LP instead of by the processing time, was tested on the $700$ remaining instances: its advantage over the cheapest-machine version is significant in sign but negligible ($-0.03\%$, $p_{\mathrm{Holm}}=7\times10^{-6}$; mean GAP $0.30\%$ to $0.27\%$), large only for machine-correlated instances ($0.31\%$ to $0.08\%$) and worse for $U[1,100]$ ($0.40\%$ to $0.50\%$); the criterion of a GAP of at most $0.25\%$ overall was not met ($0.27\%$), and neither was $0.5\%$ in every type.

% Auto-generated by analysis/make_tables_cycle5.py (Entrada 129)
\begin{table}[t]
\centering
\caption{LP-IG followed by a MILP restricted to the six cheapest machines of each job (and to the machine of the incumbent), 30\,s limit, on the 420 reserved instances of the standard $\mathrm{R}||C_{\max}$ benchmark (design on 168 tuning instances; $T=1000$). $\bar d$ in \% of the makespan (negative: first better); Wilcoxon, Holm over three. Below the line, the mean GAP (\%) to the best known solution by $n$: baseline IG / LP-IG ($T=1000$) / LP-IG ($T=8000$) / LP-IG + restricted MILP / Gurobi / LR (CPLEX in 2\,h: 0.20). Mean CPU time per run: LP-IG 3.3\,s; LP-IG ($T=8000$) 27.3\,s; + restricted MILP 23.3\,s.}
\label{tab:f34_lpigsr}
{\scriptsize\setlength{\tabcolsep}{3.5pt}
\begin{tabular}{@{}lrrrr@{}}
\toprule
Comparison & W/L/T & $\bar d$ (\%) & 90\% interval & $p_{\mathrm{Holm}}$ \\
\midrule
S1: + restricted MILP vs.\ LP-IG with $T=8000$ (similar CPU) & 196/33/191 & -0.489 & [-0.628, -0.366] & $6\times10^{-25}$ \\
S2: + restricted MILP vs.\ LP-IG ($T=1000$) & 230/0/190 & -0.675 & [-0.832, -0.540] & $5\times10^{-39}$ \\
S3: + restricted MILP vs.\ LP-IG ($T=5000$) & 204/28/188 & -0.552 & [-0.709, -0.420] & $2\times10^{-27}$ \\
\midrule
\multicolumn{5}{l}{Mean GAP (\%) to the best known solution} \\
\quad $n=100$ & \multicolumn{4}{l}{1.85 / 1.21 / 0.77 / 0.36 / 0.26 / 0.00} \\
\quad $n=200$ & \multicolumn{4}{l}{2.48 / 1.74 / 1.43 / 0.27 / 0.09 / 0.01} \\
\quad $n=500$ & \multicolumn{4}{l}{2.41 / 0.63 / 0.59 / 0.28 / 0.04 / 0.20} \\
\quad $n=1000$ & \multicolumn{4}{l}{2.18 / 0.31 / 0.31 / 0.11 / 0.02 / 0.15} \\
\quad all instances & \multicolumn{4}{l}{2.23 / 0.97 / 0.78 / 0.25 / 0.10 / 0.09} \\
\bottomrule
\end{tabular}}
\end{table}


% Auto-generated by analysis/make_tables_cycle2.py (Entrada 122)
\begin{table}[t]
\centering
\caption{$\mathrm{R}||C_{\max}$ (unrelated machines): the mechanism-versus-policy protocol on 400 instances (four processing-time distributions $\times$ five sizes $\times$ 20), $T=1000$, destruction size chosen on separate instances. $\bar d$: mean relative makespan difference (\%; negative: first better); Wilcoxon, Holm over four. M4 uses the 160 instances with a MILP reference.}
\label{tab:f18_unrelated}
{\scriptsize\setlength{\tabcolsep}{3.5pt}
\begin{tabular}{@{}lrrrrr@{}}
\toprule
Comparison & $N$ & W/L/T & $\bar d$ (\%) & 90\% interval & $p_{\mathrm{Holm}}$ \\
\midrule
M1: IG+LS vs.\ ILS (random selection) & 400 & 384/1/15 & -4.160 & [-4.480, -3.839] & $3\times10^{-64}$ \\
M2: IG+LS, \textsc{QL} vs.\ random selection & 400 & 111/108/181 & -0.145 & [-0.234, -0.054] & 0.011 \\
M3: ILS, \textsc{QL} vs.\ random selection & 400 & 78/222/100 & +0.956 & [+0.730, +1.191] & $3\times10^{-15}$ \\
M4: IG+LS-\textsc{QL} vs.\ MILP (60\,s) & 160 & 14/115/31 & +1.592 & [+1.310, +1.880] & $5\times10^{-18}$ \\
\bottomrule
\end{tabular}}
\end{table}

% Auto-generated by analysis/make_tables_cycle2.py (Entrada 122)
\begin{table}[t]
\centering
\caption{$\mathrm{R}||C_{\max}$ by cell (20 instances each): destruction strategies' headroom against random selection (best fixed strategy chosen in-sample; omniscient per-instance choice), effect of the Q-Learning policy, and IG+LS versus ILS; all in \% of the makespan (negative: first better).}
\label{tab:f19_unrelated_cells}
{\scriptsize\setlength{\tabcolsep}{3.5pt}
\begin{tabular}{@{}llrrrrr@{}}
\toprule
Distribution & Size & Best fixed & Oracle & \textsc{QL} policy & ILS vs.\ IG+LS & GAP IG+LS-\textsc{QL} \\
\midrule
$U[1,100]$ & 50x10 & -0.14 & -0.41 & -0.38 & +8.2 & 16.66 \\
$U[1,100]$ & 100x10 & -0.61 & -1.01 & -0.42 & +6.2 & 8.16 \\
$U[1,100]$ & 100x20 & -0.67 & -0.85 & +0.25 & +8.6 & 20.09 \\
$U[1,100]$ & 200x20 & -0.63 & -0.96 & -0.55 & +9.8 & 10.59 \\
$U[1,100]$ & 500x30 & -1.12 & -1.12 & -0.61 & +9.6 & 10.03 \\
\midrule
$U[10,100]$ & 50x10 & -0.10 & -0.40 & +0.10 & +5.5 & 9.56 \\
$U[10,100]$ & 100x10 & -0.49 & -0.55 & -0.35 & +3.3 & 5.48 \\
$U[10,100]$ & 100x20 & -0.17 & -0.42 & -0.10 & +8.4 & 11.07 \\
$U[10,100]$ & 200x20 & -0.86 & -0.86 & -0.30 & +5.1 & 7.27 \\
$U[10,100]$ & 500x30 & -0.82 & -0.82 & -0.73 & +3.0 & 6.36 \\
\midrule
$U[100,120]$ & 50x10 & +0.00 & +0.00 & -0.01 & +0.2 & 0.71 \\
$U[100,120]$ & 100x10 & +0.00 & -0.02 & +0.04 & +0.1 & 0.39 \\
$U[100,120]$ & 100x20 & -0.01 & -0.02 & -0.01 & +0.3 & 0.52 \\
$U[100,120]$ & 200x20 & +0.02 & +0.01 & +0.02 & +0.1 & 0.41 \\
$U[100,120]$ & 500x30 & +0.05 & +0.04 & +0.11 & +0.1 & 1.98 \\
\midrule
$U[100,200]$ & 50x10 & -0.02 & -0.06 & -0.03 & +3.3 & 2.95 \\
$U[100,200]$ & 100x10 & -0.04 & -0.07 & -0.10 & +2.5 & 1.67 \\
$U[100,200]$ & 100x20 & +0.04 & -0.03 & +0.07 & +6.5 & 2.44 \\
$U[100,200]$ & 200x20 & -0.04 & -0.04 & +0.03 & +6.3 & 1.58 \\
$U[100,200]$ & 500x30 & +0.09 & +0.03 & +0.07 & +3.1 & 2.02 \\
\bottomrule
\end{tabular}}
\end{table}

% Auto-generated by analysis/make_tables_cycle3.py (Entrada 125)
\begin{table}[t]
\centering
\caption{$\mathrm{R}||C_{\max}$ with larger budgets on the 160 instances that have a MILP reference. $\bar d$ in \% of the makespan (negative: first better; for RB1 the entry is the mean change of the cell-level policy effect in percentage points). Wilcoxon, Holm over four. The realised CPU ratio ILS($T=8000$)/IG+LS($T=1000$) was $1.25$.}
\label{tab:f23_unrelated_budget}
{\scriptsize\setlength{\tabcolsep}{3.5pt}
\begin{tabular}{@{}lrrrr@{}}
\toprule
Comparison & $N$ & $\bar d$ (\%) & 90\% interval & $p_{\mathrm{Holm}}$ \\
\midrule
RB1: policy effect, $T=5000$ vs.\ $1000$ (20 cells) & 20 & +0.085 & -- & 0.312 \\
RB2: IG+LS ($T=1000$) vs.\ ILS ($T=8000$, equal CPU) & 160 & -2.158 & [-2.520, -1.783] & $4\times10^{-19}$ \\
RB3: IG+LS \textsc{QL} vs.\ random, $T=5000$ & 160 & -0.087 & [-0.212, +0.042] & 0.128 \\
RB4: IG+LS \textsc{QL} ($T=5000$) vs.\ MILP (60\,s) & 160 & +0.951 & [+0.743, +1.171] & $4\times10^{-13}$ \\
\bottomrule
\end{tabular}}
\end{table}


\paragraph{The policy gain follows the headroom of the strategies, and needs state and learning together}
The observation that the policy gains where the headroom is large was made on the
first sample and therefore tested on 400 \emph{new} instances (same cells and
parameters, nothing re-tuned; Table~\ref{tab:f24_unrelated_law}). It is
confirmed: the Spearman correlation between the headroom of a cell and the
policy gain in it is $0.59$ over the 20 cells ($p_{\mathrm{Holm}}=0.015$; $0.68$ in
the first sample), and Q-Learning captures $30\%$ of the headroom overall
(90\% interval $0$--$54\%$). The magnitude, however, is unstable: on the
new sample the pooled gain is $0.073\%$ (not significant on its own,
$p_{\mathrm{Holm}}=0.21$), half the first one, because $U[1,100]$, where the first
sample captured a large part of a $0.5$--$1.1\%$ headroom, shows none now; the two
independent samples together give $0.109\%$ (90\% interval $0.052$--$0.165$,
$n=800$; computed after seeing both). Decomposing the agent as in
Section~\ref{sec:igvariants}, neither ingredient is enough alone: Q-Learning without
state is no better than random selection ($+0.008\%$), a state-conditioned bandit is
not either, and the full agent beats the latter by $0.136\%$
($p_{\mathrm{Holm}}=0.009$); UCB1 without state is slightly worse than random
($+0.098\%$). The gain thus requires both the state and temporal-difference
learning, is proportional to the headroom of the action space, and, in these
experiments, is at most about a tenth of a percent.

% Auto-generated by analysis/make_tables_cycle3.py (Entrada 125)
\begin{table}[t]
\centering
\caption{$\mathrm{R}||C_{\max}$ on 400 \emph{new} instances (same cells and parameters): the policy gain follows the headroom (L1, Spearman $\rho$ over the 20 cells) and its decomposition into state and temporal-difference learning (L2--L6; IG with local search, $T=1000$). $\bar d$ in \% of the makespan (negative: first better); Wilcoxon, Holm over six.}
\label{tab:f24_unrelated_law}
{\scriptsize\setlength{\tabcolsep}{3.5pt}
\begin{tabular}{@{}lrrrr@{}}
\toprule
Test & W/L/T & $\bar d$ (\%) or $\rho$ & 90\% interval & $p_{\mathrm{Holm}}$ \\
\midrule
L1: Spearman(headroom, policy gain), 20 cells & -- & +0.59 & -- & 0.015 \\
L2: \textsc{QL} vs.\ random selection & 102/102/196 & -0.073 & [-0.146, -0.000] & 0.208 \\
L3: \textsc{QL} vs.\ \textsc{QL} without state & 110/108/182 & -0.075 & [-0.151, +0.007] & 0.208 \\
L4: \textsc{QL} vs.\ UCB1 by state & 124/106/170 & -0.136 & [-0.218, -0.054] & 0.009 \\
L5: \textsc{QL} without state vs.\ random & 88/101/211 & +0.008 & [-0.074, +0.086] & 0.631 \\
L6: UCB1 (no state) vs.\ random & 92/107/201 & +0.098 & [+0.023, +0.175] & 0.208 \\
\bottomrule
\end{tabular}}
\end{table}



\subsection{A Third Problem and a Synthesis across Environments}
\label{sec:pfsp}

To test whether the attribution is a property of parallel machine scheduling, the same protocol was applied to the permutation flowshop (makespan), the problem in which Iterated Greedy originated \citep{RuizStutzle2007}. The IG starts from NEH, removes $d$ random jobs, reinserts them in the best position (Taillard's acceleration), applies an insertion local search and accepts with a constant temperature; the action is the destruction size $d\in\{2,3,4,6,8\}$. The ILS has the same local search and perturbs by two random insertions. The instances are Taillard-style, generated here ($n\in\{50,100\}$, $m\in\{10,20\}$, $30$ test and $10$ tuning instances per size; $T=300$); the comparisons are paired, the CPU times are equalised by sequential calibration, and the design, the Holm family of five tests and the predictions were recorded before the test (Table~\ref{tab:f36_pfsp}).

\paragraph{The active ingredient depends on the problem}
In the flowshop the destruction--reconstruction cycle is not what separates the methods: the IG and the ILS, with the same local search, are equivalent at equal CPU time ($-0.03\%$, $p_{\mathrm{Holm}}=0.18$; the prediction that the IG would win was wrong). The local search, however, pays its cost: the IG without it, with $17$ to $57$ times more iterations for the same CPU time, is worse by $0.21\%$ ($65$ wins against $17$, $p_{\mathrm{Holm}}=10^{-6}$), the opposite of what was found for identical machines (Section~\ref{sec:f3f4-ablation}). With the results for unrelated machines, the active ingredient is thus the destruction--reconstruction cycle for identical machines, its combination with a descent for unrelated machines and the local search for the flowshop, and only an ablation at equal CPU time reveals which.

\paragraph{The policy over the destruction size is harmful again}
Q-Learning over $d$ is worse than random choice of $d$ ($+0.14\%$, $p_{\mathrm{Holm}}=5\times10^{-4}$), the validated size beats it by $0.16\%$ ($p_{\mathrm{Holm}}=8\times10^{-5}$), and the headroom of the size over random choice is null ($-0.02\%$): the signs of the predictions recorded beforehand were right, except that the validated size was expected only to match the agent, not to beat it.

\paragraph{Synthesis}
Table~\ref{tab:f35_synthesis} puts the seven environments side by side. The oracle headroom is negative (some fixed action beats random selection per instance) in six of the seven, and the validated choice captures a part of it wherever it is real. The Q-Learning agent behaves differently by the type of action. Over qualitatively different destruction strategies its gain is small but never positive, and it recovers $0.44$ to $0.68$ of the validated gain at $1000\times50$ and on the standard benchmark; over the destruction size, an ordered parameter, it is worse than random choice in all three problems ($+0.015\%$, $+0.152\%$, $+0.142\%$), although the size has headroom that validation captures in two of them. This regularity was noticed after the results and is a hypothesis, not a finding: headroom is necessary for a policy gain, and whether an online agent can use it seems to depend on the action type (and, in the diagnostic above, on how early the reward arrives).

% Auto-generated by analysis/make_tables_cycle5.py (Entrada 129)
\begin{table}[t]
\centering
\caption{A third problem: the permutation flowshop (Taillard-style instances generated here, $n\in\{50,100\}$, $m\in\{10,20\}$; 30 test instances per size, 10 tuning instances per size; Iterated Greedy of Ruiz and St\"utzle with destruction size $d\in\{2,3,4,6,8\}$ as the action; $T=300$). $\bar d$ in \% of the makespan (negative: first better); Wilcoxon, Holm over five. The ILS has the same insertion local search; the IG without local search receives $T_{\mathrm{eq}}$ iterations (5{,}177 to 17{,}076) so that the CPU time is the same, and the ILS $217$ to $250$.}
\label{tab:f36_pfsp}
{\scriptsize\setlength{\tabcolsep}{3.5pt}
\begin{tabular}{@{}lrrrr@{}}
\toprule
Comparison & W/L/T & $\bar d$ (\%) & 90\% interval & $p_{\mathrm{Holm}}$ \\
\midrule
P1: IG vs.\ ILS at equal CPU (mechanism) & 53/31/36 & -0.027 & [-0.070, +0.014] & 0.185 \\
P2: IG with vs.\ without local search at equal CPU & 65/17/38 & -0.212 & [-0.279, -0.148] & $1\times10^{-6}$ \\
P3: \textsc{QL} over $d$ vs.\ random $d$ (equal $T$) & 25/68/27 & +0.142 & [+0.073, +0.210] & $5\times10^{-4}$ \\
P4: validated $d$ vs.\ \textsc{QL} & 64/25/31 & -0.161 & [-0.225, -0.101] & $8\times10^{-5}$ \\
P5: validated $d$ vs.\ random $d$ & 47/38/35 & -0.020 & [-0.083, +0.043] & 0.466 \\
\bottomrule
\end{tabular}}
\end{table}


% Auto-generated by analysis/make_table_synthesis.py (Entrada 154)
\begin{table}[t]
\centering
\caption{Headroom and policy gain across the environments of the study (mean $\bar d$ in \% of the makespan relative to random selection with the same budget; negative: better). Oracle: best of the fixed actions per instance (an upper bound that includes selection noise); validated: the action chosen offline on separate instances; Q-Learning: the online agent (tabular, $T=1000$ except the flowshop, $T=300$). The last column is the CPU time of Q-Learning relative to random selection.}
\label{tab:f35_synthesis}
{\scriptsize\setlength{\tabcolsep}{3pt}
\begin{tabular}{@{}lrrrrrr@{}}
\toprule
Environment & $N$ & Actions & Oracle & Validated & Q-Learning & CPU ratio \\
\midrule
$\mathrm{P}||C_{\max}$, destruction strategy & 1{,}130 & 3 & -0.039 & +0.001 & -0.016 & -- \\
$\mathrm{P}||C_{\max}$, destruction size & 1{,}200 & 5 & -0.127 & -0.089 & +0.015 & -- \\
$\mathrm{R}||C_{\max}$ synthetic, destruction strategy & 400 & 4 & -0.240 & -- & -0.073 & 0.97 \\
$\mathrm{R}||C_{\max}$ $1000{\times}50$, destruction strategy & 96 & 4 & +0.092 & -0.461 & -0.312 & 1.42 \\
$\mathrm{R}||C_{\max}$ standard benchmark, destruction strategy & 700 & 4 & -0.149 & -0.124 & -0.055 & 1.12 \\
$\mathrm{R}||C_{\max}$ standard benchmark, destruction size & 700 & 5 & -0.325 & -0.201 & +0.152 & 0.65 \\
Permutation flowshop, destruction size & 120 & 5 & -0.193 & -0.020 & +0.142 & 0.89 \\
\bottomrule
\end{tabular}}
\end{table}


\subsection{The Companion \texttt{ils\_ql} Study}
\label{sec:ils-summary}

The \texttt{ig\_ql} design grew out of an earlier study of the same agent
selecting among four neighbourhood operators inside ILS, reported in full
in Supplementary Section~\ref{S-sec:ils-chapter}; its findings motivate the
decomposition of Section~\ref{sec:f3f4-ablation}. Every ILS variant reduces
GAP by $55$--$99\%$ relative to LPT, yet \texttt{ils\_ql} is statistically
indistinguishable from uniform random operator selection in all six
original classes ($p > 0.07$, $N=30$) and from a random-forest selector in
five of six. UCB1's apparent loss to both ($p \leq 0.008$) was an artefact
of raw rewards: normalised to $[0,1]$ it is indistinguishable from
\texttt{ils\_ql} ($p=0.78$, $N=100$) and from random selection; Probability
Matching, Adaptive Pursuit and SARSA show the same equivalence. The
apparent operator preference of the learned policy is largely an artefact
of coverage: only $26.5\%$ of the 81 nominal states are visited in a typical
run, and zero-initialised tables break ties toward the first operator.
Extended-budget random selection beats extended-budget Q-Learning for
$n \geq 30$ ($p \leq 0.028$): at this scale additional iterations
contribute more than adaptive sophistication. A preliminary exploration of
hybrid IG/ILS mechanisms is in Supplementary Section~\ref{S-sec:hybrid}.


% =====================================================================
%% Secoes Sintese + Discussao + Conclusoes fundidas na Entrada 95
%% (04/10/2026): as tres repetiam as mesmas alegacoes; versao anterior em
%% learnheuristic_pmsp_pre_entrada95_cut_backup_20261004.tex.
\section{Discussion}
\label{sec:discussion}

The title poses a question: when a learnheuristic wins, is it the search
mechanism or the learned policy? Table~\ref{tab:evidence} maps the claims to their
evidence; the experiments answer in layers and also mark the limits of what can be
claimed.

\paragraph{Where the gain comes from}
The search structure accounts for most of the advantage of \texttt{ig\_ql}. With
learning removed from both sides, \texttt{ig\_random} already beats
\texttt{ils\_random} in ten of twelve classes (Section~\ref{sec:f3f4-ablation}),
by $0.5\%$ on average in the seed-averaged data, also with a strong start in both
mechanisms and at equal CPU time; for unrelated machines the IG with local search
beats the ILS by $4.2\%$. The Q-Learning policy layered on top adds about
$0.01$--$0.05$ percentage points of GAP for identical machines, a difference that no
single class resolves but that a pooled, pre-specified test confirms in three of four
published families, that replicates on $2{,}690$ instances never run before, and that
is the same with an informative state; on seed averages of the twelve classes it is
real but strictly negligible ($-0.011\%$). The companion \texttt{ils\_ql} study
reached the same direction in a stronger form: learned operator selection is
indistinguishable from random selection, and for unrelated machines it is worse
($+0.96\%$), as it is for identical machines with an informative state. That
unrelated mechanisms, sharing the state and agent but not the action set, show
learning contributing little or nothing is stronger evidence about this agent at
this scale than any single result.

\paragraph{Why the learned policy adds so little}
Five candidate explanations were tested. (i) \emph{Nothing to learn.} The three
destruction strategies are nearly interchangeable: an oracle choosing the best of
the three for every instance would gain only $0.04\%$. (ii) \emph{State.} The
original state was nearly degenerate; an informative one did not help in IG, hurt
in ILS, and added nothing in a second confirmatory test. (iii) \emph{Signal.} The
reward registers on average between $0.7$ and $4.6$ improving steps per run, yet
a dense reward made the agent worse and pre-training left it at the level of random
choice. (iv) \emph{Budget.} Tightening the budget creates headroom ($0.10\%$ at
$T=100$) but not a policy gain; enlarging it five-fold does not help either.
(v) \emph{Learnability.} Even when the action space has real headroom (the
destruction size), Q-Learning captures $2\%$ of it at $T=1000$ and not significantly
more with $T=5000$ or a pre-trained table, while choosing the parameter by
validation on separate instances captures it. The common thread is that what a tabular agent can learn
online in a thousand iterations is small relative to what is available; the gain
lies in offline tuning across instances. For unrelated machines, where the
headroom is larger ($0.24$--$0.38\%$), the agent captures about $30\%$ of it and the
gain follows the headroom across cells (Spearman $0.59$ on new instances), but only the
combination of state and temporal-difference learning captures it and the size is
unstable across samples ($0.07$--$0.15\%$). At $1000\times50$ it captures all of a $1.7\%$
headroom in $U[1,100]$, where one strategy dominates, and is worse than random selection
in two of four distributions, at equal CPU time too, so a general gain is not
confirmed there (pooled $0.35\%$, $p_{\mathrm{Holm}}=0.19$). On the standard benchmark of 700 instances the headroom law holds for the destruction strategies (Spearman $0.74$ over 28 cells), the policy gain is $0.03$--$0.06\%$ (not confirmed) and validation beats the agent by $0.066\%$. It does not hold, however, when the destruction size is the action: with a headroom that validation captures ($0.20\%$ over random choice of $d$; $0.33\%$ for a per-instance oracle), Q-Learning is worse than random choice ($+0.15\%$), also when pre-trained offline, and the correlation between headroom and policy gain over cells is negative ($-0.37$). The headroom of the actions is therefore a necessary condition for a policy gain, not a sufficient one; an exploratory diagnostic suggests why: the immediate reward ranks the actions consistently with their final quality, but more than $90\%$ of it arrives in the first tenth of the run. This is a statement about this agent and
these action sets, not evidence that learning cannot help with richer action spaces,
other learners or tighter budgets.

\paragraph{What works in practice}
The strongest method of the study contains no learned component: Iterated Greedy
with the destruction size chosen on separate instances, started from the best of
LPT, MULTIFIT and LDM. It is never worse than any of them, improves the best
constructive solution in one instance in four (synthetic) to two in five
(published), and is provably optimal on $74.7\%$ of the synthetic instances
against $48.5\%$ for MULTIFIT. Its budget matters more than any added component: a
local search, which makes an iteration $2$--$80$ times more expensive, improves the recipe by
$0.036\%$ at an equal number of iterations, but at equal CPU time it is worse by
$0.034\%$ than the recipe without it, which then reaches $89.0\%$ of provably
optimal solutions; in the plain IG the same descent is worth $0.150\%$ at an equal
number of iterations and $0.010\%$ at equal CPU time. MULTIFIT's own strength is distribution-dependent:
it is worse than LPT on narrow-range times. The exact method remains the right tool
for the instances it solves in a second ($85\%$ of families C and F;
\citealp{OstojicEtAl2025}), and for unrelated machines at the sizes studied, up to
$1000\times50$, the MILP beats the IG with local search and is matched only by the IG started from the linear relaxation.

\paragraph{A portable design for studies of adaptive selection}
Six elements of the protocol are inexpensive to add to other studies and changed
the conclusions here. A random-selection control at matched budget separates the
mechanism from the policy; a family-level pooled test, and averaging over seeds,
detect small effects that per-class single-run tests cannot; hypotheses and
decision rules frozen before the runs separate a confirmed effect from one
generated after looking at the data (one effect that looked large in the
exploration shrank by more than half on held-out instances, and a headroom law
suggested by one sample was tested on a fresh one); equivalence margins state
when an effect is negligible rather than merely significant; a headroom measure
(best fixed action, omniscient choice) tells what any selector could gain; and a
certified lower bound shows that a sizeable part of the conventional gap to a
classical bound is bound slack.

\paragraph{Problem structure decides the mechanism}
Which mechanism is stronger depends on the instance, not on a single size ratio:
the granularity of the processing times decides the direction between IG and ILS
while $n/m$ shifts the crossover. The finding does not depend on the budget
convention (an ILS iteration is costlier, but the ILS advantage on fine-grained
times persists at equal CPU time), nor on the distribution (it holds on six of
them), nor on the start (it holds with the best constructive solution in both). A
one-feature rule, tested on held-out instances and again on new distributions,
beats both fixed policies, but its pooled gain over always using IG is only $0.013$
percentage points on the published families, with a cost of $0.083$ points in
family C.

\paragraph{Limits}
These results do not make \texttt{ig\_ql} a competitive solver for the pure PMSP, and
the paper does not claim it. (i) Exact methods are strong on the published
benchmark; where the iterative arc-flow method needs more than a second,
\texttt{ig\_ql} is $1.6$--$2.1\%$ above the optimum. (ii) The default destruction size
was not tuned for the main comparisons, to avoid tuning on the evaluation
instances; the tuned recipe of Section~\ref{sec:recipe} is reported separately,
with validation on separate instances. (iii) The agent is tabular, one dense reward and
one pre-training scheme were tried, and the action spaces are small; function
approximation and richer action spaces are outside the evidence. (iv) The unrelated-machine
heuristics are not state of the art (with the linear-relaxation start and a restricted MILP, $0.25\%$ above the best known solutions of the standard benchmark, against $0.1\%$ for the exact references; $0.67\%$ for the job-correlated instances; without the MILP, $0.97\%$),
most of the instances of that section are synthetic (the standard benchmark was used for one validation), and the policy gain there is unstable across samples. (v) The confirmed
learning effect on published instances was measured with the original state
(the second test added the informative one); family F had no unseen instances.
(vi) Each instance is run once in the main batteries, with seed averaging on twelve
classes. (vi\,b) The main comparisons use $T=1{,}000$ iterations; the equal-CPU analyses
of the recipe and of the hybrid show that conclusions about added components depend
on the budget convention, and the policy was not studied at budgets much beyond
$T=5{,}000$ or, for identical machines, at sizes where the search is far from saturated
(for unrelated machines at $1000\times50$ it was, with a mixed result). (vii) The pre-specification is a dated, hash-stamped diary, not a public
registry, and the number of analyses is large: all families and verdicts, including
every non-confirmation, are in Supplementary Table~\ref{S-tab:s_prereg_all}.

\section{Conclusions}
\label{sec:conclusions}

This paper asked where the gain of Q-Learning-guided metaheuristics for parallel
machine scheduling comes from, and the answer is the search structure, not the
learned policy. \texttt{ig\_ql} is never significantly beaten by LPT or by seven
ILS variants across twelve synthetic classes, beats a simulated-annealing method
and a Bee Colony Optimisation algorithm, also when their budget rules are relaxed,
but loses to MULTIFIT and to LDM at large $n$. Random-selection controls attribute
most of the gain to the search structure in both problems studied (identical and
unrelated machines); the learned selection adds about $0.01$--$0.05$ percentage
points for identical machines, a small effect that replicates on held-out and fresh
published instances and does not depend on the state, and about $0.1\%$ for
unrelated machines, where it follows the headroom of the destruction strategies but not that of the destruction size. It is limited by
that headroom and by what a tabular agent can learn online: for the destruction
size, which has headroom, validation on separate instances beats the agent even with
five times the budget or a pre-trained table. The resulting learning-free recipe,
a validated destruction size and a start from the best of three constructive
heuristics, is the strongest method of the study, provably optimal on $75\%$ of the
synthetic instances against $49\%$ for MULTIFIT, and on $89\%$ with the CPU time of a
local-search hybrid, which gains over it only at an equal number of iterations. Which of IG and ILS is stronger
depends on the granularity of the processing times, at equal CPU time, on six
distributions and with strong starts, and a one-feature dispatch rule exploits it
for a small pooled gain. A companion study shows that Q-Learning-based operator
selection inside ILS is indistinguishable from random selection, with an apparently
interpretable policy that is largely an artefact of sparse state coverage.

The limits are as informative as the gains: \texttt{ig\_ql} is not competitive with
exact methods on the instances they solve quickly, and for unrelated machines the
MILP beats the IG with local search up to $1000\times50$ and is matched only by the IG started from the linear relaxation. Future work should
(i) test richer action spaces, other learners, rewards that arrive early enough and tighter budgets, where a policy
might have room, and refine the headroom law, which proved necessary but not sufficient, with pre-specified action spaces of
controlled quality; (ii) pair the decomposition with the stronger unrelated-machine heuristic of this study (LP start and restricted MILP, $0.25\%$ above the best known solutions of the standard benchmark) or with one of the literature; (iii) combine the validated destruction size with learned
selection of other decisions; and (iv) apply the decomposition protocol to problem
variants where exact methods do not reach, such as heterogeneous machines with
setup times.



%% =====================================================================
%% Secoes administrativas obrigatorias do EJOR, adicionadas na Entrada 85
%% (09/09/2026, Task #45) -- pendencia sinalizada desde a Fase A (Entrada
%% 74) e deixada propositalmente em aberto ate' agora, por decisao de
%% Carlos. Modelo de estrutura seguindo o artigo irmao do LCPP (mesmo
%% grupo/coautoria); texto de cada secao especifico deste artigo, nao
%% copiado do LCPP, exceto o boilerplate legal generico (Declaration of
%% Competing Interest), que e' texto padrao da Elsevier.
%% TODO(Carlos): todo campo entre colchetes [Inserir ...] abaixo e' um
%% placeholder que so' voce pode preencher (repositorio de dados/codigo,
%% agencia/numero de bolsa). A "Declaration of Generative AI" precisa da
%% sua leitura e ajuste antes da submissao -- e' uma declaracao factual
%% sobre o seu proprio processo de pesquisa, nao uma formalidade que eu
%% deva decidir sozinho por voce. O CRediT abaixo lista so' os papeis de
%% autor unico (sem Supervision/Funding acquisition/Project
%% administration); ajuste ou adicione linhas se/quando entrar coautoria.
%% =====================================================================
\begingroup
\small
\begin{spacing}{1.0}

\section*{CRediT authorship contribution statement}
\textbf{Carlos Vinicio Dantas Ara\'{u}jo} (Conceptualization, Methodology,
Software, Validation, Formal analysis, Investigation, Data curation,
Writing original draft, Writing review and editing, Visualization).

\section*{Declaration of Competing Interest}
The author declares no known competing financial interests or personal
relationships that could have appeared to influence the work reported
in this paper.

\section*{Declaration of Generative AI and AI-Assisted Technologies in the Writing Process}
During the preparation of this work, the author used a generative
AI-based assistant (Claude, Anthropic) to support software implementation
and testing, the writing of data-analysis scripts, and the drafting,
translation and editing of the manuscript text, based on the author's own
research design, data, and analysis. After using this tool, the author reviewed
and edited the content as needed and takes full responsibility for the
content of the published article.

\section*{Data and Code Availability}
The source code, the script of every experiment (with the md5 hashes of the
versions frozen before each pre-specified analysis), all result files, the
LaTeX sources and the research diary are available at [Insert Repository
Link/DOI, e.g., GitHub, Zenodo], under the MIT licence (code) and CC BY 4.0
(results). The external benchmark instances are publicly available from the
sources cited in Section~\ref{sec:external} and, for unrelated machines, in Section~\ref{sec:unrelated} (the copy distributed with \citealp{AbladEtAl2021}), and are not redistributed; the
synthetic instances are regenerated from their seeds.

\section*{Funding}
This research received no specific grant from any funding agency in the
public, commercial, or not-for-profit sectors.

\section*{Acknowledgements}
The author thanks the Universidade Estadual do Cear\'{a} (UECE) for
institutional support during the development of this work.

\end{spacing}
\endgroup

\begingroup
\scriptsize
\setlength{\bibsep}{0pt}
\begin{spacing}{0.87}
\bibliographystyle{elsarticle-harv}
\bibliography{refs}
\end{spacing}
\endgroup

\end{document}
